%%%%%%%%%%%%%%%%%%%%%%%%%%%%%%%%%

% \begin{instructbox}[\incode{{bench_pred_name}}]
% MIPROv2 Prompt:
% \begin{lstlisting}[breakindent=0pt, breaklines=true, breakatwhitespace=true, frame=single, backgroundcolor=\color{green!10}]
% {mipro_prompt}
% \end{lstlisting}

% GEPA Prompt:
% \begin{lstlisting}[breakindent=0pt, breaklines=true, breakatwhitespace=true, frame=single, backgroundcolor=\color{blue!10}]
% {gepa_prompt}
% \end{lstlisting}
% \end{instructbox}

%%%%%%%%%%%%%%%%%%%%%%%%%%%%%%%%%%

\subsection{HotpotQA, GPT-4.1 Mini}
\begin{instructbox}[\incode{HotpotQA GPT-4.1 Mini create_query_hop2.predict}]
Base Prompt:
\begin{lstlisting}[breakindent=0pt, breaklines=true, breakatwhitespace=true, frame=single, backgroundcolor=\color{gray!10}]
Given the fields `question`, `summary_1`, produce the fields `query`.
\end{lstlisting}

MIPROv2 Prompt:
\begin{lstlisting}[breakindent=0pt, breaklines=true, breakatwhitespace=true, frame=single, backgroundcolor=\color{green!10}]
You are an expert multi-hop question answering system designed to refine retrieval queries for complex questions. Given the original question and an initial summary of retrieved documents (summary_1), think step by step to analyze how the summary relates to the question. Use this reasoning to generate a focused and precise query that will guide the retrieval of additional relevant information needed to answer the question completely. Your output should include a clear chain-of-thought reasoning process followed by the refined query that targets complementary evidence for the next retrieval hop.
Demos:
Example 1:
question:
    Simone Benedetti plays for a team that currently competes in which Serie?
summary_1:
    Simone Benedetti plays for Virtus Entella, but the passages do not specify which Serie the team currently competes in.
reasoning:
    The summary indicates that Simone Benedetti plays for Virtus Entella, but it does not state the current Serie (league level) in which Virtus Entella competes. To accurately answer which Serie the team currently competes in, it is necessary to look up the most recent league information for Virtus Entella.
query:
    What Serie (league level) does Virtus Entella currently compete in?


##################################
3 other demos omitted for brevity.
##################################

\end{lstlisting}

GEPA Prompt generated by config GEPA:
\begin{lstlisting}[breakindent=0pt, breaklines=true, breakatwhitespace=true, frame=single, backgroundcolor=\color{blue!10}]
You will be given two input fields: `question` and `summary_1`.

Your task is to generate a new search query (`query`) optimized for the **second hop** of a multi-hop retrieval system. The original user question is typically complex and requires information from multiple documents to answer. The first hop query is the original question used to retrieve an initial set of documents. Your goal is to generate a **second hop query** that retrieves *additional relevant documents* that were *not* found in the first hop but are necessary to answer the original question completely.

Detailed task instructions and hints:

1. **Input Understanding:**
   - `question` is the original multi-hop question posed by the user.
   - `summary_1` is a concise summary of information from a document retrieved in the first hop, which partially addresses the question.

2. **Purpose and Context:**
   - Your generated `query` aims to find the *missing pieces* of information needed to fully answer the `question`.
   - The multi-hop retrieval system works in stages:
     - First hop: The original question returns some documents.
     - Second hop: Your query must help retrieve any *other relevant documents* NOT found in the first hop that hold complementary or broader context necessary for final answer extraction.

3. **Key Observations from Examples and Feedback:**
   - First-hop documents often cover one entity or aspect in the question.
   - Remaining relevant documents often involve connected or higher-level concepts mentioned in `summary_1` but not explicitly asked in the original question.
   - The `query` should be formulated to explicitly target these *missing*, but logically linked, documents.
   - Avoid merely paraphrasing the original question or restating known facts from `summary_1`.
   - Instead, infer what broader or related entities/concepts might provide the crucial missing information.
   - For example, if `summary_1` describes a population for a small civil parish, but the question wants total population of the wider region, your `query` should target that wider region (e.g., "Madeira archipelago population in 2011").
   - Similarly, if `summary_1` covers a song and the question wants the album it came from, but first hop got song-level documents, your query should retrieve documents about the album itself.

4. **How to Build the Query:**
   - Identify the entities or topics mentioned in `summary_1` that appear related but different from first-hop documents.
   - Reframe the query to explicitly mention these broader or related entities connected to the original question.
   - Include relevant key context from the question to maintain specificity, but shift focus to the missing piece.
   - The goal is to retrieve documents that link or complement what was retrieved initially.

5. **Practical Strategy:**
   - Read the `summary_1` carefully to spot references to bigger contexts or other entities not covered in the first hop.
   - Ask yourself, "What entity or aspect does this summary hint at that could answer the original question but was not found yet?"
   - Formulate a precise, focused factual query targeting that entity or concept to retrieve the missing documents.

6. **Output:**
   - Produce only the field `query` as a clear, concise question or keyword phrase designed for efficient retrieval of **second-hop documents**.
   - Ensure the query relates logically to the original question while targeting the broader or complementary knowledge identified in `summary_1`.
   - Do **not** include the original question or simply rephrase it.
   - Do **not** duplicate information already well-covered by the first hop retrieval.

By following these principles, you will help the multi-hop retrieval system find all necessary documents to answer the multi-faceted original question completely.
\end{lstlisting}
\end{instructbox}\begin{instructbox}[\incode{HotpotQA GPT-4.1 Mini final_answer.predict}]
Base Prompt:
\begin{lstlisting}[breakindent=0pt, breaklines=true, breakatwhitespace=true, frame=single, backgroundcolor=\color{gray!10}]
Given the fields `question`, `summary_1`, `summary_2`, produce the fields `answer`.
\end{lstlisting}

MIPROv2 Prompt:
\begin{lstlisting}[breakindent=0pt, breaklines=true, breakatwhitespace=true, frame=single, backgroundcolor=\color{green!10}]
You are an expert multi-hop reasoning assistant skilled in synthesizing information from multiple summaries to answer complex questions. Given the `question`, along with two intermediate summaries `summary_1` and `summary_2` that contain relevant evidence, carefully analyze and integrate the information step-by-step to produce a clear, logical reasoning process followed by a concise and accurate final answer.
Demos:
Example 1:
question:
    are Machaeranthera and Prumnopitys both the plants?
summary_1:
    Yes, both Machaeranthera and Prumnopitys are plants. Machaeranthera is a genus of flowering plants in the daisy family, while Prumnopitys is a genus of coniferous evergreen trees in the podocarp family.
summary_2:
    Yes, both Machaeranthera and Prumnopitys are plants; Machaeranthera is a genus of flowering plants in the daisy family, while Prumnopitys is a genus of coniferous evergreen trees in the podocarp family.
reasoning:
    Both summaries agree that Machaeranthera and Prumnopitys are plants. Specifically, Machaeranthera is a genus of flowering plants in the daisy family, and Prumnopitys is a genus of coniferous evergreen trees in the podocarp family. Since both belong to plant genera, the answer is yes.
answer:
    yes


##################################
3 other demos omitted for brevity.
##################################

\end{lstlisting}

GEPA Prompt generated by config GEPA:
\begin{lstlisting}[breakindent=0pt, breaklines=true, breakatwhitespace=true, frame=single, backgroundcolor=\color{blue!10}]
Task Description:

You are given three fields as input: `question`, `summary_1`, and `summary_2`. Your goal is to produce an `answer` field that directly and explicitly responds to the question using the information from the two summaries, enhanced by your authoritative domain knowledge when needed.

Input Format:

- `question`: A natural language question that may require a fact, definition, name, date, yes/no response, or other specific information.
- `summary_1` and `summary_2`: Two independently generated summaries or snippets containing information related to the question. They may vary in completeness, accuracy, and specificity.

Requirements and Approach:

1. **Understand the question precisely.** Determine exactly what is asked—whether a name, a specific fact, a date, or a yes/no answer.

2. **Compare both summaries.** Analyze the content of `summary_1` and `summary_2`:
   - If they agree and directly answer the question, use this as primary evidence.
   - If one summary provides a fact that the other does not mention, carefully evaluate its plausibility.
   - If the summaries conflict, use domain expertise and authoritative knowledge to resolve or explicitly state uncertainty.

3. **Domain-specific factual verification and nuance:**
   - **Names and nicknames:** Provide only the specific nickname or name when asked, without extra phrasing. For example, when asked for the nickname of a person or entity, respond with the nickname alone, not a full sentence.
   - **Nationality and identity distinctions:** Use the most precise terms aligned with factual correctness and common usage (e.g., “English” vs. “British”) based on domain knowledge.
   - **Dates and historical facts:** Verify dates or historical claims with domain knowledge to pick the correct fact, especially when there might be confusion between franchise start dates vs. event dates etc.
   - **Yes/no questions:** Prefer concise answers of “yes” or “no” only, unless the question demands elaboration.
   - **Types or categories:** If a question asks about the type or category (e.g., type of company), provide the most direct concise phrase without including adjectives like nationality unless asked explicitly.

4. **Answer conciseness and relevance:**
   - Provide a brief and direct answer to the question.
   - Avoid repeating or restating the question.
   - Avoid unnecessary context unless requested or needed for clarity.
   - Avoid constructing full sentences unless needed; for example, answers to nickname or yes/no questions should be as short and specific as possible.

5. **When authoritative knowledge supplements the summaries:**
   - If the summaries are incomplete or potentially inaccurate, incorporate trusted knowledge from your training about the topic to provide the correct and precise answer.
   - For example, when a summary gives a year that conflicts with known release dates or factual details, prefer the verified date.
   - When the summaries differ in style (one uses a formal phrase, another provides just the nickname), respond with the correct, clean answer format (e.g., just the nickname alone).

6. **Examples of correct reasoning and answers:**
   - Question: “What is the nickname of the 2005 Toyota Grand Prix of Long Beach Polesitter ?”
     - Correct answer: `the thrill from West Hill`
   - Question: “What type of company is Zipcar led by Scott Griffith from 2003-2013?”
     - Correct answer: `car-sharing company`
   - Question: “Who was the partner of British comic book artist, Henry Flint, that helped create Zombo?”
     - Correct answer: `Al Ewing`

Summary:

- Use both summaries as primary but not sole evidence.
- Reliably verify and contextualize facts using domain knowledge, especially for nationality, dates, nicknames, company types, and yes/no questions.
- Provide short, direct answers matching the specificity requested.
- Avoid unnecessary elaboration unless explicitly required.
- Explicitly resolve conflicts or ambiguity using your knowledge or state uncertainty when appropriate.

This approach ensures that answers are both accurate and concise, suitable for direct consumption or integration in knowledge bases or question-answering systems.
\end{lstlisting}
\end{instructbox}\begin{instructbox}[\incode{HotpotQA GPT-4.1 Mini summarize1.predict}]
Base Prompt:
\begin{lstlisting}[breakindent=0pt, breaklines=true, breakatwhitespace=true, frame=single, backgroundcolor=\color{gray!10}]
Given the fields `question`, `passages`, produce the fields `summary`.
\end{lstlisting}

MIPROv2 Prompt:
\begin{lstlisting}[breakindent=0pt, breaklines=true, breakatwhitespace=true, frame=single, backgroundcolor=\color{green!10}]
Given a question and a set of related passages, carefully analyze the information by thinking through the relevant facts step-by-step. Produce a clear and concise summary that synthesizes the key points from the passages directly relevant to answering the question, ensuring the summary is focused, accurate, and grounded in the evidence provided.
Demos:
Example 1:
question:
    The architectural style of a church that stands in front of the Palazzo Ghisilardi Fava originated in what city?
passages:
    ['Palazzo Ghisilardi Fava | Palazzo Ghisilardi Fava is a Renaissance style palace, located on via Manzoni 4 in Bologna, region of Emilia Romagna, Italy; it houses the Medieval Civic Museum of Bologna.', 'Madonna di Galliera, Bologna | The Madonna di Galliera is a church with a Renaissance facade and Baroque interiors, located on Via Manzoni, in central Bologna, Italy.  It stands in front of the Palazzo Ghisilardi Fava.  The present name over the portal is the "Chiesa di Filippini Madonna di Galliera e Filippo Neri".', 'Palazzo Pretorio, Prato | The Palazzo Pretorio of the Tuscan town of Prato was the old city hall located town center, standing in front of the current Palazzo Comunale.  It now accommodates the Civic Museum of Prato, which was reopened on September 2013.', 'Palazzo Pisani Gritti | The Palazzo Pisani Gritti is a Gothic-style palace located across the Grand Canal from the Church of the Salute, and stands along the Campo del Traghetto where the church of Santa Maria del Giglio, in the Sestieri of San Marco, Venice, Italy.', 'Dömötör Tower | The Dömötör tower is the oldest building in Szeged, Hungary.  The foundation was most probably laid during the 11th century, while the lower part was built (in Romanesque style) from the 12th century, and the upper part (Gothic style) from the 13th century.  The tower was once part of the former "Szent Dömötör templom" (St. Demetrius church), but today it stands in Dóm Square, in front of the much larger Votive Church of Szeged.  The upper part was rebuilt from the original stones in 1926.  The architecture of the tower is similar to another found in Southern France, or in the territory of the former Byzantine Empire.', 'Palazzo Leoni, Bologna | The Palazzo Leoni is a Renaissance style palace located on Via Marsala #31, in front of the outlet of Via Mentana, in central Bologna, region of Emilia-Romagna, Italy.  The apse of the church of San Martino, is across the street.', 'Palazzo Altieri | Palazzo Altieri is a palace in Rome, which was the home of the Altieri family in the city.  The palace faces the square in front of the Church of the Gesù.']
reasoning:
    The question asks for the city where the architectural style of a church that stands in front of the Palazzo Ghisilardi Fava originated. Passage [2] identifies the church standing in front of the Palazzo Ghisilardi Fava as "Madonna di Galliera" and mentions it has a Renaissance facade and Baroque interiors. Passage [1] states the Palazzo Ghisilardi Fava is located in Bologna, Italy. Since the church is located in Bologna and features a Renaissance architectural style, it can be inferred that the architectural style of this church originated in the city of Florence, which is historically recognized as the birthplace of Renaissance architecture. Thus, the architectural style (Renaissance) of the church in front of the Palazzo Ghisilardi Fava originated in Florence.
summary:
    The Renaissance architectural style of the church standing in front of the Palazzo Ghisilardi Fava originated in Florence.


##################################
3 other demos omitted for brevity.
##################################

\end{lstlisting}

GEPA Prompt generated by config GEPA:
\begin{lstlisting}[breakindent=0pt, breaklines=true, breakatwhitespace=true, frame=single, backgroundcolor=\color{blue!10}]
You are a first-hop **summarization module** in a multi-hop question answering (QA) system. Your task is to generate a concise, informative `summary` given two input fields: a `question` and a list of relevant `passages`.

Your goal is to extract and synthesize key information from the retrieved passages that:

1. Directly relates to the initial question.
2. Captures the core facts and entities needed to understand the scope and context of the question.
3. Includes relevant connections, bridging entities, dates, locations, or descriptions that enable the system to devise focused and effective follow-up queries in subsequent hops.
4. Provides a strong factual foundation for downstream answer generation modules.

**Task specifics and best practices:**

- The `summary` must represent a distilled synthesis, not just a compression or extractive snippet.
- Explicitly include cited passage titles or key entity labels (e.g., "Children in Need 2006 | ..." or "Anthony Levandowski | ...") in your summary to highlight the origin of information.
- Incorporate sufficient context to hint at missing or un-retrieved supporting facts, thus enhancing the multi-hop retrieval process.
- When the question asks for an attribute (e.g., nationality, location, company origin), ensure you provide:
   - Identification of the relevant subject or entity mentioned in the passages.
   - The extracted attribute or relevant information as stated or implied.
   - Bridging details that could help the system pursue remaining information in the next retrieval step.
- Avoid forming a final answer; instead, focus on "what is known now" from the input documents to facilitate further query refinement.

**Examples of critical elements to include:**

- Entity names, roles, titles, and dates tied to the question.
- Names of organizations or locations connected through intermediary entities.
- Distinctive identifiers or clarifications that can help narrow down next-step retrieval (such as "Natasha Kaplinsky is an English presenter," "Kapolei is a city on Oahu," or "Waymo spun out of Alphabet").

**Format of output:**

Provide a paragraph or a few sentences that cohesively summarize the key passages in relation to the question, referencing passage titles or entities to frame facts clearly.

---

This approach ensures the summary is both informative for next-hop retrieval and foundational for final answer extraction in multi-hop QA.
\end{lstlisting}
\end{instructbox}\begin{instructbox}[\incode{HotpotQA GPT-4.1 Mini summarize2.predict}]
Base Prompt:
\begin{lstlisting}[breakindent=0pt, breaklines=true, breakatwhitespace=true, frame=single, backgroundcolor=\color{gray!10}]
Given the fields `question`, `context`, `passages`, produce the fields `summary`.
\end{lstlisting}

MIPROv2 Prompt:
\begin{lstlisting}[breakindent=0pt, breaklines=true, breakatwhitespace=true, frame=single, backgroundcolor=\color{green!10}]
Given a `question`, relevant `context`, and a list of `passages`, provide a clear, concise summary that integrates the key information from the passages in relation to the question and context. Use step-by-step reasoning to explain how the summary is derived from the evidence before presenting the final synthesized summary.
Demos:
Example 1:
question:
    What bank was founded by the great-great-great grandfather of the second Duke of Florence?
context:
    The bank founded by the great-great-great grandfather of the second Duke of Florence is the Medici Bank, established by Giovanni di Bicci de' Medici.
passages:
    ['Giovanni di Bicci de\' Medici | Giovanni di Bicci de\' Medici (c. 1360 – February 20/28, 1429) was an Italian banker, a member of Medici family of Florence, and the founder of the Medici Bank.  While other family members, such as Chiarissimo di Giambuono de\' Medici, who served in the Signoria in 1201, and Salvestro de\' Medici, who was implicated in the Ciompi Revolt of 1378, are historically significant, Giovanni\'s founding of the family bank truly began the family\'s rise to power in Florence.  He was the father of Cosimo de\' Medici ("Pater Patriae"), grandfather of Piero di Cosimo de\' Medici, great-grandfather of Lorenzo de\' Medici (the Magnificent) and great-great-great-grandfather of Cosimo I de\' Medici, Grand Duke of Tuscany.', 'Averardo de\' Medici | Averardo de\' Medici (1320-1363), 
    also known as Everard De Medici, was the son of Salvestro de\' Medici (died 1346), "il Chiarissimo" (English meaning "the very clear.")  and the father of three children: Giovanni, Francesco, and Antonia.  Giovanni di Bicci de\' Medici would later become the first historically relevant member of Medici family of Florence and the eventual founder of the Medici bank.', "Villa Medici at Cafaggiolo | The Villa Medicea di Cafaggiolo is a villa situated near the Tuscan town of Barberino di Mugello in the valley of the River Sieve, some 25\xa0kilometres north of Florence, central Italy.  It was one of the oldest and most favoured of the Medici family estates, having been in the possession of the family since the 14th century, when it was owned by Averardo de' Medici.  Averardo's son, Giovanni di Bicci de' Medici, is considered to be the founder of the Medici dynasty.", 'Medici: Masters of Florence | Medici: Masters of Florence is an Italian-British television drama series about the Medici dynasty set in 15th-century Florence, starring Dustin Hoffman as Giovanni di Bicci de\' Medici, Richard Madden as Cosimo de\' Medici, and Stuart Martin as Lorenzo de\' Medici ("The Elder").  The series was co-created by Frank Spotnitz ("The X-Files" and "Man in the High Castle") and Nicholas Meyer ("").  Sergio Mimica-Gezzan ("The Pillars of the Earth") directed all eight episodes.  Episodes 1 and 2 aired on Rai 1 (Italian TV) on 18 October 2016.  According to Italian ratings compiler Auditel, it attracted a record 7.6\xa0million viewers.  The first season consists of eight episodes.', 'Neri di Bicci | Neri di Bicci (1419–1491) was an Italian painter of the Renaissance.  A prolific painter of mainly religious themes, he was active mainly in Florence and in the medium of tempera.  His father was Bicci di Lorenzo.  His grandfather, Lorenzo di Bicci was also a painter in Florence, a pupil of Spinello Aretino.', 'Cosimo de\' Medici | Cosimo di Giovanni de\' Medici (called \'the Elder\' (Italian "il Vecchio") and posthumously "Father of the Nation" (Latin "pater patriae"); 27 September 1389 – 1 August 1464) was an Italian banker and politician, the first of the Medici political dynasty, "de facto" rulers of Florence during much of the Italian Renaissance.  Despite his influence, his power was not absolute; Florence\'s legislative councils at times resisted his proposals, something which would not have been tolerated by the Visconti of Milan, for instance; throughout his life he was always "primus inter pares" ("first among equals").  His power derived from his wealth as a banker, and he was a great patron of learning, the arts and architecture.', "Giovanni il Popolano | Giovanni de' Medici, in full Giovanni di Pierfrancesco de' Medici, later known as il Popolano (21 October 1467 – 14 September 1498) was an Italian nobleman of the Medici House of Florence.  He was the son of Pierfrancesco di Lorenzo de' Medici, and therefore a member of a secondary branch of the family."]
reasoning:
    The question asks for the bank founded by the great-great-great-grandfather of the second Duke of Florence. Passage [1] states that Giovanni di Bicci de' Medici founded the Medici Bank and was the great-great-great-grandfather of Cosimo I de' Medici, who was the Grand Duke of Tuscany. The Dukedom of Florence eventually transformed into the Grand Duchy of Tuscany under Cosimo I, who is identified as the second Duke of Florence in the question context.

Thus, Giovanni di Bicci de' Medici, as the great-great-great-grandfather of the second Duke of Florence (Cosimo I de' Medici), founded the Medici Bank.
summary:
    The Medici Bank was founded by Giovanni di Bicci de' Medici, the great-great-great-grandfather of the second Duke of Florence.


##################################
3 other demos omitted for brevity.
##################################

\end{lstlisting}

GEPA Prompt generated by config GEPA:
\begin{lstlisting}[breakindent=0pt, breaklines=true, breakatwhitespace=true, frame=single, backgroundcolor=\color{blue!10}]
Given the fields `question`, `context`, `passages`, produce the fields `summary`.
\end{lstlisting}
\end{instructbox}
\subsection{HotpotQA, Qwen3 8B}
\begin{instructbox}[\incode{HotpotQA Qwen3 8B create_query_hop2.predict}]
Base Prompt:
\begin{lstlisting}[breakindent=0pt, breaklines=true, breakatwhitespace=true, frame=single, backgroundcolor=\color{gray!10}]
Given the fields `question`, `summary_1`, produce the fields `query`.
\end{lstlisting}

MIPROv2 Prompt:
\begin{lstlisting}[breakindent=0pt, breaklines=true, breakatwhitespace=true, frame=single, backgroundcolor=\color{green!10}]
Given the question and the first summary, synthesize the key elements of the question and the summary. Identify the specific information that needs to be retrieved or confirmed. Formulate a focused and precise query that will guide the next step in the multi-hop reasoning process, ensuring it directly addresses the gap in knowledge or requires further clarification from additional context.
Demos:

\end{lstlisting}

GEPA Prompt generated by config GEPA+Merge:
\begin{lstlisting}[breakindent=0pt, breaklines=true, breakatwhitespace=true, frame=single, backgroundcolor=\color{blue!10}]
Given the fields `question` and `summary_1`, produce the field `query` that optimizes the retrieval of additional documents for a multi-hop system.  

**Task Details:**  
1. **Objective:** Your query must target documents not retrieved in the first hop, using clues from the summary and the original question.  
2. **Key Strategy:**  
   - Identify gaps in the first hop's retrieved documents (e.g., missing entities, relationships, or specific details).  
   - Use explicit information from the summary (e.g., names, locations, quantities) to rephrase the question into a query that surfaces new relevant documents.  
   - Avoid restating the answer directly; instead, structure the query to explore connections or unresolved details.  
3. **Domain-Specific Guidance:**  
   - If the summary explicitly answers the question, the query should still focus on retrieving documents that provide deeper context or verify the answer (e.g., "What is the headquarters location of [Company]?" instead of "The answer is [Location]").  
   - Leverage entities mentioned in the summary (e.g., "Carhartt," "Aubrey O'Day") to anchor the query.  
   - If no documents are missing, rephrase the query to explicitly request the answer (e.g., "Which has more acts, Elektra or From the House of the Dead?").  
4. **Avoid:**  
   - Generating queries that duplicate the original question.  
   - Assuming the summary contains all necessary information for the second hop.
\end{lstlisting}
\end{instructbox}\begin{instructbox}[\incode{HotpotQA Qwen3 8B final_answer.predict}]
Base Prompt:
\begin{lstlisting}[breakindent=0pt, breaklines=true, breakatwhitespace=true, frame=single, backgroundcolor=\color{gray!10}]
Given the fields `question`, `summary_1`, `summary_2`, produce the fields `answer`.
\end{lstlisting}

MIPROv2 Prompt:
\begin{lstlisting}[breakindent=0pt, breaklines=true, breakatwhitespace=true, frame=single, backgroundcolor=\color{green!10}]
Given the question, summary_1, and summary_2, generate a step-by-step reasoning process that logically connects the information from both summaries to arrive at the final answer, and then provide the answer. Ensure the reasoning explicitly references how the summaries support the conclusion and maintains consistency with the question's requirements.
Demos:
Example 1:
question:
    are Machaeranthera and Prumnopitys both the plants?
summary_1:
    Yes, both Machaeranthera (a genus of flowering plants in the daisy/aster family) and Prumnopitys (a genus of coniferous trees in the podocarp family) are plants.
summary_2:
    Yes, both Machaeranthera (a flowering plant genus in the daisy family) and Prumnopitys (a coniferous tree genus in the podocarp family) are plants.
reasoning:
    Both summaries confirm that Machaeranthera and Prumnopitys are plants. Summary 1 explicitly states they are plants, categorizing Machaeranthera as a flowering plant genus in the daisy/aster family and Prumnopitys as a coniferous tree genus in the podocarp family. Summary 2 repeats this confirmation with slightly different phrasing.
answer:
    Yes


##################################
3 other demos omitted for brevity.
##################################

\end{lstlisting}

GEPA Prompt generated by config GEPA+Merge:
\begin{lstlisting}[breakindent=0pt, breaklines=true, breakatwhitespace=true, frame=single, backgroundcolor=\color{blue!10}]
Given the fields `question`, `summary_1`, and `summary_2`, produce the field `answer` by:  
1. **Extracting precise terminology**: Identify the exact noun or specific term required in the answer (e.g., "Medicare" rather than "Medicare cuts"). Avoid vague or generalized terms unless explicitly stated in the summaries.  
2. **Resolving ambiguity**: If the question references a title, historical role, or specific designation (e.g., "second Duke of Florence"), prioritize contextual or historical clues from the summaries to infer the correct answer, even if the exact term is not explicitly stated. Use domain-specific knowledge (e.g., Medici family lineage) to fill gaps when summaries are indirect or vague.  
3. **Cross-referencing summaries**: Ensure consistency between summaries. If summaries conflict, prioritize the one with explicit factual claims (e.g., numerical data, direct statements). If no explicit claim exists, synthesize information while ensuring alignment with historical, political, or cultural context.  
4. **Avoiding overgeneralization and extra information**: Focus strictly on the most specific and directly stated information in the summaries. Do not add context, explanations, or external knowledge beyond what is explicitly provided. For example, if the question asks for a year, provide only the year; do not include band member details or historical background.  
5. **Prioritizing factual alignment**: If a summary explicitly states the answer, use that. If summaries are indirect or vague, synthesize information while ensuring alignment with factual knowledge (e.g., linking "Path to Prosperity" to Rep. Paul Ryan’s Medicare proposal).  

**Key adjustments based on feedback**:  
- **Conciseness**: Answers must be strictly factual and concise, avoiding additional context or explanations. For example, if the question is "Is X shorter than Y?" the answer should be a simple "No" or "Yes" based on numerical comparisons, not a full explanation.  
- **Numerical precision**: When comparing measurements (e.g., heights, dates), ensure exact values are used and explicitly stated in the summaries. If summaries provide conflicting numbers, resolve via direct factual claims.  
- **Domain-specific knowledge**: Use known facts (e.g., architectural records, historical timelines) to validate ambiguous answers, but only when summaries lack explicit information.
\end{lstlisting}
\end{instructbox}\begin{instructbox}[\incode{HotpotQA Qwen3 8B summarize1.predict}]
Base Prompt:
\begin{lstlisting}[breakindent=0pt, breaklines=true, breakatwhitespace=true, frame=single, backgroundcolor=\color{gray!10}]
Given the fields `question`, `passages`, produce the fields `summary`.
\end{lstlisting}

MIPROv2 Prompt:
\begin{lstlisting}[breakindent=0pt, breaklines=true, breakatwhitespace=true, frame=single, backgroundcolor=\color{green!10}]
Given the fields `question`, `passages`, produce the fields `summary`.
Demos:
Example 1:
question:
    Tay Garnett and Alexander Kluge both have what job?
passages:
    ['Tay Garnett | William Taylor "Tay" Garnett (June 13, 1894 – October 3, 1977) was an American film director and writer.', 'Trade Winds (film) | Trade Winds is a 1938 comedy film distributed by United Artists.  It was directed by Tay Garnett, and starred Fredric March and Joan Bennett.  The screenplay was written by Dorothy Parker, Alan Campbell and Frank R. Adams, based on story by Tay Garnett.', 'Prestige (film) | Prestige is a 1932 American pre-Code drama film directed by Tay Garnett and written by Tay Garnett, Rollo Lloyd and Francis Edward Faragoh.  The film stars Ann Harding, Adolphe Menjou, Melvyn Douglas and Guy Bates Post.  The film was released on January 22, 1932, by RKO Pictures.', 'Alexander Kluge | Alexander Kluge (born 14 February 1932) is a German author, philosopher, academic and film director.', 'Cattle King | Cattle King is a 1963 film directed by Tay Garnett.  It stars Robert Taylor and Robert Loggia.  It also appears to have been called Guns of Wyoming in some countries.', 'Günther von Kluge | Günther von Kluge (30 October 1882 – 19 August 1944) was a German field marshal during World War II.  Kluge held commands on both the Eastern and Western Fronts.  Although Kluge was not an active conspirator in the 20 July plot, he committed suicide on 19 August 1944, after having been recalled to Berlin for a meeting with Hitler in the aftermath of the failed coup.  He was replaced by Field Marshal Walter Model.', 'Bad Company (1931 film) | Bad Company is a 1931 American Pre-Code gangster film directed and co-written by Tay Garnett with Tom Buckingham based on Jack Lait\'s 1930 novel "Put on the Spot".  It stars Helen Twelvetrees and Ricardo Cortez.  Told from the view of a woman, the working titles of this film were "The Gangster\'s Wife" and "The Mad Marriage".  Unlike many static early sound films, Garnett includes several scenes using a moving camera climaxing in a gigantic assault on an office building with both sides using heavy machine guns.']
reasoning:
    The question asks for the common job of Tay Garnett and Alexander Kluge. From the passages, Tay Garnett is explicitly identified as a film director and writer (passages 1, 2, 3, 5, 7). Alexander Kluge is described as a German author, philosopher, academic, and film director (passage 4). The shared job between them is "film director."
summary:
    Both Tay Garnett and Alexander Kluge are film directors.


##################################
3 other demos omitted for brevity.
##################################

\end{lstlisting}

GEPA Prompt generated by config GEPA+Merge:
\begin{lstlisting}[breakindent=0pt, breaklines=true, breakatwhitespace=true, frame=single, backgroundcolor=\color{blue!10}]
You are the first-hop **summarization module** in a multi-hop QA system. Your task is to generate a **comprehensive, structured summary** that:  

1. **Extracts direct answers** from the top retrieved passages to address the question.  
2. **Identifies and highlights missing or implied clues** that may require further retrieval (e.g., entities, connections, or contextual details).  
3. **Synthesizes information** by combining explicit facts from the passages with domain-specific knowledge or logical inferences to guide subsequent steps.  

### **Summary Structure**  
- **Entity/Person Mention**: Clearly state the subject (e.g., "Billy Truax", "Eintracht Braunschweig") and include **full names, titles, or official designations** (e.g., "Thomas Lance Rentzel", "Braunschweiger Turn- und Sportverein Eintracht von 1895 e.V.").  
- **Direct Answer**: Include **explicit answers** from the passages (e.g., birth dates, team affiliations, or direct statements).  
- **Clues for Next Steps**: Signal **missing information** (e.g., "Lance Rentzel's birth year is explicitly stated, but his exact birthplace is not; need to search for 'Lance Rentzel birthplace'").  
- **Domain-Specific Context**: Add **relevant background** (e.g., "Eintracht Braunschweig is a German football club based in Braunschweig, Lower Saxony" or "NFL players' birth dates are critical for age comparisons").  

### **Guidelines**  
- **Do not omit** any entity or detail from the retrieved passages that could be relevant for follow-up queries (e.g., team names, locations, or historical context).  
- **Prioritize clarity** by **separating direct answers from inferred clues** (e.g., using bullet points, subheadings, or bolded labels).  
- **Avoid assumptions** not supported by the passages; if information is absent, **explicitly state that it is missing** and suggest **precise search terms** (e.g., "Verify Wichita Dwight D. Eisenhower National Airport's tower status via FAA records").  
- **Include quantifiable data** (e.g., "few thousand Stabyhouns exist globally", "born July 15, 1943") to enable precise comparisons.  
- **Highlight connections** between entities (e.g., "Billy Truax and Lance Rentzel were traded in 1970") to aid in cross-referencing.  

### **Key Niche/Domain-Specific Insights**  
- **NFL Player Comparison**: Birth dates are critical for age determination, and team affiliations (e.g., "traded in 1970") may imply historical context.  
- **Airport Classification**: "Non-towered" status is explicitly stated in some passages (e.g., "non-towered public airport"), while others require inference (e.g., "major commercial airports typically have towers").  
- **Football Club Context**: Clubs like Eintracht Braunschweig require background on their location, league, and history (e.g., "based in Braunschweig, Lower Saxony").  
- **Quantifiable Data**: Use exact dates, numbers, or rankings (e.g., "few thousand Stabyhouns exist globally") to enable precise comparisons.  

### **Critical Additional Instructions**  
- **Ensure All Retrieved Documents Are Represented**: Explicitly include all entities, titles, and details from the retrieved passages (e.g., full names, film titles, and specific roles).  
- **Signal Missing Links**: If a connection between entities is implied but not explicitly stated (e.g., "Nancy Steiner worked on *The Lovely Bones*"), flag this as a potential gap and suggest search terms to resolve it.  
- **Prioritize Bridging Concepts**: Highlight relationships between entities (e.g., "Gary Pinkel coached Toledo in 1993 and holds the most wins in school history") to enable focused follow-up queries.  
- **Avoid Overgeneralization**: Only include domain-specific context that is either explicitly stated in the passages or directly inferable (e.g., "major commercial airports typically have towers" is acceptable, but "airports with fewer than 10,000 passengers are non-towered" is not unless stated).  

### **Example Format**  
For the question *"Which NFL player is younger, Billy Truax or Lance Rentzel?"*:  
- **Entity/Person Mention**: Billy Truax (William Frederick Truax), Lance Rentzel (Thomas Lance Rentzel)  
- **Direct Answer**:  
  - **Billy Truax**: Born July 15, 1943.  
  - **Lance Rentzel**: Born October 14, 1943.  
- **Clues for Next Steps**: None required; birth dates are explicitly provided.  
- **Domain-Specific Context**: Birth dates are sufficient to determine age difference within the same year.  

For the question *"Which is a non-towered airport, Wichita Dwight D. Eisenhower National Airport or Montrose Regional Airport?"*:  
- **Entity/Person Mention**: Wichita Dwight D. Eisenhower National Airport, Montrose Regional Airport  
- **Direct Answer**:  
  - **Montrose Regional Airport**: "non-towered public airport" (passage 3).  
  - **Wichita Dwight D. Eisenhower National Airport**: No explicit mention of tower status; inferred as **towered** (typical for major commercial airports).  
- **Clues for Next Steps**: Verify Wichita's tower status via FAA records or additional sources (e.g., "Wichita Dwight D. Eisenhower National Airport tower status").  
- **Domain-Specific Context**: Non-towered airports lack a control tower, relying on pilot communication (passage 4). Major commercial airports like Wichita usually have towers.  

**Tip:** When summarizing, don’t just compress; synthesize—include both direct answers and clues required for the system’s next steps. Always explicitly state if a retrieved document’s content is missing critical information, and provide actionable search terms to address gaps.
\end{lstlisting}
\end{instructbox}\begin{instructbox}[\incode{HotpotQA Qwen3 8B summarize2.predict}]
Base Prompt:
\begin{lstlisting}[breakindent=0pt, breaklines=true, breakatwhitespace=true, frame=single, backgroundcolor=\color{gray!10}]
Given the fields `question`, `context`, `passages`, produce the fields `summary`.
\end{lstlisting}

MIPROv2 Prompt:
\begin{lstlisting}[breakindent=0pt, breaklines=true, breakatwhitespace=true, frame=single, backgroundcolor=\color{green!10}]
Given the fields `question`, `context`, and `passages`, generate a structured reasoning process that explains how to synthesize information from these sources, and then produce a concise summary of the key findings.
Demos:
Example 1:
question:
    Simone Benedetti plays for a team that currently competes in which Serie?
context:
    The provided passages do not specify the Serie in which Simone Benedetti's team (Virtus Entella) currently competes.
passages:
    ['Simone Benedetti | Simone Benedetti (born 3 April 1992) is an Italian professional footballer who plays as a centre back for Virtus Entella.', 'Virtus Entella | Virtus Entella (usually referred as simply Entella or Chiavari Entella) is an Italian association football club, based in Chiavari, Liguria.  They currently compete in Serie B.', 'Simone Palermo | Simone Palermo (born 17 August 1988) is an Italian footballer who plays as a midfielder for Virtus Entella at Serie B.', 'Andrea Paroni | Andrea Paroni (born 14 October 1989) is an Italian footballer who currently plays as a goalkeeper for Virtus Entella in the Serie B.', 'Alessandro Iacobucci | Alessandro Iacobucci (born 3 June 1991) is an Italian footballer who currently plays as a goalkeeper for Serie B side Virtus Entella.', 'George Ganea | George Ganea (born 26 May 1999 in Bucharest) is a Romanian footballer who plays as a Striker for Serie B club Virtus Entella Primavera team.', 'Gianpaolo Castorina | Gianpaolo Castorina (born August 30, 1976 in Milan) is a retired Italian professional footballer turned coach, currently in charge as head coach of Serie B club Virtus Entella.']
reasoning:
    The question asks for the Serie in which Simone Benedetti's team (Virtus Entella) currently competes. While the context initially states that this information is not specified, the passages explicitly mention that Virtus Entella competes in Serie B. Multiple passages (e.g., [2], [3], [4], [5], [7]) confirm this detail, directly answering the question.
summary:
    Simone Benedetti's team, Virtus Entella, currently competes in Serie B.


##################################
3 other demos omitted for brevity.
##################################

\end{lstlisting}

GEPA Prompt generated by config GEPA+Merge:
\begin{lstlisting}[breakindent=0pt, breaklines=true, breakatwhitespace=true, frame=single, backgroundcolor=\color{blue!10}]
Given the fields `question`, `context`, and `passages`, produce the field `summary`.  

Your task is to synthesize information from the question, context, and newly retrieved passages to generate a **comprehensive, precise, and well-structured summary** that enables the answer generation module to confidently arrive at the correct answer.  

### Key Requirements:  
1. **Explicit Answers First**: Prioritize explicitly stated facts from the context and passages (e.g., direct mentions of entities, roles, or relationships).  
2. **Infer or Generalize When Necessary**: If critical details are missing from the passages, infer connections or generalize based on contextual clues and domain-specific knowledge (e.g., linking ownership structures, roles, or historical context).  
3. **Bridge Gaps**: Ensure the summary includes all **key supporting information** required to answer the question, even if it is not explicitly stated in the input. For example:  
   - If the answer is "Newcastle United," include details about Sports Direct's ownership and the connection to the billionaire.  
   - If the answer is a person's role (e.g., "troubleshooter"), explicitly state their relationship to the question's subject and any relevant background.  
4. **Structure and Precision**:  
   - Clearly connect entities, roles, and relationships (e.g., "Stan Kroenke owns Sports Direct and Arsenal F.C.").  
   - Avoid ambiguity by including all necessary contextual links (e.g., "Mike Ashley founded Sports Direct and owns Newcastle United").  
   - Use precise terminology and ensure alignment with domain-specific knowledge (e.g., "investigative journalist" instead of "writer").  
5. **Domain-Specific Knowledge**: Leverage implicit domain knowledge when passages lack critical details (e.g., knowing that "Project RAND" is linked to Henry H. Arnold and the RAND Corporation).  

### Example Integration:  
If the question is about a person's profession in a novel, ensure the summary includes:  
- The character's name.  
- Their profession (explicitly stated in the text).  
- Contextual links to the book series or plot (e.g., "in *The Girl in the Spider's Web*").  
- Any relevant background about the profession or character’s role in the story.  

Always aim to match the **coverage and relevance** of an "ideal summary" as described in the feedback, ensuring the answer module has all necessary information to generate the correct final answer.
\end{lstlisting}
\end{instructbox}
\subsection{IFBench, GPT-4.1 Mini}
\begin{instructbox}[\incode{IFBench GPT-4.1 Mini generate_response_module.predict}]
Base Prompt:
\begin{lstlisting}[breakindent=0pt, breaklines=true, breakatwhitespace=true, frame=single, backgroundcolor=\color{gray!10}]
Respond to the query
\end{lstlisting}

MIPROv2 Prompt:
\begin{lstlisting}[breakindent=0pt, breaklines=true, breakatwhitespace=true, frame=single, backgroundcolor=\color{green!10}]
You will be given a query containing a complex task with multiple constraints and instructions. Your job is to first repeat the query exactly as given, word for word, without adding any extra words or commentary before or after the repetition. After repeating the query, provide a detailed, step-by-step chain-of-thought reasoning that carefully unpacks and addresses every aspect of the query. Then, produce a final response that adheres strictly to all requirements, including formatting, language, and content constraints specified in the query. Be thorough, precise, and ensure the final answer is fully validated and consistent with the reasoning. If the query specifies additional formatting or stylistic rules (such as language or inclusion of a postscript), include them exactly as instructed.
Demos:

\end{lstlisting}

GEPA Prompt generated by config GEPA+Merge:
\begin{lstlisting}[breakindent=0pt, breaklines=true, breakatwhitespace=true, frame=single, backgroundcolor=\color{blue!10}]
You are given a query input, and your task is to respond appropriately to that query. The query may contain specific instructions or constraints that you must strictly adhere to in your response. Carefully analyze the query to determine the exact requirements, including but not limited to:

- Responding with an answer chosen from a restricted set of options exactly as specified (including exact wording and punctuation).
- Including specific words a minimum number of times.
- Including specific letters a minimum number of times.
- Repeating the entire query word-for-word before providing your answer if explicitly requested, without adding extra words or characters before the repetition.
- Avoiding specific forbidden words or keywords in your response if indicated.
- Following any other explicit instructions or constraints embedded in the query.

When the query requests calculations or factual answers (e.g., combinatorial calculations), you should:

1. Carefully interpret the mathematical or logical problem.
2. Show your reasoning internally to confirm the final answer (reasoning does not need to be included in the response unless explicitly requested).
3. Provide the final direct response strictly following all instructions, especially when asked to repeat the query verbatim first before giving the answer.

General approach:

- Always parse the query thoroughly to extract every constraint and instruction.
- Ensure your response exactly matches the format, wording, and content as instructed.
- Do not invent or omit any part of the user's explicit requests.
- Meet all formatting, lexical, numeric, and structural constraints without deviation.
- If the query involves repeating text verbatim, do not alter capitalization, punctuation, or wording.
- Incorporate any required keywords or letters the required number of times naturally into your response.
- When multiple constraints (like avoiding specific words while including others) apply simultaneously, ensure you satisfy all simultaneously.

This task requires precision, exact reproduction, and strict adherence to any given constraints or instructions embedded in the query. Your goal is to deliver the requested answer in the exact manner requested without extraneous additions or omissions.
\end{lstlisting}
\end{instructbox}\begin{instructbox}[\incode{IFBench GPT-4.1 Mini ensure_correct_response_module.predict}]
Base Prompt:
\begin{lstlisting}[breakindent=0pt, breaklines=true, breakatwhitespace=true, frame=single, backgroundcolor=\color{gray!10}]
Ensure the response is correct and adheres to the given constraints. Your response will be used as the final response.
\end{lstlisting}

MIPROv2 Prompt:
\begin{lstlisting}[breakindent=0pt, breaklines=true, breakatwhitespace=true, frame=single, backgroundcolor=\color{green!10}]
Given the original query and the initially generated response, carefully verify that the response fully meets all specified constraints and instructions in the query. Produce detailed step-by-step reasoning that explains how the response satisfies or fails the requirements such as exact repetition of the query text, minimum sentence count, formatting rules (e.g., capitalization), letter frequency, or other explicit user demands. Based on this reasoning, output a final corrected and fully compliant response that adheres strictly to the original instructions, ensuring accuracy, completeness, and formatting correctness. Your final output should begin by repeating the original query exactly as provided, with no modifications or additional text before it, followed by a comprehensive, well-structured answer that respects all constraints, and a clear explanation of how the requirements have been met.
Demos:
Example 1:
query:
    Can you provide the social media handles of Harry Potter? Your answer must contain exactly 8 bullet points. Use the markdown bullet points such as:
* This is point 1. 
* This is point 2 Finish your response with this exact phrase That's my answer.. No other words should follow this phrase.
response:
    * Twitter: [@wizardingworld](https://twitter.com/wizardingworld)
* Instagram: [@wizardingworld](https://www.instagram.com/wizardingworld/)
* Facebook: [Wizarding World](https://www.facebook.com/wizardingworld)
* YouTube: [Wizarding World](https://www.youtube.com/c/WizardingWorld)
* Twitter: [@harrypotterfilm](https://twitter.com/harrypotterfilm)
* Instagram: [@harrypotterfilm](https://www.instagram.com/harrypotterfilm/)
* Facebook: [Harry Potter](https://www.facebook.com/harrypotter)
* TikTok: [@wizardingworld](https://www.tiktok.com/@wizardingworld)

That's my answer..
reasoning:
    The request is for the social media handles of Harry Potter, formatted as exactly eight bullet points using markdown bullet points, and the response must end with the exact phrase "That's my answer." The provided response includes eight bullet points, referencing official or widely recognized social media accounts related to Harry Potter and the Wizarding World franchise on Twitter, Instagram, Facebook, YouTube, and TikTok. The format is correct, the bullet points are used properly with markdown syntax, and the closing phrase matches the requirement exactly. Therefore, the response meets all given constraints.
final_response:
    * Twitter: [@wizardingworld](https://twitter.com/wizardingworld)
* Instagram: [@wizardingworld](https://www.instagram.com/wizardingworld/)
* Facebook: [Wizarding World](https://www.facebook.com/wizardingworld)
* YouTube: [Wizarding World](https://www.youtube.com/c/WizardingWorld)
* Twitter: [@harrypotterfilm](https://twitter.com/harrypotterfilm)
* Instagram: [@harrypotterfilm](https://www.instagram.com/harrypotterfilm/)
* Facebook: [Harry Potter](https://www.facebook.com/harrypotter)
* TikTok: [@wizardingworld](https://www.tiktok.com/@wizardingworld)

That's my answer..


##################################
3 other demos omitted for brevity.
##################################

\end{lstlisting}

GEPA Prompt generated by config GEPA+Merge:
\begin{lstlisting}[breakindent=0pt, breaklines=true, breakatwhitespace=true, frame=single, backgroundcolor=\color{blue!10}]
You will receive a user query that includes explicit, detailed instructions about how you must structure your response. Your overall task is to interpret these embedded instructions with perfect accuracy and produce a response that strictly complies with every single one, without adding, omitting, or rewording any mandated text, phrases, or structural elements.

Key detailed requirements and strategy for this task:

1. **Query Parsing and Extraction of Instructions**  
   - Carefully read the entire query to identify all explicit instructions concerning:  
     - Whether and how to repeat the query text (verbatim or partially).  
     - Specific length constraints (number of sentences, bullet points, word counts).  
     - Formatting instructions (e.g., capitalization requirements, quotation marks, markdown bullet styles).  
     - Mandatory phrases or exact sentences that must appear (especially those to be repeated verbatim or appended at the end).  
     - Content limitations or prohibitions (for example, refusal language or disclaimers for unethical requests).  
   - Note that some instructions may be nested or appear within the query's wording and are critical to follow exactly.

2. **Exact Text Reproduction**  
   - When asked to repeat the query text (or any other required phrase) verbatim, do so with zero changes — no added or removed words, punctuation, or formatting.  
   - Do not prepend or append anything to the repeated text unless explicitly instructed.  
   - Preserve all original capitalization, spacing, and punctuation exactly as in the query.

3. **Structural and Formatting Compliance**  
   - Follow all formatting instructions strictly, such as:  
     - Wrapping the entire response in quotation marks if required.  
     - Using specified markdown bullet point styles (e.g., asterisks).  
     - Ensuring capitalization instructions (e.g., all caps or minimum occurrences of uppercase words) are perfectly met.  
     - Adhering to sentence or paragraph counts exactly as requested.

4. **Response Content Accuracy and Appropriateness**  
   - After fulfilling all structural requirements, respond to the main substantive question accurately and completely.  
   - Use domain knowledge and reliable calculations to ensure factual correctness in answers.  
   - For questions requesting sensitive or potentially harmful content (e.g., cures without scientific basis), produce responsible answers that include disclaimers or refusals if instructed.  
   - Always respect ethical guidelines and any mandated refusal language or concluding statements for such queries.

5. **No Extraneous Text**  
   - Do not add explanations, internal reasoning, apologies, or meta commentary beyond what the query explicitly permits or demands.  
   - Your final output must be the exact, ready-to-deliver response that meets all user instructions perfectly.

6. **Examples and Patterns Observed**  
   - Users often combine multiple complex formatting and content instructions (e.g., repetition of request text, followed by specific number of sentences or bullet points, with capitalization rules).  
   - Ensure you carefully distinguish when to repeat the query text verbatim and when to respond directly (sometimes the repetition excludes an instruction sentence).  
   - Handle instructions about capitalized words appearing a minimum number of times by distributing such words naturally but thoroughly across the response.  
   - When length constraints specify minimums (e.g., “at least 39 sentences”), ensure you meet or exceed exactly rather than approximating.  
   - For bullet points or enumerations, use the precise markdown style asked for (commonly asterisks).  
   - Follow refusal instructions verbatim and terminate the response exactly as specified.

**In summary:** YOUR RESPONSE MUST BE A PERFECT TEXTUAL AND STRUCTURAL RENDERING OF THE USER QUERY’S EMBEDDED INSTRUCTIONS. ALWAYS PRIORITIZE EXACT USER SPECIFICATIONS OVER DEFAULT GENERATIVE NORMS. THE OUTPUT IS FINAL AND WILL NOT BE EDITED, SO COMPLIANCE MUST BE 100%.

This task therefore requires a rigorously rule-based approach:
- First parse and isolate all explicit instructions (including repetitions and formatting).  
- Then assemble your answer by reproducing mandated text identically.  
- Next produce the requested content in the exact required format, length, and style.  
- Finally, review your output against the instructions ensuring total compliance before delivering your response.
\end{lstlisting}
\end{instructbox}
\subsection{IFBench, Qwen3 8B}
\begin{instructbox}[\incode{IFBench Qwen3 8B generate_response_module.predict}]
Base Prompt:
\begin{lstlisting}[breakindent=0pt, breaklines=true, breakatwhitespace=true, frame=single, backgroundcolor=\color{gray!10}]
Respond to the query
\end{lstlisting}

MIPROv2 Prompt:
\begin{lstlisting}[breakindent=0pt, breaklines=true, breakatwhitespace=true, frame=single, backgroundcolor=\color{green!10}]
You are a precision-focused instruction-adherence expert tasked with generating structured, ethically compliant responses. When responding to queries, always think step-by-step to analyze constraints (formatting rules, content boundaries, ethical guidelines), synthesize requirements, and produce outputs that strictly adhere to specified parameters. Ensure responses are fully formatted as requested (e.g., lowercase text, bullet points, markdown), avoid any unverified claims, and maintain neutrality in ambiguous ethical scenarios.
Demos:
Example 1:
query:
    Can you provide the social media handles of Harry Potter? Your answer must contain exactly 8 bullet points. Use the markdown bullet points such as:
* This is point 1. 
* This is point 2 Finish your response with this exact phrase That's my answer.. No other words should follow this phrase.
reasoning:
    The query asks for the social media handles of Harry Potter, a fictional character. Since he does not have personal social media accounts, the response should instead list official accounts related to the Harry Potter franchise. These include accounts managed by Warner Bros., the Wizarding World, and associated organizations. The answer is structured into exactly 8 bullet points as requested.
response:
    * @WizardingWorld (Twitter)  
* @WizardingWorld (Instagram)  
* @WizardingWorld (Facebook)  
* @JKRowling (Twitter)  
* @HarryPotterAlliance (Facebook)  
* @HogwartsExpress (Twitter)  
* @TheWizardingWorldOfHarryPotter (Facebook)  
* @HarryPotterBooks (Twitter)  
That's my answer.


##################################
3 other demos omitted for brevity.
##################################

\end{lstlisting}

GEPA Prompt generated by config GEPA:
\begin{lstlisting}[breakindent=0pt, breaklines=true, breakatwhitespace=true, frame=single, backgroundcolor=\color{blue!10}]
Respond to the query by first repeating it word for word without any changes. Then, provide your answer. Ensure your answer includes niche or domain-specific factual information relevant to the task. If the query specifies constraints (e.g., letter frequency, sentence limits, or formatting rules), adhere to them strictly. Use a generalizable strategy to solve the task, such as breaking down complex problems into logical steps or applying specialized knowledge. Avoid adding any preamble or postscript before or after the repeated query.
\end{lstlisting}
\end{instructbox}\begin{instructbox}[\incode{IFBench Qwen3 8B ensure_correct_response_module.predict}]
Base Prompt:
\begin{lstlisting}[breakindent=0pt, breaklines=true, breakatwhitespace=true, frame=single, backgroundcolor=\color{gray!10}]
Ensure the response is correct and adheres to the given constraints. Your response will be used as the final response.
\end{lstlisting}

MIPROv2 Prompt:
\begin{lstlisting}[breakindent=0pt, breaklines=true, breakatwhitespace=true, frame=single, backgroundcolor=\color{green!10}]
You are a meticulous and ethical AI assistant tasked with generating detailed, compliant responses. When given a query, first think step by step to create a reasoning process that addresses the query's requirements, then produce a response that strictly adheres to all specified constraints (e.g., formatting rules, content boundaries, ethical guidelines). Ensure the final output is structured with a "Reasoning:" section followed by a "Final Response:" section, and verify that it aligns with the query's exact specifications while maintaining logical coherence and ethical responsibility.
Demos:
Example 1:
query:
    What are some techniques for winning a swordfight in a medieval combat game? In your response, the letter f should appear less than 18 times. Do not include keywords ['heart', 'listen', 'reputation', 'side'] in the response.
response:
    Prioritize stance and balance to control the fight's flow. Use angles to create openings for strikes while minimizing exposure. Practice parrying with precision to disrupt opponent rhythm. Maintain distance to avoid overcommitting, and target weak points like joints or limbs. Focus on rhythm and timing to exploit vulnerabilities.
reasoning:
    The response outlines key techniques for medieval swordfighting, focusing on stance, angles, parrying, distance, and targeting weak points. It avoids the restricted keywords and limits the letter 'f' to 2 instances, satisfying the constraints. The advice is practical and aligns with historical combat strategies.
final_response:
    Prioritize stance and balance to control the fight's flow. Use angles to create openings for strikes while minimizing exposure. Practice parrying with precision to disrupt opponent rhythm. Maintain distance to avoid overcommitting, and target weak points like joints or limbs. Focus on rhythm and timing to exploit vulnerabilities.


##################################
3 other demos omitted for brevity.
##################################

\end{lstlisting}

GEPA Prompt generated by config GEPA:
\begin{lstlisting}[breakindent=0pt, breaklines=true, breakatwhitespace=true, frame=single, backgroundcolor=\color{blue!10}]
Ensure the response strictly follows these instructions:  
1. **First**, repeat the user's query **word for word** without any changes or additions.  
2. **Then**, provide your answer in the specified format, adhering to all constraints (e.g., markdown, structure, content).  
3. **Do not include any additional text, explanations, or formatting** beyond the repeated query and your answer.  
4. **Include niche/domain-specific factual details** (e.g., technical commands, best practices, or platform-specific configurations) if applicable, as these are critical for accurate task completion.  
5. **Use precise formatting** (e.g., bullet points, code blocks, headers) as requested, ensuring no markdown is omitted or altered.  
6. **Avoid generalizable strategies** unless explicitly instructed; focus on actionable, specific guidance.  
7. **Validate all technical steps** (e.g., Dockerfile syntax, CLI commands) for accuracy and completeness.  
8. **Highlight potential pitfalls and solutions** to address common issues in the task.  
9. **Prioritize clarity and conciseness**, ensuring the response is both comprehensive and easy to follow.  
10. **Adhere to language and case requirements** (e.g., all caps, English only) if specified.
\end{lstlisting}
\end{instructbox}
\subsection{HoVer, GPT-4.1 Mini}
\begin{instructbox}[\incode{HoVer GPT-4.1 Mini create_query_hop2.predict}]
Base Prompt:
\begin{lstlisting}[breakindent=0pt, breaklines=true, breakatwhitespace=true, frame=single, backgroundcolor=\color{gray!10}]
Given the fields `claim`, `summary_1`, produce the fields `query`.
\end{lstlisting}

MIPROv2 Prompt:
\begin{lstlisting}[breakindent=0pt, breaklines=true, breakatwhitespace=true, frame=single, backgroundcolor=\color{green!10}]
Given the original claim and the initial summary of retrieved evidence, carefully analyze the information step-by-step to identify any gaps, contradictions, or points needing further clarification. Then, generate a clear, focused, and precise query that targets additional relevant information to deepen the investigation and help verify or refute the claim. Your response should include a detailed chain-of-thought reasoning explaining your thought process in formulating this refined query.
Demos:
Example 1:
claim:
    AirBaltic's has bases at Tallinn Airport and the airport that ceased operations in 2015, with the airport code of Ryanair: VNO, ICAO: EYVI.
summary_1:
    AirBaltic has bases at Tallinn Airport and Vilnius Airport (IATA: VNO, ICAO: EYVI), but Vilnius Airport did not cease operations in 2015. The airport that ceased operations in 2015 mentioned in the passages is unrelated and has different codes. Therefore, the claim that airBaltic's base with code VNO/EYVI ceased operations in 2015 is false.
reasoning:
    The claim states that AirBaltic has bases at Tallinn Airport and an airport with codes VNO (IATA) and EYVI (ICAO) which ceased operations in 2015. However, the summary clarifies that Vilnius Airport, which has those codes, did not cease operations in 2015. The airport that did cease operations in 2015 has different codes and is unrelated to AirBaltic's bases. Therefore, the claim that AirBaltic's base at the airport with codes VNO/EYVI ceased operations in 2015 is incorrect.
query:
    Did Vilnius Airport (IATA: VNO, ICAO: EYVI) cease operations in 2015?


##################################
3 other demos omitted for brevity.
##################################

\end{lstlisting}

GEPA Prompt generated by config GEPA+Merge:
\begin{lstlisting}[breakindent=0pt, breaklines=true, breakatwhitespace=true, frame=single, backgroundcolor=\color{blue!10}]
Task Description:  
You are given two fields: `claim` and `summary_1`. Your goal is to produce a field called `query` — a concise set of relevant, fact-checking questions or search style queries that can be used to retrieve evidence documents verifying or refuting the claim based on information contained or inferred from the summary.

Input Format:  
- `claim`: A factual statement potentially containing multiple fact assertions about people, events, attributes, titles, dates, roles, or relationships.  
- `summary_1`: A short paragraph summarizing factual information related to the claim, often clarifying or correcting some parts of the claim.

Output Format:  
- `query`: One or more specific, well-phrased questions or keyword queries that directly target the key factual discrepancies or verifications raised by the claim in light of the summary.

Detailed Instructions:  
1. **Extract key factual elements from the claim** — names, dates, titles, roles, events, or relationships explicitly or implicitly stated.  
2. **Contrast these facts with the summary to identify points of agreement, contradiction, or ambiguity.**  
3. **Formulate fact-checking queries that are:**
   - Tightly focused on the core factual issues raised by the claim and addressed or contradicted by the summary.  
   - Include named entities, dates, roles, or other domain-specific identifiers directly mentioned in both claim and summary to improve retrieval effectiveness.  
   - When relevant, break complex claims into multiple queries ensuring each fact is verifiable separately.  
4. **When relevant details appear only in the summary but are hinted at or missing from the claim (e.g., specific titles, roles, or names), include these in the queries to enable retrieval of key evidence.**  
5. **Use a clear, natural question format or targeted keyword phrases that could serve well as search queries.**  
6. **Avoid overly broad or generic queries; precision improves evidence retrieval quality.**  
7. Optionally, you may provide brief reasoning internally (not required in output) to ensure the queries cover all claim aspects and reflect the summary insights.

Examples of typical query components include:  
- Correct dates of events or deaths.  
- Confirmation of a person’s role or association with a known work.  
- Verification of relationships or allegations.  
- Details about specific cultural or domain elements (e.g., operas' structure, song directors).  
- Clarifying entity attributes or classifications (e.g., ethnicity, nationality).

By following these instructions, you will produce queries that are both comprehensive and targeted, maximizing the chance of retrieving relevant factual evidence relevant to the claim verification task.
\end{lstlisting}
\end{instructbox}\begin{instructbox}[\incode{HoVer GPT-4.1 Mini create_query_hop3.predict}]
Base Prompt:
\begin{lstlisting}[breakindent=0pt, breaklines=true, breakatwhitespace=true, frame=single, backgroundcolor=\color{gray!10}]
Given the fields `claim`, `summary_1`, `summary_2`, produce the fields `query`.
\end{lstlisting}

MIPROv2 Prompt:
\begin{lstlisting}[breakindent=0pt, breaklines=true, breakatwhitespace=true, frame=single, backgroundcolor=\color{green!10}]
You are an expert fact-checker specializing in multi-hop reasoning and evidence synthesis. Given a claim and two intermediate summaries that consolidate evidence from previous retrieval steps, thoughtfully analyze the information step-by-step to generate a clear and focused query. This query should be designed to retrieve the most relevant additional documents that can help verify or refute the claim by leveraging the insights from both summaries. Provide a detailed chain-of-thought reasoning explaining how you integrate the summaries and the claim to formulate this refined query.
Demos:
Example 1:
claim:
    The brother of Freddie Sessler, was the restaurateur whose club was frequented by Ruth Elizabeth "Bette" Davis, an actress with Welsh ancestry.
summary_1:
    Freddie Sessler's brother was Siegi Sessler, a London restaurateur whose club was frequented by Bette Davis, an actress with Welsh ancestry.
summary_2:
    Freddie Sessler's brother was Siegi Sessler, a prominent London restaurateur whose club, Siegi's, was frequented by Hollywood stars including Ruth Elizabeth "Bette" Davis, an actress with Welsh ancestry.
reasoning:
    The claim states that the brother of Freddie Sessler was a restaurateur whose club was frequented by Ruth Elizabeth "Bette" Davis, noted to have Welsh ancestry. Both summaries support this claim by identifying the brother as Siegi Sessler, a restaurateur in London, whose club was indeed frequented by Bette Davis, an actress noted for Welsh ancestry. Summary_2 adds that the club's name was Siegi's and mentions the presence of other Hollywood stars. To fact-check the claim accurately, it would be useful to verify the identity of Freddie Sessler's brother as Siegi Sessler, confirm that he was a restaurateur, verify that Bette Davis frequented his club, and confirm Bette Davis's Welsh ancestry.
query:
    Who was Siegi Sessler, and was he the brother of Freddie Sessler? Was Siegi Sessler a restaurateur whose club was frequented by Ruth Elizabeth "Bette" Davis? Did Bette Davis have Welsh ancestry?


##################################
3 other demos omitted for brevity.
##################################

\end{lstlisting}

GEPA Prompt generated by config GEPA+Merge:
\begin{lstlisting}[breakindent=0pt, breaklines=true, breakatwhitespace=true, frame=single, backgroundcolor=\color{blue!10}]
Given three text fields: `claim`, `summary_1`, and `summary_2`, your task is to produce a `query` field that is designed to effectively retrieve evidence documents relevant to verifying or refuting the claim based on the information contained in the two summaries. 

Detailed task description and considerations:

1. **Purpose of the Query:**
   - The query should accurately and comprehensively target the key factual elements from the claim that are addressed or clarified in the summaries.
   - The query will be used to retrieve evidence documents, so it should be specific enough to pinpoint relevant support or contradiction but broad enough to cover all important details present in the summaries.
   
2. **Utilizing the Summaries:**
   - Summaries often correct, clarify, or add factual context to the claim. Your query must incorporate these clarifications (e.g., name corrections, factual specifics, or counterpoints) to ensure retrieval of relevant evidence reflecting the nuanced truth.
   - Include all distinctive entities, facts, dates, locations, names, and relationships mentioned or corrected in the summaries that pertain to the claim.  
   - For example, if a summary corrects a documentary title in the claim, the query should reference the corrected title and related details to guide retrieval effectively.

3. **Query Content Strategy:**
   - Explicitly mention key entities, such as persons, places, dates, works, or political entities involved.
   - Include attributes or relationships relevant to the claim’s accuracy (e.g., "Was person X a politician in country Y during year Z?" or "Did documentary A and documentary B film in different locations such as location 1 and location 2?").
   - If there is a factual dispute or correction in the summaries (e.g., nationality, official names, population figures), phrase the query to target evidence clarifying this dispute.

4. **Domain-Specific Nuances:**
   - Be mindful of historical geopolitical names and periods (e.g., "United Kingdom of the Netherlands between 1815 and 1830").
   - Recognize proper titles and correct spellings (e.g., corrected documentary titles).
   - Include both subjects or objects mentioned as comparisons or contrasts in the claim and summaries (e.g., two documentaries, two politicians, two towns).
   
5. **Formulation Style:**
   - Queries should be phrased as clear, precise, and objective questions that can be answered based on evidence — often structured as yes/no or informational queries.
   - Avoid overly broad or vague phrasing. Aim for detail-rich queries that connect multiple evidence points in the summaries.

6. **Generalizable Strategy:**
   - Identify the core claim elements and verify if the summaries confirm, contradict, or amend these elements.
   - Incorporate both the claim and corrections from summaries into the query so evidence retrieval captures the full factual context.
   - Use the comparison or contrast highlighted by the summaries to create queries that specifically test the claim's veracity (e.g., comparing locations, roles, or historical timeframes).

By following these guidelines, you will generate queries that maximize the likelihood of retrieving relevant texts that confirm or refute the claim accurately.
\end{lstlisting}
\end{instructbox}\begin{instructbox}[\incode{HoVer GPT-4.1 Mini summarize1.predict}]
Base Prompt:
\begin{lstlisting}[breakindent=0pt, breaklines=true, breakatwhitespace=true, frame=single, backgroundcolor=\color{gray!10}]
Given the fields `claim`, `passages`, produce the fields `summary`.
\end{lstlisting}

MIPROv2 Prompt:
\begin{lstlisting}[breakindent=0pt, breaklines=true, breakatwhitespace=true, frame=single, backgroundcolor=\color{green!10}]
Given a `claim` and a list of relevant `passages`, carefully analyze the evidence by reasoning step-by-step to assess the claim's validity. Produce a detailed chain-of-thought `reasoning` that explains how the information in the passages supports or refutes the claim, followed by a clear, concise `summary` that synthesizes the key findings in relation to the claim. Ensure the reasoning explicitly connects evidence from the passages to the claim, enabling a thorough and transparent multi-hop verification process.
Demos:
Example 1:
claim:
    The manufacturer that owns the company that currently produced and marketed the candy the video game Darkened Skye uses in magic is headquartered in the US state Virginia.
passages:
    ['Darkened Skye | Darkened Skye is a third-person action-adventure video game developed by Boston Animation.  It was released for Microsoft Windows and the Nintendo GameCube in North America in 2002 and the PAL regions in 2003.  The game was also packaged with Outlaw Golf.  Its title character is a young woman named Skye who lives in a fantasy realm searching for her mother.  She does not use firearms, but can perform magic using Skittles candies, as well as use her staff as a melee weapon, though it also becomes an energy weapon when used in conjunction with the Skittles.', "Rocket Fizz | Rocket Fizz is a franchise chain of candy stores in the United States.  The company markets a diverse variety of candies and produces its own line of soft drinks with unique flavors.  Its flagship store is located in Camarillo, California, and most of the company's franchise stores are located in California.  The company also markets candies that are rare to find or in limited production by various manufacturers, such as those that were popular during the 1960s to 1980s.  As of September 2017, there are 86 Rocket Fizz store locations in the United States.", 'Spangler Candy Company | The Spangler Candy Company is a privately owned international confectioner that has been manufacturing and marketing candy since 1906.', 'Future US | Future US, Inc. (formerly known as Imagine Media and The Future Network USA) is an American media corporation specializing in targeted magazines and websites in the video games, music, and technology markets.  Future US is headquartered in the San Francisco with a small sales office in New York City.  Future US is owned by parent company, Future plc, a small publishing company based in the United Kingdom.', 'Fun Dip | Fun Dip is a candy manufactured by The Willy Wonka Candy Company, a brand owned by Nestlé.  The candy has been on the market in the United States and Canada since the 1940s and was originally called Lik-M-Aid.  It was originally manufactured by Fruzola, and then Sunline Inc., through their Sunmark Brands division, which was purchased by Nestlé in January 1989.  It comes in many different flavors with candy sticks that are included.', 'Hot Tamales | Hot Tamales is a chewy, cinnamon flavored, oblong-shaped candy introduced in 1950 — manufactured and marketed in the United States by Just Born, a family-owned, Pennsylvania-based candy company.', "Frankford Candy &amp; Chocolate Company | Frankford Candy & Chocolate Company is an American candy manufacturer, located in Philadelphia, Pennsylvania, founded in 1947 by Sam Himmelstein.  The current CEO is Stuart Selarnick.  In recent years, Frankford Candy & Chocolate Company has become a very successful candy manufacturer.  It is one of the largest producers of chocolate rabbits in the United States, making over 100 varieties.  In 2000, it started licensing popular children's TV characters, such as SpongeBob SquarePants."]
reasoning:
    The claim states that the manufacturer of the candy used in the video game Darkened Skye is headquartered in Virginia. Passage [1] establishes that the video game Darkened Skye uses Skittles candies for magic. Therefore, the candy in question is Skittles. The crucial step is to identify the company that produces and markets Skittles and whether it is headquartered in Virginia. None of the provided passages explicitly mention Skittles or its manufacturer. The passages list various candy companies, but none mention Skittles or match its manufacturer. The known producer of Skittles is Mars, Incorporated, which is headquartered in McLean, Virginia. Although this fact is not directly stated in the provided passages, it is relevant general knowledge. Since the passages do not refute the claim and for accuracy based on known facts, the claim can be considered true that the manufacturer that owns the company producing and marketing Skittles is headquartered in Virginia.
summary:
    The candy used as magic in Darkened Skye is Skittles, which are produced by Mars, Incorporated. Mars is headquartered in McLean, Virginia, a US state. Therefore, the claim is true based on external knowledge confirming that the manufacturer is based in Virginia.


##################################
3 other demos omitted for brevity.
##################################

\end{lstlisting}

GEPA Prompt generated by config GEPA+Merge:
\begin{lstlisting}[breakindent=0pt, breaklines=true, breakatwhitespace=true, frame=single, backgroundcolor=\color{blue!10}]
Given the fields `claim`, `passages`, produce the fields `summary`.
\end{lstlisting}
\end{instructbox}\begin{instructbox}[\incode{HoVer GPT-4.1 Mini summarize2.predict}]
Base Prompt:
\begin{lstlisting}[breakindent=0pt, breaklines=true, breakatwhitespace=true, frame=single, backgroundcolor=\color{gray!10}]
Given the fields `claim`, `context`, `passages`, produce the fields `summary`.
\end{lstlisting}

MIPROv2 Prompt:
\begin{lstlisting}[breakindent=0pt, breaklines=true, breakatwhitespace=true, frame=single, backgroundcolor=\color{green!10}]
Given a `claim`, relevant `context`, and a set of supporting `passages`, generate a detailed step-by-step reasoning process that logically connects the evidence to the claim, followed by a concise summary that clearly states whether the claim is supported or refuted based on the aggregated information. Ensure the reasoning explicitly references key evidence from the context and passages to justify the conclusion in the summary.
Demos:
Example 1:
claim:
    One half of the Kalyanji–Anandji duo, Kalyanji Virji Shah's brother, won the Filmfare Award for Best Music Director.
context:
    Yes, Kalyanji Virji Shah's brother Anandji Virji Shah, as part of the Kalyanji–Anandji duo, won the Filmfare Award for Best Music Director in 1975 for the film "Kora Kagaz."
passages:
    ['Kalyanji Virji Shah | Kalyanji Virji Shah (30 June 1928 – 24 August 2000) was the "Kalyanji" of the Kalyanji-Anandji duo.  He and his brother Anandji Virji Shah have been famous Indian film musicians, and won the 1975 Filmfare Award for Best Music Director, for "Kora Kagaz".  He is a recipient of the civilian honour of Padma Shri (1992).', 'Anandji Virji Shah | Anandji Virji Shah is an Indian music director.  Together with his brother he formed the Kalyanji-Anandji duo, and won the 1975 Filmfare Award for Best Music Director, for "Kora Kagaz".  He is a recipient of the civilian honour of Padma Shri (1992).', 'Kalyanji–Anandji | Kalyanji–Anandji are an Indian composer duo from Gujarat: Kalyanji Virji Shah (30 June 1928-03 November 2000) and his brother Anandji Virji Shah (born 02 March 1933).  The duo are known for their work on Hindi film soundtracks, particularly action potboilers in the 1970s.  Some of their best-known works are "Don", "Bairaag", "Saraswatichandra", "Qurbani", "Tridev" and "Safar".  They won the 1975 Filmfare Award for Best Music Director for "Kora Kagaz".', 'Kalyanji–Anandji discography | This is a discography of Bollywood composer duo Kalyanji Anandji, consisting of Kalyanji Virji Shah and Anandji Virji Shah.  They have composed music for over 200 films in their 46 year career.', 'Viju Shah | Viju Shah (born, Vijay Kalyanji Shah, 5 June 1959) is a music director of Hindi cinema.  He is the son of music director Kalyanji Virji Shah of composer duo Kalyanji Anandji.  He composed music for movies such as "Vishwatama" (1992), "Mohra" (1994), "Tere Mere Sapne" (1996) and "" (1997) for which got his second nomination for the Filmfare Award for Best Music Director and he won the 1998 Filmfare Award for Best Background Score for "".', 'Bombay the Hard Way: Guns, Cars and Sitars | Bombay the Hard Way: Guns, Cars and Sitars is a 1998 project of producer Dan the Automator in which music from 1970s funk Bollywood composer duo Kalyanji-Anandji (Kalyanji Virji Shah and Anandji Virji Shah) was taken and given a funky remix.  The album was eventually withdrawn (possibly due to copyright concerns).  A sequel, "Bombay 2: Electric Vindaloo", was also produced.', 'Kora Kagaz | Kora Kagaz (Hindi: कोरा कागज़; translation: Blank Paper) is a 1974 Hindi drama film produced by Sanath Kothari and directed by Anil Ganguly.  The film stars Vijay Anand, Jaya Bhaduri, A.K. Hangal, Achala Sachdev and Deven Verma.  The films music is by Kalyanji Anandji.  The famous title song "Mera Jeevan Kora Kagaz" was by Kishore Kumar.']
reasoning:
    The claim states that one half of the Kalyanji–Anandji duo, Kalyanji Virji Shah's brother, won the Filmfare Award for Best Music Director. From the passages and context, it is clear that the duo consists of Kalyanji Virji Shah and his brother Anandji Virji Shah. Both together won the Filmfare Award for Best Music Director in 1975 for the film "Kora Kagaz." Therefore, Anandji Virji Shah, who is Kalyanji's brother, indeed won the Filmfare Award for Best Music Director as part of this duo.
summary:
    Anandji Virji Shah, the brother of Kalyanji Virji Shah, won the Filmfare Award for Best Music Director in 1975 as part of the Kalyanji–Anandji duo.


##################################
3 other demos omitted for brevity.
##################################

\end{lstlisting}

GEPA Prompt generated by config GEPA+Merge:
\begin{lstlisting}[breakindent=0pt, breaklines=true, breakatwhitespace=true, frame=single, backgroundcolor=\color{blue!10}]
Given the input fields:

- `claim`: A factual statement or assertion that may be true or false.
- `context`: A brief explanation or summary relating to the claim that typically clarifies the claim’s accuracy.
- `passages`: A list of relevant textual evidence or knowledge snippets containing facts, descriptions, or biographies related to entities or concepts mentioned in the claim.

Your task is to produce a `summary` that meets the following requirements:

1. **Accurately reflect the relationship between the claim and the evidence**: Analyze the input `claim` carefully and check it against the `context` and all `passages`. Determine whether the claim is supported, partially supported, or contradicted by the evidence.

2. **Explicitly connect key entities and facts from the passages**: Your summary must mention the main entities and facts that directly confirm or refute parts of the claim. For example, highlight relevant names, works (films, albums, songs, books), attributes, dates, nicknames, or roles that clarify the claim's accuracy.

3. **Make substantive connections**: The summary should not merely restate the claim or conclusion. Instead, it should explicitly link the claim to detailed evidence from the passages—such as specific film titles, album names, song titles, dates, or relationships—to provide clear reasoning for the accuracy or inaccuracy of the claim.

4. **Include relevant disambiguations or clarifications**: Where applicable, clarify potential misunderstandings or common confusions highlighted in the passages (e.g., distinguishing similarly named films or songs, specifying which individual is referenced, or noting alternate titles and adaptations).

5. **Use concise, factual language**: The summary should be clear and succinct but include all relevant evidence necessary to understand how the claim is supported or refuted.

6. **Support future query generation**: As the summary will be used to generate queries to find evidence, it must contain explicit mentions of key entities, titles, and facts that can guide retrieval and reasoning.

**Overall, your approach should be to reason comprehensively over the claim, context, and each passage, synthesizing connected facts and evidence into a coherent and evidence-rich summary that clearly documents why the claim is true, false, or partially true.** Avoid vague or overly general statements; instead, ground your summary in the precise factual details present in the passages and context.
\end{lstlisting}
\end{instructbox}
\subsection{HoVer, Qwen3 8B}
\begin{instructbox}[\incode{HoVer Qwen3 8B create_query_hop2.predict}]
Base Prompt:
\begin{lstlisting}[breakindent=0pt, breaklines=true, breakatwhitespace=true, frame=single, backgroundcolor=\color{gray!10}]
Given the fields `claim`, `summary_1`, produce the fields `query`.
\end{lstlisting}

MIPROv2 Prompt:
\begin{lstlisting}[breakindent=0pt, breaklines=true, breakatwhitespace=true, frame=single, backgroundcolor=\color{green!10}]
In a high-stakes scenario where the accuracy of your query directly determines the validity of a complex claim, generate a precise query that builds on the provided claim and summary to uncover critical evidence. Use step-by-step reasoning to identify gaps in the initial summary, then craft a query that explicitly targets these gaps to retrieve additional supporting information. Ensure your query is unambiguous and directly addresses the claim's unresolved aspects.
Demos:
Example 1:
claim:
    Sandi Simcha DuBowski directed a documentary released in 2001. The horror film Zeitgeist was not also released in 2001.
summary_1:
    The claim is correct: Sandi Simcha DuBowski directed the 2001 documentary *Trembling Before G-d*, and there is no evidence in the passages that the horror film *Zeitgeist* was released in 2001.
reasoning:
    The claim is correct because Sandi Simcha DuBowski indeed directed the 2001 documentary *Trembling Before G-d*, as stated in summary_1. The assertion that *Zeitgeist* was not released in 2001 is supported by the absence of evidence in the provided summaries linking the film to that year.
query:
    Is there any evidence in the provided summaries that the horror film *Zeitgeist* was released in 2001?


##################################
3 other demos omitted for brevity.
##################################

\end{lstlisting}

GEPA Prompt generated by config GEPA:
\begin{lstlisting}[breakindent=0pt, breaklines=true, breakatwhitespace=true, frame=single, backgroundcolor=\color{blue!10}]
Given the fields `claim` and `summary_1`, generate a precise and focused query that identifies the specific evidence needed to verify or refute the claim.  

### Key Requirements:  
1. **Target Missing Evidence**: Identify and explicitly ask about the unverified or unconfirmed details in the summary (e.g., names, dates, locations, or connections) that are critical to the claim.  
   - Example: If the summary mentions "Massimo Giordano" as a potential counterpart but lacks birthplace details, the query should ask for evidence confirming their ties to Naples.  
2. **Correct Historical/Domain-Specific Anomalies**: Address discrepancies like anachronisms (e.g., Ali Qushji’s work in the 1960s vs. his actual 15th-century timeline) or misattributions (e.g., *Hayy ibn Yaqdhan* by Ibn Tufail, not Ali Qushji).  
3. **Link to Summary Context**: Ensure the query references the summary’s key points (e.g., "the claim states X, but the summary notes Y is unverified").  
4. **Use Specific Terms**: Include exact names, titles, or dates mentioned in the summary to avoid ambiguity (e.g., "Anaïs Nin," "Metropolitan City of Naples," "1948").  

### Example Strategy:  
If the summary states:  
- "The claim is partially supported. While [Fact A] is confirmed, [Fact B] lacks verification."  
- The query should ask:  
  *"Is there evidence confirming [Fact B], such as [specific name/date/location]?"*  

### Domain-Specific Notes:  
- Verify historical timelines (e.g., Ali Qushji died in 1474, not the 1960s).  
- Confirm authorships (e.g., *Hayy ibn Yaqdhan* is attributed to Ibn Tufail, not Ali Qushji).  
- Check for geographic or biographical details (e.g., birthplaces, cities, or professional connections).  

Ensure your query directly addresses the summary’s unverified claims and includes all critical terms from the summary to maximize evidence retrieval.
\end{lstlisting}
\end{instructbox}\begin{instructbox}[\incode{HoVer Qwen3 8B create_query_hop3.predict}]
Base Prompt:
\begin{lstlisting}[breakindent=0pt, breaklines=true, breakatwhitespace=true, frame=single, backgroundcolor=\color{gray!10}]
Given the fields `claim`, `summary_1`, `summary_2`, produce the fields `query`.
\end{lstlisting}

MIPROv2 Prompt:
\begin{lstlisting}[breakindent=0pt, breaklines=true, breakatwhitespace=true, frame=single, backgroundcolor=\color{green!10}]
You are a fact-checking assistant specializing in multi-hop reasoning and information synthesis. Given the fields `claim`, `summary_1`, and `summary_2`, generate a precise query that synthesizes the original claim with the two summaries to probe for deeper contextual relationships, resolving ambiguities and confirming supporting details through targeted evidence retrieval.
Demos:

\end{lstlisting}

GEPA Prompt generated by config GEPA:
\begin{lstlisting}[breakindent=0pt, breaklines=true, breakatwhitespace=true, frame=single, backgroundcolor=\color{blue!10}]
Given the fields `claim`, `summary_1`, and `summary_2`, produce the field `query` that:  
1. Explicitly asks whether the claim is supported by the provided summaries.  
2. Includes **specific evidence** from the summaries (e.g., names, titles, dates, or factual details) to guide retrieval of relevant documents.  
3. Focuses on **key disputed points** in the claim (e.g., incorrect attributions, missing evidence, or conflicting statements) to ensure the query targets the most relevant information.  

**Key considerations for the query:**  
- If the summaries mention specific works (e.g., *Planes, Trains and Automobiles*), include the title.  
- If the summaries reference named entities (e.g., John Candy, KLM, Richard Ford), include them.  
- If the summaries clarify contradictions (e.g., "not X but Y"), structure the query to highlight this contrast.  
- Avoid vague phrasing; instead of asking "Is the claim true?" use "Does the evidence in the summaries confirm [specific detail]?"  

**Example:**  
If the claim is about a film and the summaries mention *Planes, Trains and Automobiles* (1987), the query should include that title to retrieve relevant evidence.  
```
\end{lstlisting}
\end{instructbox}\begin{instructbox}[\incode{HoVer Qwen3 8B summarize1.predict}]
Base Prompt:
\begin{lstlisting}[breakindent=0pt, breaklines=true, breakatwhitespace=true, frame=single, backgroundcolor=\color{gray!10}]
Given the fields `claim`, `passages`, produce the fields `summary`.
\end{lstlisting}

MIPROv2 Prompt:
\begin{lstlisting}[breakindent=0pt, breaklines=true, breakatwhitespace=true, frame=single, backgroundcolor=\color{green!10}]
Given the fields `claim` and `passages`, generate a structured reasoning process that analyzes the claim step-by-step using the provided evidence, and produce a concise summary that distills the key findings and evaluates the claim's validity based on the evidence.
Demos:
Example 1:
claim:
    One half of the Kalyanji–Anandji duo, Kalyanji Virji Shah's brother, won the Filmfare Award for Best Music Director.
passages:
    ['Kalyanji Virji Shah | Kalyanji Virji Shah (30 June 1928 – 24 August 2000) was the "Kalyanji" of the Kalyanji-Anandji duo.  He and his brother Anandji Virji Shah have been famous Indian film musicians, and won the 1975 Filmfare Award for Best Music Director, for "Kora Kagaz".  He is a recipient of the civilian honour of Padma Shri (1992).', 'Kalyanji–Anandji | Kalyanji–Anandji are an Indian composer duo from Gujarat: Kalyanji Virji Shah (30 June 1928-03 November 2000) and his brother Anandji Virji Shah (born 02 March 1933).  The duo are known for their work on Hindi film soundtracks, particularly action potboilers in the 1970s.  Some of their best-known works are "Don", "Bairaag", "Saraswatichandra", "Qurbani", "Tridev" and "Safar".  They won the 1975 Filmfare Award for Best Music Director for "Kora Kagaz".', 'Anandji Virji Shah | Anandji Virji Shah is an Indian music director.  Together with his brother he formed the Kalyanji-Anandji duo, and won the 1975 Filmfare Award for Best Music Director, for "Kora Kagaz".  He is a recipient of the civilian honour of Padma Shri (1992).', 'Viju Shah | Viju Shah (born, Vijay Kalyanji Shah, 5 June 1959) is a music director of Hindi cinema.  He is the son of music director Kalyanji Virji Shah of composer duo Kalyanji Anandji.  He composed music for movies such as "Vishwatama" (1992), "Mohra" (1994), "Tere Mere Sapne" (1996) and "" (1997) for which got his second nomination for the Filmfare Award for Best Music Director and he won the 1998 Filmfare Award for Best Background Score for "".', 'Kalyanji–Anandji discography | This is a discography of Bollywood composer duo Kalyanji Anandji, consisting of Kalyanji Virji Shah and Anandji Virji Shah.  They have composed music for over 200 films in their 46 year career.', 'Bombay the Hard Way: Guns, Cars and Sitars | Bombay the Hard Way: Guns, Cars and Sitars is a 1998 project of producer Dan the Automator in which music from 1970s funk Bollywood composer duo Kalyanji-Anandji (Kalyanji Virji Shah and Anandji Virji Shah) was taken and given a funky remix.  The album was eventually withdrawn (possibly due to copyright concerns).  A sequel, "Bombay 2: Electric Vindaloo", was also produced.', 'Bairaag | Bairaag is a 1976 Bollywood film.  Produced by the duo Mushir-Riaz, it is directed by Asit Sen.  The film stars Dilip Kumar, who received a Filmfare nomination for Best Actor for playing a triple role.  "Bairaag" also stars Saira Banu, Leena Chandavarkar, Ruma Guha Thakurta, Prem Chopra, Helen, Sujit Kumar, Madan Puri, Paintal, Kader Khan, Asit Sen and Sachin.  Its music is by the duo Kalyanji Anandji, who received a Filmfare nomination for Best Music.  The film received a Filmfare Award for Best Sound for P. Harikishan.  According to one source, the film did not do well at the box office.']
reasoning:
    The claim states that Kalyanji Virji Shah's brother won the Filmfare Award for Best Music Director. The passages confirm that the Kalyanji–Anandji duo (Kalyanji Virji Shah and his brother Anandji Virji Shah) won the 1975 Filmfare Award for Best Music Director for "Kora Kagaz." Since the award was shared by both brothers, the claim is accurate as it identifies Anandji Virji Shah (Kalyanji's brother) as the winner.
summary:
    Anandji Virji Shah, the brother of Kalyanji Virji Shah, won the 1975 Filmfare Award for Best Music Director as part of the Kalyanji–Anandji duo for their work on "Kora Kagaz."


##################################
3 other demos omitted for brevity.
##################################

\end{lstlisting}

GEPA Prompt generated by config GEPA:
\begin{lstlisting}[breakindent=0pt, breaklines=true, breakatwhitespace=true, frame=single, backgroundcolor=\color{blue!10}]
Given the fields `claim` and `passages`, produce the fields `summary` by:  
1. **Analyzing the claim** to identify key elements (e.g., names, locations, relationships, specific terms).  
2. **Cross-referencing the passages** to locate explicit or implicit evidence that supports, contradicts, or partially addresses the claim.  
3. **Explicitly connecting the claim to the evidence** in the passages, even if the evidence is not directly stated in the claim (e.g., linking "Roger Yuan" to "Shanghai Noon" in Example 1).  
4. **Highlighting domain-specific details** (e.g., "Sentinelese language is unclassifiable" in Example 3) that are critical for accuracy.  
5. **Stating whether the claim is supported, unsupported, or partially supported**, and explaining the reasoning with clear references to the passages.  
6. **Prioritizing key evidence** that would be relevant for further queries (e.g., "Marlborough" in Example 2) to ensure retrievability.  

**Additional Guidance:**  
- If the claim contains unverified details (e.g., "Roger Yuan" in Example 1), explicitly note this in the summary.  
- For niche or ambiguous information (e.g., "uncontacted people" in Example 3), reference specific passages that clarify the context.  
- Ensure summaries are concise but include all critical evidence needed to validate or refute the claim.
\end{lstlisting}
\end{instructbox}\begin{instructbox}[\incode{HoVer Qwen3 8B summarize2.predict}]
Base Prompt:
\begin{lstlisting}[breakindent=0pt, breaklines=true, breakatwhitespace=true, frame=single, backgroundcolor=\color{gray!10}]
Given the fields `claim`, `context`, `passages`, produce the fields `summary`.
\end{lstlisting}

MIPROv2 Prompt:
\begin{lstlisting}[breakindent=0pt, breaklines=true, breakatwhitespace=true, frame=single, backgroundcolor=\color{green!10}]
Given the fields `claim`, `context`, and `passages`, perform multi-hop reasoning to generate a structured summary that validates or refutes the claim. First, analyze the claim and contextual information to identify key relationships. Next, evaluate evidence from the retrieved passages to build a coherent narrative. Construct a logical reasoning chain connecting the claim to supporting or contradictory evidence. Finally, produce a concise summary that explicitly confirms or contradicts the claim based on your analysis, citing relevant evidence from the context and passages.
Demos:
Example 1:
claim:
    The recorded history of the state, for which Baptist George Ruby was a prominent black Republican leader in the Reconstruction-era, officially started in 1519 in the same state that holds the city Augustus Chapman Allen used his inheritance to fund the founding of.
context:
    The claim is not supported by the passages. While George Ruby was associated with Texas during Reconstruction and Augustus Chapman Allen founded Houston in Texas, there is no evidence in the passages that Texas's recorded history began in 1519.
passages:
    ["Merchants and Manufacturers Building | The One Main Building, formerly the Merchants and Manufacturers Building (commonly referred to as the M&M Building), is a building on the campus of the University of Houston–Downtown.  The building is recognized as part of the National Register of Historic Places, is a Recorded Texas Historic Landmark, and considered a Contributing Building in Downtown Houston's Main Street/Market Square Historic District.  The building was built above Allen's Landing—an area where Houston's founders John Kirby Allen and Augustus Chapman Allen originally settled.", 'Augustus Chapman Allen | Augustus Chapman Allen (July 4, 1806 – January 11, 1864), along with his younger brother, John Kirby Allen, founded the City of Houston in the U.S. state of Texas.  He was born on July 4, 1806, in Canaseraga Village, New York (the present day hamlet of Sullivan in the Town of Sullivan, New York), to Sarah (Chapman) and Roland Allen.', "History of Texas | The recorded History of Texas begins with the arrival of the first Spanish conquistadors in the region of North America now known as Texas in 1519, who found the region populated by numerous Native American /Indian tribes.  
    Native Americans' ancestors had been there for more than 10,000 years as evidenced by the discovery of the remains of prehistoric Leanderthal Lady.  
    During the period of recorded history from A.D. 1519 to 1848, all or parts of Texas were claimed by five countries: France, Spain, Mexico, the Republic of Texas, the United States of America—as well as the secessionist Confederate States of America in 1861–1865, during the American Civil War.", 'Charlotte Baldwin Allen | Charlotte Baldwin Allen (July 14, 1805 – August 3, 1895) is known in Texan history as the "mother of Houston".  She was the wife of Augustus Chapman Allen, who used her inheritance to finance the founding of this city.', 'Allen Parkway | Allen Parkway is an arterial road west of Downtown Houston, Texas.  It has a distance of approximately 2.3 miles (3.7\xa0km), running from Interstate 45 west to Shepherd Drive, where it becomes Kirby Drive.  Originally known as Buffalo Parkway, it was later named after John Kirby Allen and Augustus Chapman Allen, the founders of Houston.', "Allen's Landing | Allen's Landing is the birthplace of the city of Houston—the largest city in the U.S. state of Texas.  In August 1836, just months after the Republic of Texas won its independence from Mexico, two brothers (and real estate developers) from New York—John Kirby Allen and Augustus Chapman Allen—purchased 6,642\xa0acres\xa0(27\xa0km²) in the area and settled there on the banks of Buffalo Bayou.  Allen's Landing is located south of the University of Houston–Downtown Commerce Street Building, at the intersection of Commerce Street and Main Street.", "Keith Allen (American football) | Keith D. Allen (born July 12, 1974) is an American football coach.  He served as the head football coach at Southwest Baptist University from 2007 to 2012, compiling a record of 27–39.  His 27 wins are the most in the history of Southwest Baptist's football program, which began play in 1983."]
reasoning:
    The claim is not supported by the passages. While the passages confirm that Texas's recorded history began in 1519 (passage 3) and that Augustus Chapman Allen founded Houston in Texas (passage 2), they do not explicitly link the start of Texas's recorded history to the same state where Augustus Chapman Allen used his inheritance to fund the founding of a city. The claim conflates these two facts without direct evidence in the passages that the 1519 start of recorded history is tied to the state where Allen's inheritance was used for founding Houston.
summary:
    The claim is not supported by the provided passages, as they confirm Texas's recorded history began in 1519 and Augustus Chapman Allen founded Houston in Texas but do not explicitly connect the two facts as stated in the claim.


##################################
3 other demos omitted for brevity.
##################################

\end{lstlisting}

GEPA Prompt generated by config GEPA:
\begin{lstlisting}[breakindent=0pt, breaklines=true, breakatwhitespace=true, frame=single, backgroundcolor=\color{blue!10}]
Given the fields `claim`, `context`, and `passages`, produce the field `summary` that:  
1. **Evaluates the claim's validity** by explicitly connecting it to **specific evidence** from the `passages` (e.g., citing exact names, dates, or terms that support/refute the claim).  
2. **Highlights factual inaccuracies** in the claim, such as incorrect origins, dates, or relationships, and links them to **direct contradictions** in the `passages` or `context`.  
3. **Mentions key terms or entities** critical to the claim (e.g., "Sly Stone," "Large Münsterländer," or "Hit Parader's 2006 ranking") to ensure summaries are actionable for query generation.  
4. **Balances accuracy** by acknowledging both correct and incorrect aspects of the claim, as seen in the examples (e.g., validating "The Wild Bunch" compilation while correcting the magazine's nationality and ranking year).  
5. **Avoids vague statements** like "no evidence" and instead specifies **which parts of the passages** are relevant (e.g., "Passage 6 confirms Hit Parader was American").  

**Generalizable Strategy**:  
- Always reference **specific passages** (e.g., "Passage 1 states...") to ground claims in the text.  
- Use **precise terminology** (e.g., "German Shorthaired Pointer" instead of "cousins") to align with the data.  
- Prioritize **clarity and specificity** to ensure summaries can guide further evidence retrieval.
\end{lstlisting}
\end{instructbox}
\subsection{PUPA, GPT-4.1 Mini}
\begin{instructbox}[\incode{PUPA GPT-4.1 Mini craft_redacted_request.predict}]
Base Prompt:
\begin{lstlisting}[breakindent=0pt, breaklines=true, breakatwhitespace=true, frame=single, backgroundcolor=\color{gray!10}]
Given a private user query, create a privacy-preserving request for a powerful external LLM.
The LLM may assist without learning private information about the user.
\end{lstlisting}

MIPROv2 Prompt:
\begin{lstlisting}[breakindent=0pt, breaklines=true, breakatwhitespace=true, frame=single, backgroundcolor=\color{green!10}]
You will be given a private user query. Your task is to analyze the query carefully and think step-by-step to understand the user's intent while identifying and removing or abstracting any sensitive, proprietary, or personally identifiable information. Then, generate a clear, privacy-preserving prompt suitable for an external large language model (LLM) that requests useful, relevant guidance or information without exposing private details. Additionally, provide a detailed chain-of-thought explanation describing your reasoning process for how you transformed the original query into the redacted, generalized LLM request. Ensure the resulting prompt is professional, contextually appropriate, and maintains user privacy while enabling the LLM to assist effectively.
Demos:
Example 1:
user_query:
    Imagine you are a stock researcher and advisor, I will give you a table about stocks and I want you to give me advice on what to invest next. This is a project I am working on an I need help with predictions I am making. Here is the table: Stock Price Dividends 5D 5D % 1M 1M % YTD YTD % 1Y 1Y % 5Y 5Y % IBM $146,62 $6,64 $3,81 2,66 $2,63 1,83 $5,07 3,58 $16,31 12,51 $0,03 0,02 MCD $284,90 $6,08 $5,07 1,81 -$7,20 -2,46 $20,57 7,78 $29,24 11,44 $123,26 76,26 JNJ $163,98 $4,76 -$3,00 -1,80 -$3,90 -2,32 -$14,21 -7,97 $0,98 0,60 $29,66 22,08 MSFT $324,89 $2,72 $1,25 0,39 -$9,88 -2,95 $85,31 35,61 $59,66 22,49 $217,33 202,05 WMT $158,36 $2,28 $1,39 0,89 -$1,85 -1,15 $14,75 10,27 $25,47 19,16 $62,29 64,84 KO $60,50 $1,84 $0,12 0,20 -$1,59 -2,55 -$2,45 -3,89 -$2,23 -3,55 $15,13 33,35 PFE $36,07 $1,64 -$1,06 -2,84 -$0,67 -1,81 -$15,19 -29,64 -$10,16 -21,98 -$5,43 -13,09 AAPL $179,96 $0,96 $3,45 1,95 -$16,57 -8,43 $54,89 43,89 $18,58 11,51 $125,04 227,65 NVDA $468,23 $0,16 -$3,71 -0,79 $7,05 1,53 $325,08 227,09 $310,22 196,33 $401,52 601,89 CRM $212,10 $0,00 $2,75 1,31 -$9,86 -4,44 $77,32 57,37 $51,89 32,39 $59,11 38,64 PYPL $61,16 $0,00 $1,06 1,76 -$13,82 -18,43 -$13,43 -18,00 -$31,51 -34,00 -$29,46 -32,51 BIDU $137,47 $0,00 $10,24 8,05 -$15,63 -10,21 $18,34 15,39 -$9,85 -6,69 -$88,25 -39,10
reasoning:
    The user wants investment advice based on a provided table of stock data including price, dividends, and percentage changes over various time intervals (5 days, 1 month, year-to-date, 1 year, and 5 years). The primary request is for recommendations on what to invest in next and help with making predictions.

To preserve privacy, I should not include the user's raw data or specific details about individual stocks directly in the prompt to the external LLM. Instead, I should request general advice on how to analyse stock data tables that include price, dividend, and percentage changes over multiple time periods. Also, I can ask for advice on what indicators to weigh more heavily for investment decisions and for predictive strategies based on such data structures.

This approach ensures that the external LLM does not receive any private or identifying user data, only a generic request. I will ask the LLM to provide guidance on interpreting such stock performance indicators and suggestions for predictive analytics or investment selection criteria.
llm_request:
    Please provide general guidance for analyzing stock investment opportunities based on a data table format that includes stock prices, dividend payouts, and percentage changes over various time intervals such as 5 days, 1 month, year-to-date, 1 year, and 5 years. Outline key indicators that investors typically weigh when making decisions, how to interpret short-term versus long-term performance trends, and suggest strategies or metrics to use when making predictions or recommendations about which stocks to invest in next. The request is for generic advice without any specific stock names or values.


##################################
3 other demos omitted for brevity.
##################################

\end{lstlisting}

GEPA Prompt generated by config GEPA+Merge:
\begin{lstlisting}[breakindent=0pt, breaklines=true, breakatwhitespace=true, frame=single, backgroundcolor=\color{blue!10}]
New Instruction for Privacy-Preserving Query Reformulation Assistant (Enhanced):

Task Overview:  
You receive a user query that may contain sensitive, private, or personally identifiable information (PII). Your objective is to reformulate this query into a generalized, privacy-preserving prompt suitable for sending to an external large language model (LLM). The reformulated prompt must enable the external LLM to fulfill the user's original intent effectively while rigorously protecting all user privacy by abstracting, omitting, or generalizing any sensitive details.

Input Format:  
- A single user query string.  
- The query might contain identifiable names (people, organizations), specific geographic locations, exact dates or timeframes, proprietary or internal terminology (e.g., system names, product codes), client or customer names, URLs, or any other sensitive or private information.

Output Format:  
- Part (a) Reasoning:  
   Provide a concise paragraph that explains:  
   * How you identified sensitive or private information in the input  
   * What strategies you applied to protect privacy (generalization, omission, abstraction, replacement with placeholders)  
   * How the reformulated prompt preserves the original intent and task requirements without risking data leakage  
- Part (b) LLM Request:  
   A concise, carefully constructed privacy-safe prompt that:  
   * Removes or anonymizes all PII and proprietary/internal details  
   * Abstracts locations, names, dates, and technical terms as needed  
   * Produces a clear and contextually rich instruction for the LLM to generate a relevant and informative response aligned with the user's original task

Detailed Domain-Specific Guidance and Best Practices:

1. Identification and Treatment of Sensitive Data:
   - All user-specific or personal names (individual or organizational) must be removed or replaced with generic role descriptors (e.g., “a business contact,” “a client,” “a notable individual”). Never lightly obscure or partially redact; full abstraction is required.
   - All geographic mentions must be abstracted unless the location is publicly known, essential to the task, and can be generalized (e.g., “a region known for eco-tourism” instead of naming a country or city explicitly).
   - Exact dates or durations must never be retained; instead, use relative or approximate temporal references (e.g., “recently,” “over the past year”).
   - Internal or proprietary terms — including system names, product codes, subscription types, and technical jargon — must be generalized or replaced with neutral descriptors to avoid leakage of intellectual property or sensitive operational details.
   - Avoid direct quotes or verbatim inclusion of user-supplied texts unless obfuscated by generalization.

2. Task Understanding and Reformulation:
   - Identify the functional intent of the query: Is it creative writing, translation, summarization, professional communication drafting, technical explanation, or other?
   - Preserve the thematic and informational core of the query (e.g., request for educational quality analysis, technical translation of a passage, biographical summary).
   - Do not reproduce the original input verbatim; instead, frame the LLM prompt around the essential thematic elements extracted from the input.
   - For queries regarding individuals, avoid direct reference to real personal information unless the name is widely public and essential; even then, use a generic or hypothetical framing for the individual profile.

3. Strategies for High-Quality and Privacy-Preserving Prompts:
   - Strike a balance between sufficient contextual detail and privacy abstraction to maintain prompt clarity and relevance.
   - Use neutral, context-aware formulations that clearly instruct the LLM on the content and style expected.
   - Avoid vague or overly generic prompts that could result in less useful or lower-quality responses.
   - When system or proprietary content is mentioned, instruct the LLM to generalize specific terms and maintain the technical meaning without revealing sensitive info.
   - When a direct translation is requested on specialized text, specify to replace or abstract internal nomenclature.

4. Explanation Requirements:
   - The reasoning must transparently explain how privacy risks were identified (e.g., presence of names, locations, dates, proprietary terms).
   - It must describe the abstraction or omission methods applied (e.g., replacing “Jonah Van Beijnen” with “a notable individual,” substituting “Makau” with “a specific region,” or “Yoda” with “a system name”).
   - Clarify how the essential task and user intent were preserved despite these generalizations (e.g., focusing on educational quality, technical translation, biographical summary).
   - This explanation justifies your design choices and demonstrates adherence to privacy-preserving principles.

Common Pitfalls to Avoid:
- Do not merely lightly obscure or partially redact sensitive details; full anonymization or abstraction is required.
- Do not repeat any user-supplied PII or proprietary content verbatim.
- Avoid including URLs, exact dates, or direct quotes without modification.
- Do not leave ambiguity that could degrade the quality or contextual clarity of the reformulated prompt.
- Do not include any real personal or organizational names unless they are public figures and the query requires it, then use generic descriptors instead.

Example Summary of Effective Approach (Informed by Prior Examples):
- For geographic queries: replace exact place names with general regions and provide a brief contextual descriptor.
- For technical texts containing system names or subscription types: instruct the LLM to translate or process the text while replacing or abstracting proprietary system identifiers.
- For biographical summaries about specific individuals: remove the real name and request a generic, well-structured four-paragraph summary about “a notable individual,” preserving the overall intent without leaking PII.

Summary:
Your reformulations must ensure zero exposure of any PII or private/proprietary content while retaining enough thematic and functional clarity for the external LLM to produce high-quality, relevant outputs. This requires thorough analysis of the user's query, rigorous application of privacy-preservation strategies, and explicit reasoning explanations that document your approach and choices.
\end{lstlisting}
\end{instructbox}\begin{instructbox}[\incode{PUPA GPT-4.1 Mini respond_to_query}]
Base Prompt:
\begin{lstlisting}[breakindent=0pt, breaklines=true, breakatwhitespace=true, frame=single, backgroundcolor=\color{gray!10}]
Respond to a user query.
For inspiration, we found a potentially related request to a powerful external LLM and its response.
\end{lstlisting}

MIPROv2 Prompt:
\begin{lstlisting}[breakindent=0pt, breaklines=true, breakatwhitespace=true, frame=single, backgroundcolor=\color{green!10}]
Given a privacy-preserving redacted prompt sent to an external untrusted LLM and its response, along with the original user query containing potentially sensitive information, generate a final, coherent, and contextually appropriate user-facing response. Your output should integrate the anonymized LLM output with the original query’s intent, ensuring privacy is maintained while providing a complete and polished answer that directly addresses the user's request.
Demos:
Example 1:
related_llm_request:
    Please provide a detailed explanation of David Easton's system analysis in political science, outlining the key components of his systems theory approach and its significance in the study of political systems.
related_llm_response:
    Certainly! David Easton's system analysis is a foundational framework in political science that applies systems theory to the study of political systems. His approach provides a structured way to understand how political systems function, maintain stability, and respond to changes in their environment. Below is a detailed explanation of Easton's system analysis, including its key components and significance.

---

### Overview of David Easton's System Analysis

David Easton (1917–2014) was a prominent political scientist who sought to develop a general theory of political systems by borrowing concepts from systems theory, which originated in biology and engineering. His goal was to create a scientific and analytical framework to study politics as an integrated and dynamic system rather than as isolated events or institutions.

Easton's system analysis views a political system as a set of interactions abstracted from the totality of social behavior, through which authoritative allocation of values (decisions) is made and implemented. In other words, politics is about how societies decide who gets what, when, and how, and Easton's model helps explain how these decisions are made and sustained.

---

### Key Components of Easton's Systems Theory Approach

Easton's model conceptualizes the political system as a **black box** that processes inputs and produces outputs, embedded within a larger environment. The key components are:

#### 1. **Inputs**
Inputs are demands and supports that come from the environment and society into the political system.

- **Demands:** These are claims or requests made by individuals or groups for the political system to allocate values or resources in a particular way (e.g., demands for law enforcement, education, welfare).
- **Supports:** These are actions or attitudes that indicate approval, allegiance, or compliance with the political system (e.g., voting, paying taxes, obeying laws).

Inputs represent the pressures and resources that the political system must process.

#### 2. **Political System (The Black Box)**
The political system is the core of Easton's model. It receives inputs, processes them, and produces outputs.

- It includes institutions, organizations, and processes involved in authoritative decision-making (e.g., legislatures, executives, courts).
- The system transforms demands and supports into decisions and policies.
- The system must maintain its boundaries and adapt to changes to survive.

#### 3. **Outputs**
Outputs are the decisions, policies, and actions that the political system produces in response to inputs.

- These include laws, regulations, public policies, and authoritative allocations of values.
- Outputs affect the environment and society, influencing future inputs.

#### 4. **Feedback**
Feedback is the process by which outputs influence future inputs.

- Positive feedback can reinforce the system’s stability.
- Negative feedback can signal problems or dissatisfaction, prompting adjustments.
- Feedback helps the system adapt and maintain equilibrium.

#### 5. **Environment**
The political system exists within a broader environment that includes social, economic, cultural, and international factors.

- The environment provides inputs and is affected by outputs.
- The system must interact with and adapt to environmental changes to survive.

---

### Diagrammatic Representation (Simplified)

```
Environment
   ↑          ↓
[Demands & Supports] → [Political System] → [Decisions & Policies]
          ↑                                ↓
          ←----------- Feedback -----------
```

---

### Significance of Easton's System Analysis in Political Science

1. **Holistic Understanding of Politics:**
   Easton's model encourages viewing politics as an interconnected system rather than isolated events or institutions. It highlights the dynamic interactions between society and political authority.

2. **Focus on Process and Stability:**
   The model emphasizes how political systems maintain stability and order through continuous processing of inputs and outputs, and through feedback mechanisms.

3. **Analytical Clarity:**
   By abstracting political activity into inputs, system processing, outputs, and feedback, Easton provided a clear framework for empirical research and comparative analysis.

4. **Foundation for Further Theories:**
   Easton's systems theory influenced later developments in political science, including structural-functionalism, policy analysis, and the study of political development and change.

5. **Adaptability:**
   The model can be applied to various political systems (democracies, authoritarian regimes, etc.) and scales (local, national, international), making it versatile.

6. **Emphasis on Environment:**
   Recognizing the environment’s role underscores the importance of external factors (economic conditions, social movements, international pressures) in shaping political outcomes.

---

### Criticisms and Limitations

- **Overly Abstract:** Some critics argue that Easton's model is too abstract and general, lacking specificity about power relations, conflict, and individual agency.
- **Neglect of Conflict:** The model tends to emphasize stability and equilibrium, potentially underplaying political conflict and change.
- **Limited Normative Guidance:** It is primarily descriptive and analytical, offering little normative guidance on what political systems ought to do.

---

### Summary

David Easton's system analysis
user_query:
    System analysis by David Easton
response:
    David Easton's system analysis is a foundational framework in political science that conceptualizes the political system as a dynamic and interactive system. His approach uses systems theory to explain how political systems function, maintain stability, and adapt to changes in their environment.

At its core, Easton's model views the political system as a "black box" that processes inputs from society and the environment, transforms them into outputs, and uses feedback to adjust and sustain itself. The key components of Easton's system analysis include:

1. **Inputs:** These are the demands and supports that society directs towards the political system. Demands are requests or claims for resource allocation or policy change, while supports are actions or attitudes such as compliance, legitimacy, or loyalty that sustain the system.

2. **Political System:** This is the set of institutions and processes that processes the inputs. It interprets demands and supports to make authoritative decisions, including the formulation and enforcement of policies.

3. **Outputs:** These are the decisions, policies, and actions that the political system produces in response to inputs. They affect society and its environment, shaping future interactions.

4. **Feedback:** Outputs generate responses that feed back into the system as new inputs, signaling whether the outputs were effective or require modification. This feedback enables the system to adapt to internal and external changes.

5. **Environment:** The system operates within a broader social, economic, cultural, and international context that influences and is influenced by the political system.

Easton's theory is significant because it offers a holistic, process-oriented lens to analyze politics, emphasizing interaction between society and state institutions. It provides a framework to understand political stability, adaptation, and change by focusing on the continuous circulation of inputs, processing, outputs, and feedback. The model has been influential in shaping subsequent political theories and comparative political analysis.

However, critics note that Easton's model can be overly abstract, sometimes downplaying conflict, power struggles, and individual agency within political processes.

In summary, David Easton's system analysis offers a powerful conceptual tool to study political systems as dynamic entities that manage societal demands and supports through authoritative decision-making to maintain order and adapt over time.


##################################
3 other demos omitted for brevity.
##################################

\end{lstlisting}

GEPA Prompt generated by config GEPA+Merge:
\begin{lstlisting}[breakindent=0pt, breaklines=true, breakatwhitespace=true, frame=single, backgroundcolor=\color{blue!10}]
Respond to a user query.
For inspiration, we found a potentially related request to a powerful external LLM and its response.
\end{lstlisting}
\end{instructbox}
\subsection{PUPA, Qwen3 8B}
\begin{instructbox}[\incode{PUPA Qwen3 8B craft_redacted_request.predict}]
Base Prompt:
\begin{lstlisting}[breakindent=0pt, breaklines=true, breakatwhitespace=true, frame=single, backgroundcolor=\color{gray!10}]
Given a private user query, create a privacy-preserving request for a powerful external LLM.
The LLM may assist without learning private information about the user.
\end{lstlisting}

MIPROv2 Prompt:
\begin{lstlisting}[breakindent=0pt, breaklines=true, breakatwhitespace=true, frame=single, backgroundcolor=\color{green!10}]
Given a private user query, generate a structured, non-sensitive LLM request that captures the user's intent while omitting personal details. Include a step-by-step reasoning process explaining how the redaction was achieved, ensuring the request is actionable and privacy-preserving. Format the reasoning as a coherent thought process and the request as a clear, specific task for the LLM to execute.
Demos:
Example 1:
user_query:
    Привет напиши команду для отправки snmpwalk на ip 10.2.35.2 с данными по доступу initi_user/d2J3BgtYdB2/d2J3BgtYdB2 (MD5) DES
reasoning:
    The user is requesting an SNMP walk command to retrieve data from an SNMPv3-enabled device at IP 10.2.35.2. The provided credentials include a username (`initi_user`), an authentication password (`d2J3BgtYdB2`), and an encryption password (`d2J3BgtYdB2`). The authentication protocol is MD5, and the encryption protocol is DES. The command must include SNMPv3 options to specify these parameters correctly.
llm_request:
    Generate a command-line instruction for executing `snmpwalk` with SNMPv3 authentication and encryption using the provided credentials, ensuring privacy by omitting sensitive details in the response.


##################################
3 other demos omitted for brevity.
##################################

\end{lstlisting}

GEPA Prompt generated by config GEPA:
\begin{lstlisting}[breakindent=0pt, breaklines=true, breakatwhitespace=true, frame=single, backgroundcolor=\color{blue!10}]
### New Instruction for the Assistant:  

**Task Overview:**  
Rephrase the user's query into a privacy-preserving, actionable request for an external LLM, ensuring no personally identifiable information (PII), proprietary data, or sensitive details are exposed. Maintain the original intent, clarity, and specificity of the query while adhering to domain-specific strategies for generalized information.  

**Key Requirements:**  
1. **Privacy Preservation**:  
   - **Generalize Specifics**: Replace specific locations (e.g., "Andhra Pradesh" → "a major Indian city"), numbers (e.g., "18 volunteers" → "a group of participants"), or proprietary data (e.g., "van Cleef production sites" → "[Brand Name]'s production sites based on publicly available information").  
   - **Avoid PII/Proprietary Data**: Remove names, addresses, dates, internal processes, or brand-specific secrets. Use placeholders like [Brand Name], [Country], or [Region] for brand/company references.  
   - **Publicly Available Focus**: For brand-related queries, emphasize publicly accessible data (e.g., "locations of [Brand Name]'s production sites" instead of "van Cleef's secret factories").  

2. **Quality of Request**:  
   - **Clarity and Actionability**: Ensure the rephrased query is precise, avoids ambiguity, and specifies constraints (e.g., "concise, platform-friendly hashtags" for creative tasks).  
   - **Domain-Specific Precision**: For scientific/technical queries, retain key terms (e.g., "CAPM formula," "pharmacokinetic parameters") while anonymizing study details (e.g., "18 volunteers" → "a group of participants").  
   - **Avoid Overly Broad Requests**: Specify constraints like keyword focus, tone, or format (e.g., "health and fitness keywords" instead of "general hashtags").  

3. **Domain-Specific Strategies**:  
   - **Geographic Queries**: Replace regions/cities with generalized terms (e.g., "major city in [Country]" or "well-known [Region]").  
   - **Brand/Company Queries**: Reference brands by name but prioritize publicly available data (e.g., "locations of [Brand Name]'s production sites").  
   - **Creative Tasks**: Use examples or constraints (e.g., "wellness, strength, motivation keywords" for hashtags; "platform guidelines" for social media content).  

4. **Avoid Leakage**:  
   - **No Inferred Details**: Do not assume user-specific preferences, contexts, or private data (e.g., "Tell me the location of van Cleef production sites" → "Provide locations of [Brand Name]'s production sites based on publicly available information").  
   - **Anonymize Numerical Data**: Replace specific counts, dosages, or metrics with generalized terms (e.g., "40 mg OA" → "a dose of OA").  

**Example Workflow:**  
- **User Query**: "Calculate D’Amico’s cost of equity using CAPM with rRF = 4.67%, β = 1.56, and market risk premium = 6.63%."  
- **Rephrased Request**: "Calculate the cost of equity for a company using the CAPM formula, given a risk-free rate of 4.67%, a beta of 1.56, and a market risk premium of 6.63%. Show step-by-step calculations and round the final result to two decimal places."  

**Niche Domain-Specific Guidance:**  
- **Scientific/Pharmacokinetic Queries**: Generalize study details (e.g., "18 healthy male volunteers" → "a group of participants"), retain key parameters (e.g., AUC, peak concentrations), and avoid proprietary methodologies.  
- **Creative/Translation Tasks**: Use placeholders for song titles/artists (e.g., "Mahmood's song [Title]") and specify tone/rhythm requirements for translations.  
- **Financial Calculations**: Retain formulas and numerical values but avoid linking to specific companies or proprietary models.  

**Final Output Format:**  
- **reasoning**: Briefly explain the rephrasing choices (e.g., "Generalized '18 volunteers' to 'a group of participants' to preserve privacy").  
- **llm_request**: Provide the anonymized, actionable request for the LLM.
\end{lstlisting}
\end{instructbox}\begin{instructbox}[\incode{PUPA Qwen3 8B respond_to_query}]
Base Prompt:
\begin{lstlisting}[breakindent=0pt, breaklines=true, breakatwhitespace=true, frame=single, backgroundcolor=\color{gray!10}]
Respond to a user query.
For inspiration, we found a potentially related request to a powerful external LLM and its response.
\end{lstlisting}

MIPROv2 Prompt:
\begin{lstlisting}[breakindent=0pt, breaklines=true, breakatwhitespace=true, frame=single, backgroundcolor=\color{green!10}]
Synthesize the external LLM's response with the user's specific query to generate a final privacy-preserving answer. Replace placeholders with actual values while omitting sensitive information, maintaining a professional tone, and ensuring contextual accuracy. Combine the technical solution from the LLM with the user's unique details to produce a clear, anonymized response that adheres to privacy constraints.
Demos:
Example 1:
related_llm_request:
    Create a PESTLE analysis framework for GE Healthcare in the US, outlining key issues in each category (Political, Economic, Social, Technological, Environmental, Legal) and their potential impact on the organization. Prioritize these issues based on urgency and significance without referencing confidential data. Structure the output as a matrix with clear criteria for evaluation.
related_llm_response:
    Certainly! Below is a **PESTLE analysis framework** tailored for **GE Healthcare in the US**, structured as a matrix. It outlines key issues in each category, evaluates their **urgency** and **significance**, and highlights their **potential impact** on the organization.

---

### PESTLE Analysis Matrix for GE Healthcare (US)

| **Category**   | **Key Issues**                                                                 | **Urgency** (High/Med/Low) | **Significance** (High/Med/Low) | **Potential Impact on GE Healthcare**                                                                                  |
|----------------|--------------------------------------------------------------
----
--------------|----------------------------|---------------------------------
|-------------------------------------------------------------------------
-----------------------------------------------|
| **Political**  | 1. Healthcare policy changes (e.g., Medicare/Medicaid reimbursement reforms)    | High                       | High                            | Directly affects revenue streams and product/service demand; requires adaptation to new reimbursement models.          |
|                | 2. Trade policies and tariffs on medical equipment imports/exports              | Medium                     | Medium                          | Could increase costs of components or affect supply chain efficiency.                                                  |
|                | 3. Government funding for healthcare innovation and research                    | Medium                     | High                            | Influences R&D opportunities and partnerships; potential for grants and subsidies.                                     |
| **Economic**   | 1. US healthcare spending trends and budget constraints                         | High                       | High                            | Impacts purchasing power of hospitals and clinics; affects sales volume and pricing strategies.                        |
|                | 2. Inflation and rising operational costs                                      | High                       | Medium                          | Increases production and operational expenses; may pressure profit margins.                                            |
|                | 3. Economic cycles affecting capital investment in healthcare infrastructure    | Medium                     | Medium                          | Influences demand for high-cost medical equipment and technology upgrades.                                             |
| **Social**     | 1. Aging population increasing demand for healthcare services and devices       | High                       | High                            | Drives demand for diagnostic imaging, monitoring, and therapeutic devices; growth opportunity.                         |
|                | 2. Growing patient awareness and demand for personalized medicine               | Medium                     | Medium                          | Necessitates innovation in tailored healthcare solutions and data-driven technologies.                                 |
|                | 3. Workforce shortages in healthcare professionals                              | High                       | Medium                          | May affect adoption and training for new technologies; impacts customer support and service delivery.                  |
| **Technological** | 1. Rapid advancements in AI and machine learning for diagnostics and imaging | High                       | High                            | Critical for maintaining competitive edge; requires continuous investment in R&D and product development.             |
|                | 2. Integration of IoT and connected devices in healthcare                       | Medium                     | High                            | Enhances product offerings but demands robust cybersecurity and interoperability solutions.                            |
|                | 3. Data privacy and cybersecurity technologies                                 | High                       | High                            | Essential to protect patient data and comply with regulations; failure risks legal penalties and reputational damage.  |
| **Environmental** | 1. Regulatory pressure to reduce carbon footprint and waste in manufacturing | Medium                     | Medium                          | Necessitates sustainable practices; potential cost implications but improves brand reputation.                         |
|                | 2. Energy efficiency in product design and operations                          | Medium                     | Medium                          | Can reduce operational costs and appeal to environmentally conscious customers.                                        |
|                | 3. E-waste management and recycling programs                                  | Low                        | Medium                          | Compliance with environmental standards; growing importance for corporate social responsibility.                       |
| **Legal**      | 1. Compliance with FDA regulations and medical device approvals                 | High                       | High                            | Critical for market access; delays or failures can halt product launches and cause financial losses.                   |
|                | 2. Intellectual property rights and patent protections                         | Medium                     | High                            | Protects innovation and competitive advantage; infringement risks costly litigation.                                   |
|                | 3. Liability and product safety regulations                                   | High                       | High                            | Ensures patient safety; non-compliance can lead to lawsuits, recalls, and reputational damage.                         |

---

### Criteria for Evaluation:

- **Urgency**: How immediate the issue requires attention (High = immediate/ongoing; Medium = within 1-2 years; Low = longer-term).
- **Significance**: The potential magnitude of impact on GE Healthcare’s operations, profitability, and strategic positioning (High = critical; Medium = moderate; Low = minor).
- **Potential Impact**: Describes how the issue could affect GE Healthcare’s business, including risks and opportunities.

---

This matrix helps prioritize strategic focus areas for GE Healthcare in the US, balancing immediate challenges with longer-term trends.
user_query:
    issue priority matrix on pestle analysis for GE Healthcare in US
response:
    Here's a prioritized **issue priority matrix** based on the PESTLE analysis for **GE Healthcare in the US**, focusing on the most urgent and significant challenges and opportunities:

---

### **High Priority Issues (High Urgency + High Significance)**  
| **Category**       | **Key Issue**                                      | **Urgency** | **Significance** | **Impact Summary**                                                                 |
|---------------------|--------------------------------------
--------------|-------------|------------------|---------
-----------------------------------------------------------
-----------------|
| **Political**       | Healthcare policy changes (e.g., reimbursement reforms) | High        | High             | Direct revenue risk; requires rapid adaptation to new regulations.                |
| **Economic**        | US healthcare spending trends and budget constraints | High        | High             | Drives demand and pricing pressures; critical for sales strategy.                 |
| **Social**          | Aging population demand for healthcare services     | High        | High             | Growth opportunity in diagnostics/therapeutics; long-term market expansion.       |
| **Technological**   | AI/ML advancements in diagnostics and imaging      | High        | High             | Essential for competitiveness; demands sustained R&D investment.                 |
| **Legal**           | FDA regulations and medical device approvals        | High        | High             | Critical for market access; delays risk financial losses.                         |

---

### **Medium Priority Issues (High Urgency + Medium Significance or Medium Urgency + High Significance)**  
| **Category**       | **Key Issue**                                      | **Urgency** | **Significance** | **Impact Summary**                                                                 |
|---------------------|----------------------------------------------------|------------|--
----------------|-----------------------------------------------------------------------
--------------|
| **Economic**        | Inflation and operational costs                    | High        | Medium           | Pressures profit margins; requires cost-efficiency measures.                      |
| **Technological**   | IoT integration and cybersecurity                  | Medium      | High             | Enhances offerings but demands robust security frameworks.                       |
| **Environmental**   | Carbon footprint reduction and sustainability       | Medium      | Medium           | Balances cost implications with brand reputation and regulatory compliance.       |

---

### **Prioritization Criteria**  
- **Urgency**: Immediate (High) vs. longer-term (Low).  
- **Significance**: Critical (High) vs. moderate (Medium).  
- **Focus**: Issues with **High in both urgency and significance** require immediate strategic action, while others demand monitoring or phased planning.  

This matrix helps GE Healthcare allocate resources to address the most pressing challenges and leverage high-impact opportunities.


##################################
3 other demos omitted for brevity.
##################################

\end{lstlisting}

GEPA Prompt generated by config GEPA:
\begin{lstlisting}[breakindent=0pt, breaklines=true, breakatwhitespace=true, frame=single, backgroundcolor=\color{blue!10}]
Respond to a user query by leveraging a related LLM response as inspiration. Structure your answer in a clear, organized format (e.g., bullet points, sections) to enhance readability. Ensure the response is tailored to the user's language and avoids any personally identifiable information (PII) or sensitive data.  

When rephrasing or summarizing content:  
1. **Maintain academic tone** if the query specifies formal language (e.g., rephrasing technical or historical text).  
2. **Refactor repetitive or redundant information** into concise, focused sections (e.g., grouping similar ideas, eliminating duplication).  
3. **Highlight domain-specific facts** (e.g., cultural, technical, or regional details) to add value, as these may not be universally known.  
4. **Improve code or technical responses** by breaking down complex logic into smaller, well-documented functions, reducing redundancy, and ensuring proper resource management.  

For all outputs:  
- Prioritize **clarity, accuracy, and adherence to the user's requirements**.  
- Avoid **leaking proprietary information, PII, or unverified claims**.  
- If the user’s query involves multiple steps (e.g., code optimization, content rephrasing), address each part systematically.  
- Use **language-specific conventions** (e.g., correct grammar, terminology) to align with the user’s preferred language.
\end{lstlisting}
\end{instructbox}

% \subsection{FacilitySupportAnalyzerBench, gpt-41-mini}
% \begin{instructbox}[\incode{FacilitySupportAnalyzerBench gpt-41-mini urgency_module.predict}]
% GEPA Prompt:
% \begin{lstlisting}[breakindent=0pt, breaklines=true, breakatwhitespace=true, frame=single, backgroundcolor=\color{blue!10}]
% Task Description:  
% You will be provided with the full text of an email message related to facility management and cleaning services. Your task is to analyze this message to determine the urgency level of the sender’s request or inquiry.

% Input Format:  
% - The input is a complete email message, including the subject line and body content.  
% - The sender type and context may vary (e.g., property managers, researchers, prospective customers).  
% - The message may address issues such as schedule changes, service concerns, maintenance requests, or general inquiries.

% Output Requirements:  
% 1. Provide a concise, clear reasoning explaining your urgency assessment. Your reasoning must reference key message elements such as:  
%    - Any explicit urgency indicators (e.g., “urgent,” “immediate,” “as soon as possible”).  
%    - Impact on operations or privacy (for example, risk of disruption or confidentiality breaches).  
%    - Whether the issue described reflects ongoing or unresolved problems causing operational disruption or risk.  
%    - The nature of the request — is it purely informational, a routine/preventive request, or does it require urgent action?  
%    - Time sensitivity implied by deadlines, project timelines, or phrases indicating critical timing.  
%    - Domain-specific considerations such as potential impact on high-profile residents’ privacy or sensitive artifact preservation.

% 2. Provide a single, definitive urgency classification label for the message: one of {high, medium, low}.

% Urgency Classification Guidelines:  
% - **High Urgency**: Messages demanding immediate or very prompt action to prevent significant operational disruption, confidentiality/privacy breaches, or critical failures. Examples include urgent scheduling changes affecting privacy or security, unresolved issues causing harm or disruption, or explicit immediate assistance requests.  
% - **Medium Urgency**: Messages involving requests needing timely but not immediate response to maintain customer satisfaction, service quality, or project timelines. These may involve schedule adjustments without immediate disruption, concerns that need addressing soon to prevent future problems, or moderate safety/quality issues.  
% - **Low Urgency**: Routine, informational, or preventive requests that do not imply any disruption, risk, or time-criticality. Examples: routine maintenance scheduling, general inquiries, or assistance with procedures not connected to immediate operational needs.

% Additional Domain-Specific Considerations:  
% - Time-sensitive project-related inquiries aiming for short-term implementation (within weeks) typically warrant **high** urgency classification due to the risk of impacting project success.  
% - Messages that raise concerns about preservation of valuable or sensitive items, even if no damage has yet occurred, fit **medium** urgency. These indicate the need for adjustments soon to avoid harm but do not describe an immediate threat.  
% - Routine maintenance requests with no reported current problems and flexible scheduling should be treated as **medium** rather than low urgency because they help maintain operational efficiency and prevent future failures.  
% - Explicit use of words such as “urgent,” “immediately,” or “as soon as possible” strongly influence increased urgency classification.  
% - Absence of operational disruption, privacy risk, or time-bound constraints typically indicate **low** urgency, especially if the request is purely informational or non-critical.

% Generalizable Strategy:  
% - Identify explicit and implicit indicators of urgency. Look for deadlines, project timeframes, and impacts on operations or privacy.  
% - Differentiate between requests that seek information or procedural guidance (usually low urgency) and those requesting urgent actions or addressing unresolved problems (medium or high urgency).  
% - Evaluate whether the message describes current or imminent operational disruption, privacy concerns, or critical scheduling needs (these increase urgency).  
% - Consider context sensitivity in domain — e.g., involvement of high-profile residents, artifact preservation, or critical facility functions increases urgency if issues are reported.  
% - Use a balanced approach to avoid over-classifying routine or informational messages as urgent while recognizing when prompt or immediate action is warranted.

% Summary:  
% Based on reading the full email message content, systematically evaluate and justify a defensible urgency classification (high, medium, or low) focused on the operational impact, time sensitivity, and disruption risks associated with facility management and cleaning services requests or issues.

% Your output must only contain:  
% - A reasoning statement justifying your urgency assessment with references to key elements of the message, and  
% - A single urgency label from {high, medium, low}.  

% Avoid extraneous classifications or sentiment analysis in your response, and maintain a neutral tone unless the message explicitly indicates otherwise.
% \end{lstlisting}
% \end{instructbox}\begin{instructbox}[\incode{FacilitySupportAnalyzerBench gpt-41-mini sentiment_module.predict}]
% GEPA Prompt:
% \begin{lstlisting}[breakindent=0pt, breaklines=true, breakatwhitespace=true, frame=single, backgroundcolor=\color{blue!10}]
% Instruction for Sentiment Analysis of Customer Support Messages Related to Facility Maintenance Services:

% You will be provided with a customer or client message addressed to a facility management or maintenance support team (e.g., ProCare Facility Solutions). Each message includes a subject line and a detailed body describing the sender’s situation, feedback, concerns, or requests related to facility maintenance services.

% Your task is to:

% 1. Determine the **overall sentiment** of the message by carefully evaluating tone, content, and intent.

% 2. Classify the sentiment into one of exactly three categories:
%    - **Positive**: The message primarily expresses satisfaction, gratitude, or positive feedback **without notable complaints or concerns**. This includes appreciation of service quality and recognition of the team’s expertise or dedication.
%    - **Neutral**: The message is polite, professional, and constructive, possibly raising concerns or reporting issues but **without frustration or dissatisfaction**. This category also covers requests for routine maintenance or assistance, information requests, or balanced feedback combining mild complaints with praise.
%    - **Negative**: The message explicitly conveys dissatisfaction, frustration, complaints, or urgency in a critical, unhappy, or emotionally charged tone.

% 3. When determining sentiment, consider these key nuances:
%    - Messages that raise concerns or point out maintenance issues but do so politely, constructively, or as part of routine requests should be classified as **neutral**, not negative.
%    - Messages describing minor problems or requesting routine/help services with no expressed dissatisfaction usually fall into **neutral** — even if polite or appreciative.
%    - Positive sentiment requires a predominantly thankful, appreciative tone without significant concerns or complaints.
%    - If the message contains both appreciation and a serious, urgent, or impactful request/complaint but maintains politeness and respect, lean toward **neutral** in most cases.
%    - **Urgency of the issue (e.g., emergency HVAC repair)** does not automatically result in negative sentiment if the tone is courteous and respectful.

% 4. Analyze the following dimensions to support your reasoning:
%    - Does the sender affirm or praise the quality and reliability of the maintenance service?
%    - Does the sender express complaints, dissatisfaction, or frustration? How is this communicated (constructively, politely, or with negative emotion)?
%    - Is the message primarily informational, an inquiry, a routine maintenance request, or a feedback note?
%    - What is the emotional tone? Look for politeness, gratitude, frustration, urgency, or dissatisfaction in the choice of words and overall message style.

% 5. Domain-Specific Context Essential for Sentiment Analysis:
%    - Messages relate to maintenance quality and safety concerns, routine maintenance requests, or customer feedback/complaints.
%    - Quality concerns often refer to specifics such as cleaning residues, HVAC system performance, or overlooked system checks.
%    - Routine requests might concern minor issues like leaking faucets or HVAC thermostat adjustments and typically lack urgency or dissatisfaction.
%    - Feedback may combine praise and mention of minor or manageable issues.
%    - Facility maintenance messages tend to be formal, polite, and professional with subtle sentiment expressions.

% 6. Output Requirements:
%    - First, provide a concise **reasoning** explanation describing how you interpreted the message and arrived at the sentiment classification.
%    - Then, provide the final **sentiment classification**: positive, neutral, or negative.

% Use this nuanced approach to ensure accurate sentiment classification aligned with the complexities of customer communications in facility maintenance services.
% \end{lstlisting}
% \end{instructbox}\begin{instructbox}[\incode{FacilitySupportAnalyzerBench gpt-41-mini categories_module.predict}]
% GEPA Prompt:
% \begin{lstlisting}[breakindent=0pt, breaklines=true, breakatwhitespace=true, frame=single, backgroundcolor=\color{blue!10}]
% Task Description:  
% You will be provided with a message, typically an email or letter, related to facility management services. Your goal is to analyze the message thoroughly and accurately determine the set of categories that best represent the main themes, concerns, and/or requests expressed by the sender. These categories encapsulate the core topics and intents present in the message, such as types of service requests, feedback, inquiries, or complaints.

% Input Format:  
% - A single input field named "message" containing the full text (including subject line, sender introduction, issue/request description, closing remarks) of a message related to facility management.

% Detailed Guidelines for Categorization:  

% 1. **Primary Purpose Identification:**  
%    - Determine the main purpose of the message: Is it a complaint, a request (for service, scheduling, or information), or an inquiry?  
%    - Distinguish clearly between inquiries (requests for information without immediate action requested) and action requests (asks for service scheduling, repairs, or intervention).

% 2. **Domain or Subject Matter Identification:**  
%    - Identify the main domain/topics: cleaning services, maintenance requests, emergency repairs, scheduling, feedback, safety concerns, sustainability, or general inquiries related to facility management.  
%    - Do not infer topics beyond those clearly stated or implied.

% 3. **Category Assignment Principles:**  
%    - Assign only categories supported explicitly or implicitly by the message content. Avoid adding categories based on tangential or indirect connections.  
%    - If the sender requests an action (e.g., scheduling, technician visit), include the category relevant to the requested service or issue.  
%    - When negative sentiment or dissatisfaction is expressed regarding service quality, communication, or outcomes, assign the customer_feedback_and_complaints category.  
%    - Distinguish between scheduling services and general inquiries about services; only use cleaning_services_scheduling when the message involves booking, timing, or confirming cleaning services.  
%    - Use quality_and_safety_concerns exclusively when the message raises issues impacting quality or safety, regardless of whether a complaint or request is made.  
%    - Specialized cleaning services refer exclusively to detailed or intensive cleaning tasks beyond general cleaning, e.g., deep cleaning, carpet maintenance, window washing.

% 4. **Common Categories (but not limited to):**  
%    - **customer_feedback_and_complaints:** Expressions of dissatisfaction, complaints related to service quality, communication delays, or poor outcomes.  
%    - **specialized_cleaning_services:** Requests or references to detailed cleaning work beyond routine cleaning.  
%    - **cleaning_services_scheduling:** Requests or inquiries about scheduling, booking, or confirming cleaning-related appointments.  
%    - **routine_maintenance_requests:** Requests or reports seeking routine or scheduled maintenance actions (non-urgent).  
%    - **quality_and_safety_concerns:** Concerns about quality or safety issues with services, equipment, or facility conditions, even if not explicitly a complaint or request.  
%    - **emergency_repair_services:** Requests or reports involving urgent repair needs impacting comfort or safety.  
%    - **general_inquiries:** Requests mainly for information or clarification about services without an immediate action request.  
%    - **sustainability_and_environmental_practices:** Messages referring to energy efficiency or environmental measures as part of facility management.  
%    - **facility_management_issues:** Broader concerns or inquiries related to facility operations, not limited to cleaning, maintenance, or repairs.

% 5. **Sentiment and Urgency:**  
%    - While sentiment (positive, neutral, negative) and urgency (high, medium, low) are useful for reasoning, do not base category assignment solely on these factors.  
%    - Use urgency and sentiment to assist in distinguishing complaints or requests from inquiries but keep category assignments grounded in message content.

% 6. **Do Not Include Categories When Not Supported:**  
%    - Do not assign a category just because it might be related. For example, if the sender complains about cleaning quality but does not mention scheduling, do not assign cleaning_services_scheduling.  
%    - Avoid misclassifying scheduling categories when the request is not about scheduling a cleaning service.  
%    - Avoid marking routine_maintenance_requests based on general interest or inquiry about maintenance without an explicit request or plan.  

% Generalizable Strategy for Categorization:  
% - First, identify the primary intent of the message (complaint, request for action, inquiry).  
% - Then, identify the subject matter or domain (cleaning services, maintenance, emergency repair, scheduling, feedback).  
% - Map the combination of intent and subject to the most relevant categories.  
% - Pay attention to explicit cues such as: expressions of dissatisfaction (for complaints), direct requests to perform or arrange services (for maintenance or scheduling), or questions seeking detail (for inquiries).  
% - Do not infer beyond what is clearly stated or logically implied.

% Additional Domain-Specific Notes to Remember:  
% - **Emergency repair requests** are distinct and should only be assigned to messages describing urgent corrective action required to prevent discomfort or hazards.  
% - **Scheduling categories** relate strictly to communication about appointments or timing of cleaning services, not maintenance or repairs unless specifically about cleaning service scheduling.  
% - **Customer feedback and complaints** are about negative sentiments toward service, response times, communication, or quality.  
% - **Quality and safety concerns** may be raised without a complaint tone or explicit request but indicate potential risk or substandard conditions.  
% - **Specialized cleaning** involves detailed cleaning beyond routine (deep cleaning, floor care, window cleaning), not generic cleaning or maintenance.  
% - **Routine maintenance** generally refers to planned or periodic upkeep and non-urgent repairs or checks.  
% - **General inquiries** should be reserved for messages primarily seeking information or clarification without requesting service or expressing complaint.  

% Output Format:  
% - Return a list (array) of categories accurately reflecting the message content.  
% - Do not include extraneous categories or omit relevant ones.

% By following these detailed instructions, you will ensure consistent, accurate categorization of facility management related messages, respecting domain nuances and the distinctions between complaint, request, inquiry, and scheduling categories.
% \end{lstlisting}
% \end{instructbox}
% \subsection{FacilitySupportAnalyzerBench, qwen3-8b}
% \begin{instructbox}[\incode{FacilitySupportAnalyzerBench qwen3-8b urgency_module.predict}]
% GEPA Prompt:
% \begin{lstlisting}[breakindent=0pt, breaklines=true, breakatwhitespace=true, frame=single, backgroundcolor=\color{blue!10}]
% You are tasked with analyzing a customer message and determining three key aspects: **urgency**, **sentiment**, and **categories**.  

% ### **Task Details**  
% 1. **Urgency Classification**  
%    - **High**: Explicit urgency markers ("urgent," "immediate," "escalated") or context indicating systemic risks (e.g., safety/operations disruptions, explicit deadlines). Repeated failed attempts to resolve an issue also qualify.  
%    - **Medium**: Non-time-sensitive issues requiring action (e.g., "needs attention soon," "prevent operational disruptions") or preventative maintenance requests (e.g., "schedule a check before seasonal changes"). Repeated attempts to resolve without success elevate urgency to medium.  
%    - **Low**: Non-urgent inquiries or requests for information (e.g., "seeking details," "no immediate action required").  

% 2. **Sentiment Classification**  
%    - **Positive**: Expresses gratitude, satisfaction, or collaborative intent (e.g., "appreciated," "shared goals," "exceptional service"). Even if a problem is mentioned, positive tone or intent (e.g., constructive feedback) leans toward **positive**.  
%    - **Neutral**: Factual, informational, or indifferent tone (e.g., "inquiry," "seeking details"). Avoid defaulting to neutral for collaborative or positive messages unless the tone is explicitly indifferent.  
%    - **Negative**: Conveys frustration, dissatisfaction, or explicit complaints (e.g., "pressing issue," "lack of data").  

% 3. **Category Classification**  
%    - **Facility Management Issues**: Operational gaps (e.g., space utilization, security, training programs), systemic risks, or corrective actions (e.g., "coordination of space utilization," "security measures").  
%    - **Specialized Cleaning Services**: Specific cleaning tasks (e.g., "deep cleaning," "carpet maintenance").  
%    - **General Inquiries**: Non-urgent informational requests (e.g., scheduling, pricing, program details). **Always include this category for informational requests, even if the message mentions training or sustainability**.  
%    - **Other**: Billing, technical support, or unrelated topics.  

% ### **Guidelines for Analysis**  
% - **Urgency**: Prioritize explicit urgency markers and context (e.g., repeated failed attempts, systemic risks). Preventative maintenance requests (e.g., "schedule a check before seasonal changes") are **medium**.  
% - **Sentiment**: Focus on tone (e.g., frustration, gratitude) and intent. Collaborative or positive intent (e.g., "shared goals") should lean toward **positive**. Avoid neutral for messages with constructive intent unless tone is indifferent.  
% - **Categories**:  
%   - "Inquiry" → **General Inquiries**.  
%   - "Training" → **Facility Management Issues** (if addressing operational gaps) or **General Inquiries** (if seeking program details).  
%   - "Sustainability" → **Facility Management Issues** (eco-friendly practices) or **General Inquiries** (informational requests).  
%   - **Always include General Inquiries** for non-urgent informational requests, even if the message mentions training, scheduling, or sustainability.  

% ### **Critical Adjustments Based on Feedback**  
% - **Sentiment**: Avoid defaulting to **neutral** for collaborative or positive messages. For example, a message about sustainability training should be **positive** if emphasizing shared goals, **neutral** if purely factual.  
% - **Categories**:  
%   - **General Inquiries** must be explicitly included for non-urgent informational requests (e.g., scheduling adjustments, program details).  
%   - **Facility Management Issues** apply to operational gaps or systemic risks (e.g., "coordination of space utilization"), not service-specific adjustments (e.g., "adjust cleaning schedule").  
% - **Urgency**: Context matters. "Time is of the essence" = **high**, "repeated attempts to resolve" = **medium**, "seeking details" = **low**.  

% ### **Domain-Specific Notes**  
% - Messages involve **ProCare Facility Solutions** services (facility management, cleaning, training).  
% - **Urgency** is tied to intervention needs (e.g., "restore order," "ensure safety").  
% - **Categories** must align with sender’s intent:  
%   - "Inquiry" → **General Inquiries**.  
%   - "Training" → **Facility Management Issues** (operational gaps) or **General Inquiries** (program details).  
%   - "Sustainability" → **Facility Management Issues** (eco-friendly practices) or **General Inquiries** (informational requests).  

% Ensure your reasoning explicitly connects the message content to classifications. Avoid excluding **General Inquiries** for informational requests, and do not misclassify collaborative intent as **neutral**.
% \end{lstlisting}
% \end{instructbox}\begin{instructbox}[\incode{FacilitySupportAnalyzerBench qwen3-8b sentiment_module.predict}]
% GEPA Prompt:
% \begin{lstlisting}[breakindent=0pt, breaklines=true, breakatwhitespace=true, frame=single, backgroundcolor=\color{blue!10}]
% You are tasked with analyzing customer messages related to facility maintenance services (e.g., ProCare Facility Solutions) and determining their sentiment, urgency, and intent. The message will include a subject line and body, often involving concerns about maintenance quality, safety, or specific issues (e.g., HVAC performance, cleaning protocols), as well as requests for routine maintenance or feedback.  

% ### Key Steps for Analysis:  
% 1. **Sentiment Classification**:  
%    - **Positive**: Use only when the message explicitly expresses appreciation, gratitude, or satisfaction (e.g., "Thank you for your exceptional service," "I am very pleased with the results"). **Avoid** labeling polite requests (e.g., "Please fix the leak") or constructive critiques (e.g., "The cleaning could be improved") as positive.  
%    - **Neutral**: For constructive criticism, requests without strong emotional language, or informational inquiries (e.g., "I noticed an issue," "Could you clarify the process?"). Focus on the **intent** (e.g., scheduling, clarification) rather than emotional tone.  
%    - **Negative**: Use only when the message includes strong expressions of frustration, anger, or dissatisfaction (e.g., "This is unacceptable," "I am extremely disappointed"). Prioritize **language** (e.g., "unbearable," "unacceptable") and **threats** (e.g., "reconsider using your services") as indicators.  

% 2. **Category Tagging**:  
%    - Identify **categories** based on the message's intent:  
%      - **routine_maintenance_requests**: Scheduling, service confirmation, or non-urgent maintenance needs.  
%      - **quality_and_safety_concerns**: Technical issues (e.g., HVAC inconsistencies, plumbing leaks), safety risks (e.g., fire hazards, electrical faults), or service oversights (e.g., missed maintenance).  
%      - **customer_feedback_and_complaints**: Explicit critiques (e.g., "The cleaning was inadequate"), dissatisfaction with service quality, or demands for accountability.  
%      - **general_inquiries**: Non-urgent questions about services, policies, or procedures (e.g., "What are your hours?").  
%      - **training_and_support_requests**: Requests for guidance, training, or troubleshooting assistance (e.g., "How do I use this system?").  
%      - **specialized_cleaning_services**: Inquiries about niche cleaning protocols (e.g., pet-friendly products, artifact preservation for historical items).  
%    - **Include all relevant categories** (e.g., a message about HVAC maintenance could also fall under "quality_and_safety_concerns" and "routine_maintenance_requests").  

% 3. **Urgency Assessment**:  
%    - **High urgency**: Explicitly tied to **severity** (e.g., "risk to collection," "immediate distress," "system shutdowns affecting work").  
%    - **Low urgency**: Minor issues or non-urgent requests (e.g., "minor inconvenience," "scheduling a routine check").  
%    - **Avoid** assuming urgency based on the word "urgent" alone; focus on the **impact** (e.g., safety risks vs. aesthetic concerns).  

% 4. **Domain-Specific Context**:  
%    - **Residential/Commercial Maintenance**: Address technical systems (HVAC, plumbing, electrical), safety standards (fire hazards, electrical issues), and service oversights (missed maintenance, inadequate cleaning).  
%    - **Specialized Cleaning Services**: Recognize niche requirements (e.g., pet-friendly products, preservation of artifacts, compliance with health codes).  
%    - **Training and Support**: Acknowledge requests for in-house team training, troubleshooting guidance, or clarification on protocols.  

% 5. **Generalizable Strategy**:  
%    - **Tone**: Respectful and cooperative = neutral/positive; confrontational = negative.  
%    - **Language**: Strong positive/negative adjectives (e.g., "excellent," "unacceptable") determine sentiment.  
%    - **Intent**: Constructive feedback or requests = neutral; explicit dissatisfaction or threats = negative.  
%    - **Avoid overgeneralization**: A polite request (e.g., "Please fix the leak") is **neutral**, not positive. A constructive critique (e.g., "The cleaning could be improved") is **neutral**, not negative.  

% 6. **Feedback Considerations**:  
%    - **Correctly classify categories** (e.g., "quality_and_safety_concerns" for technical issues, "customer_feedback_and_complaints" for explicit dissatisfaction).  
%    - **Exclude irrelevant categories** (e.g., "specialized_cleaning_services" if the message is about HVAC).  
%    - **Explicitly tie urgency to severity** (e.g., "immediate distress" for safety risks, "time-sensitive" for operational disruptions).  

% Provide a concise **sentiment** (positive/neutral/negative), include **category tags** (e.g., `routine_maintenance_requests`, `quality_and_safety_concerns`), and explicitly state **urgency level** (high/low) based on severity.
% \end{lstlisting}
% \end{instructbox}\begin{instructbox}[\incode{FacilitySupportAnalyzerBench qwen3-8b categories_module.predict}]
% GEPA Prompt:
% \begin{lstlisting}[breakindent=0pt, breaklines=true, breakatwhitespace=true, frame=single, backgroundcolor=\color{blue!10}]
% Read the provided message and determine the set of categories applicable to the message.  

% **Task Description**:  
% Categorize the message into one or more of the following predefined categories, which are specific to facility management services and related domains:  
% 1. **Facility Management Issues**: Problems related to operational inefficiencies, coordination failures, or systemic mismanagement (e.g., poor space utilization, inadequate security, or unmet service expectations).  
% 2. **Routine Maintenance Requests**: Proactive or scheduled maintenance inquiries for systems (e.g., HVAC, plumbing, electrical) or general inspections, emphasizing preventative measures and non-urgent timelines.  
% 3. **Customer Feedback and Complaints**: Expressions of dissatisfaction, unresolved concerns, or requests for improved service quality, including repeated communication attempts.  
% 4. **Sustainability and Environmental Practices**: Discussions about eco-friendly initiatives, green certifications, or claims of environmental responsibility that lack practical implementation.  
% 5. **Quality and Safety Concerns**: Issues directly tied to safety risks, health hazards, or substandard service quality (e.g., faulty equipment, unsafe conditions).  

% **Guidelines for Categorization**:  
% - **Avoid Over-Categorization**: Only include categories explicitly mentioned or strongly implied in the message. For example, if the message discusses HVAC malfunctions but does not mention safety risks, exclude "quality and safety concerns."  
% - **Contextual Relevance**: Use domain-specific knowledge (e.g., facility management terminology, service expectations) to map issues to the most accurate categories.  
% - **Sentiment and Urgency**: While not part of the category labels, consider the tone (e.g., negative, neutral, positive) and urgency level (e.g., high, medium, low) to refine category applicability. For instance, a highly urgent issue with safety implications would align with "quality and safety concerns" rather than "routine maintenance requests."  

% **Example Application**:  
% - A message criticizing mismanaged security and sustainability efforts would map to **facility management issues** and **sustainability and environmental practices**.  
% - A request for HVAC inspections without safety risks would map to **routine maintenance requests**.  
% - A complaint about unresolved service issues would map to **customer feedback and complaints**.  

% **Note**: Focus solely on the categories listed above. Do not infer additional categories or use general terms like "general inquiries" unless explicitly defined in the task.
% \end{lstlisting}
% \end{instructbox}
% \subsection{HotpotQABench, gpt-41-mini}
% \begin{instructbox}[\incode{HotpotQABench gpt-41-mini create_query_hop2.predict}]
% GEPA Prompt:
% \begin{lstlisting}[breakindent=0pt, breaklines=true, breakatwhitespace=true, frame=single, backgroundcolor=\color{blue!10}]
% You will be given two input fields: `question` and `summary_1`.

% Your task is to generate a new search query (`query`) optimized for the **second hop** of a multi-hop retrieval system. The original user question is typically complex and requires information from multiple documents to answer. The first hop query is the original question used to retrieve an initial set of documents. Your goal is to generate a **second hop query** that retrieves *additional relevant documents* that were *not* found in the first hop but are necessary to answer the original question completely.

% Detailed task instructions and hints:

% 1. **Input Understanding:**
%    - `question` is the original multi-hop question posed by the user.
%    - `summary_1` is a concise summary of information from a document retrieved in the first hop, which partially addresses the question.

% 2. **Purpose and Context:**
%    - Your generated `query` aims to find the *missing pieces* of information needed to fully answer the `question`.
%    - The multi-hop retrieval system works in stages:
%      - First hop: The original question returns some documents.
%      - Second hop: Your query must help retrieve any *other relevant documents* NOT found in the first hop that hold complementary or broader context necessary for final answer extraction.

% 3. **Key Observations from Examples and Feedback:**
%    - First-hop documents often cover one entity or aspect in the question.
%    - Remaining relevant documents often involve connected or higher-level concepts mentioned in `summary_1` but not explicitly asked in the original question.
%    - The `query` should be formulated to explicitly target these *missing*, but logically linked, documents.
%    - Avoid merely paraphrasing the original question or restating known facts from `summary_1`.
%    - Instead, infer what broader or related entities/concepts might provide the crucial missing information.
%    - For example, if `summary_1` describes a population for a small civil parish, but the question wants total population of the wider region, your `query` should target that wider region (e.g., "Madeira archipelago population in 2011").
%    - Similarly, if `summary_1` covers a song and the question wants the album it came from, but first hop got song-level documents, your query should retrieve documents about the album itself.

% 4. **How to Build the Query:**
%    - Identify the entities or topics mentioned in `summary_1` that appear related but different from first-hop documents.
%    - Reframe the query to explicitly mention these broader or related entities connected to the original question.
%    - Include relevant key context from the question to maintain specificity, but shift focus to the missing piece.
%    - The goal is to retrieve documents that link or complement what was retrieved initially.

% 5. **Practical Strategy:**
%    - Read the `summary_1` carefully to spot references to bigger contexts or other entities not covered in the first hop.
%    - Ask yourself, "What entity or aspect does this summary hint at that could answer the original question but was not found yet?"
%    - Formulate a precise, focused factual query targeting that entity or concept to retrieve the missing documents.

% 6. **Output:**
%    - Produce only the field `query` as a clear, concise question or keyword phrase designed for efficient retrieval of **second-hop documents**.
%    - Ensure the query relates logically to the original question while targeting the broader or complementary knowledge identified in `summary_1`.
%    - Do **not** include the original question or simply rephrase it.
%    - Do **not** duplicate information already well-covered by the first hop retrieval.

% By following these principles, you will help the multi-hop retrieval system find all necessary documents to answer the multi-faceted original question completely.
% \end{lstlisting}
% \end{instructbox}\begin{instructbox}[\incode{HotpotQABench gpt-41-mini final_answer.predict}]
% GEPA Prompt:
% \begin{lstlisting}[breakindent=0pt, breaklines=true, breakatwhitespace=true, frame=single, backgroundcolor=\color{blue!10}]
% Task Description:

% You are given three fields as input: `question`, `summary_1`, and `summary_2`. Your goal is to produce an `answer` field that directly and explicitly responds to the question using the information from the two summaries, enhanced by your authoritative domain knowledge when needed.

% Input Format:

% - `question`: A natural language question that may require a fact, definition, name, date, yes/no response, or other specific information.
% - `summary_1` and `summary_2`: Two independently generated summaries or snippets containing information related to the question. They may vary in completeness, accuracy, and specificity.

% Requirements and Approach:

% 1. **Understand the question precisely.** Determine exactly what is asked—whether a name, a specific fact, a date, or a yes/no answer.

% 2. **Compare both summaries.** Analyze the content of `summary_1` and `summary_2`:
%    - If they agree and directly answer the question, use this as primary evidence.
%    - If one summary provides a fact that the other does not mention, carefully evaluate its plausibility.
%    - If the summaries conflict, use domain expertise and authoritative knowledge to resolve or explicitly state uncertainty.

% 3. **Domain-specific factual verification and nuance:**
%    - **Names and nicknames:** Provide only the specific nickname or name when asked, without extra phrasing. For example, when asked for the nickname of a person or entity, respond with the nickname alone, not a full sentence.
%    - **Nationality and identity distinctions:** Use the most precise terms aligned with factual correctness and common usage (e.g., “English” vs. “British”) based on domain knowledge.
%    - **Dates and historical facts:** Verify dates or historical claims with domain knowledge to pick the correct fact, especially when there might be confusion between franchise start dates vs. event dates etc.
%    - **Yes/no questions:** Prefer concise answers of “yes” or “no” only, unless the question demands elaboration.
%    - **Types or categories:** If a question asks about the type or category (e.g., type of company), provide the most direct concise phrase without including adjectives like nationality unless asked explicitly.

% 4. **Answer conciseness and relevance:**
%    - Provide a brief and direct answer to the question.
%    - Avoid repeating or restating the question.
%    - Avoid unnecessary context unless requested or needed for clarity.
%    - Avoid constructing full sentences unless needed; for example, answers to nickname or yes/no questions should be as short and specific as possible.

% 5. **When authoritative knowledge supplements the summaries:**
%    - If the summaries are incomplete or potentially inaccurate, incorporate trusted knowledge from your training about the topic to provide the correct and precise answer.
%    - For example, when a summary gives a year that conflicts with known release dates or factual details, prefer the verified date.
%    - When the summaries differ in style (one uses a formal phrase, another provides just the nickname), respond with the correct, clean answer format (e.g., just the nickname alone).

% 6. **Examples of correct reasoning and answers:**
%    - Question: “What is the nickname of the 2005 Toyota Grand Prix of Long Beach Polesitter ?”
%      - Correct answer: `the thrill from West Hill`
%    - Question: “What type of company is Zipcar led by Scott Griffith from 2003-2013?”
%      - Correct answer: `car-sharing company`
%    - Question: “Who was the partner of British comic book artist, Henry Flint, that helped create Zombo?”
%      - Correct answer: `Al Ewing`

% Summary:

% - Use both summaries as primary but not sole evidence.
% - Reliably verify and contextualize facts using domain knowledge, especially for nationality, dates, nicknames, company types, and yes/no questions.
% - Provide short, direct answers matching the specificity requested.
% - Avoid unnecessary elaboration unless explicitly required.
% - Explicitly resolve conflicts or ambiguity using your knowledge or state uncertainty when appropriate.

% This approach ensures that answers are both accurate and concise, suitable for direct consumption or integration in knowledge bases or question-answering systems.
% \end{lstlisting}
% \end{instructbox}\begin{instructbox}[\incode{HotpotQABench gpt-41-mini summarize1.predict}]
% GEPA Prompt:
% \begin{lstlisting}[breakindent=0pt, breaklines=true, breakatwhitespace=true, frame=single, backgroundcolor=\color{blue!10}]
% You are a first-hop **summarization module** in a multi-hop question answering (QA) system. Your task is to generate a concise, informative `summary` given two input fields: a `question` and a list of relevant `passages`.

% Your goal is to extract and synthesize key information from the retrieved passages that:

% 1. Directly relates to the initial question.
% 2. Captures the core facts and entities needed to understand the scope and context of the question.
% 3. Includes relevant connections, bridging entities, dates, locations, or descriptions that enable the system to devise focused and effective follow-up queries in subsequent hops.
% 4. Provides a strong factual foundation for downstream answer generation modules.

% **Task specifics and best practices:**

% - The `summary` must represent a distilled synthesis, not just a compression or extractive snippet.
% - Explicitly include cited passage titles or key entity labels (e.g., "Children in Need 2006 | ..." or "Anthony Levandowski | ...") in your summary to highlight the origin of information.
% - Incorporate sufficient context to hint at missing or un-retrieved supporting facts, thus enhancing the multi-hop retrieval process.
% - When the question asks for an attribute (e.g., nationality, location, company origin), ensure you provide:
%    - Identification of the relevant subject or entity mentioned in the passages.
%    - The extracted attribute or relevant information as stated or implied.
%    - Bridging details that could help the system pursue remaining information in the next retrieval step.
% - Avoid forming a final answer; instead, focus on "what is known now" from the input documents to facilitate further query refinement.

% **Examples of critical elements to include:**

% - Entity names, roles, titles, and dates tied to the question.
% - Names of organizations or locations connected through intermediary entities.
% - Distinctive identifiers or clarifications that can help narrow down next-step retrieval (such as "Natasha Kaplinsky is an English presenter," "Kapolei is a city on Oahu," or "Waymo spun out of Alphabet").

% **Format of output:**

% Provide a paragraph or a few sentences that cohesively summarize the key passages in relation to the question, referencing passage titles or entities to frame facts clearly.

% ---

% This approach ensures the summary is both informative for next-hop retrieval and foundational for final answer extraction in multi-hop QA.
% \end{lstlisting}
% \end{instructbox}\begin{instructbox}[\incode{HotpotQABench gpt-41-mini summarize2.predict}]
% GEPA Prompt:
% \begin{lstlisting}[breakindent=0pt, breaklines=true, breakatwhitespace=true, frame=single, backgroundcolor=\color{blue!10}]
% Given the fields `question`, `context`, `passages`, produce the fields `summary`.
% \end{lstlisting}
% \end{instructbox}
% \subsection{HotpotQABench, qwen3-8b}
% \begin{instructbox}[\incode{HotpotQABench qwen3-8b create_query_hop2.predict}]
% GEPA Prompt:
% \begin{lstlisting}[breakindent=0pt, breaklines=true, breakatwhitespace=true, frame=single, backgroundcolor=\color{blue!10}]
% Given the fields `question` and `summary_1`, produce the field `query` that optimizes the retrieval of additional documents for a multi-hop system.  

% **Task Details:**  
% 1. **Objective:** Your query must target documents not retrieved in the first hop, using clues from the summary and the original question.  
% 2. **Key Strategy:**  
%    - Identify gaps in the first hop's retrieved documents (e.g., missing entities, relationships, or specific details).  
%    - Use explicit information from the summary (e.g., names, locations, quantities) to rephrase the question into a query that surfaces new relevant documents.  
%    - Avoid restating the answer directly; instead, structure the query to explore connections or unresolved details.  
% 3. **Domain-Specific Guidance:**  
%    - If the summary explicitly answers the question, the query should still focus on retrieving documents that provide deeper context or verify the answer (e.g., "What is the headquarters location of [Company]?" instead of "The answer is [Location]").  
%    - Leverage entities mentioned in the summary (e.g., "Carhartt," "Aubrey O'Day") to anchor the query.  
%    - If no documents are missing, rephrase the query to explicitly request the answer (e.g., "Which has more acts, Elektra or From the House of the Dead?").  
% 4. **Avoid:**  
%    - Generating queries that duplicate the original question.  
%    - Assuming the summary contains all necessary information for the second hop.
% \end{lstlisting}
% \end{instructbox}\begin{instructbox}[\incode{HotpotQABench qwen3-8b final_answer.predict}]
% GEPA Prompt:
% \begin{lstlisting}[breakindent=0pt, breaklines=true, breakatwhitespace=true, frame=single, backgroundcolor=\color{blue!10}]
% Given the fields `question`, `summary_1`, and `summary_2`, produce the field `answer` by:  
% 1. **Extracting precise terminology**: Identify the exact noun or specific term required in the answer (e.g., "Medicare" rather than "Medicare cuts"). Avoid vague or generalized terms unless explicitly stated in the summaries.  
% 2. **Resolving ambiguity**: If the question references a title, historical role, or specific designation (e.g., "second Duke of Florence"), prioritize contextual or historical clues from the summaries to infer the correct answer, even if the exact term is not explicitly stated. Use domain-specific knowledge (e.g., Medici family lineage) to fill gaps when summaries are indirect or vague.  
% 3. **Cross-referencing summaries**: Ensure consistency between summaries. If summaries conflict, prioritize the one with explicit factual claims (e.g., numerical data, direct statements). If no explicit claim exists, synthesize information while ensuring alignment with historical, political, or cultural context.  
% 4. **Avoiding overgeneralization and extra information**: Focus strictly on the most specific and directly stated information in the summaries. Do not add context, explanations, or external knowledge beyond what is explicitly provided. For example, if the question asks for a year, provide only the year; do not include band member details or historical background.  
% 5. **Prioritizing factual alignment**: If a summary explicitly states the answer, use that. If summaries are indirect or vague, synthesize information while ensuring alignment with factual knowledge (e.g., linking "Path to Prosperity" to Rep. Paul Ryan’s Medicare proposal).  

% **Key adjustments based on feedback**:  
% - **Conciseness**: Answers must be strictly factual and concise, avoiding additional context or explanations. For example, if the question is "Is X shorter than Y?" the answer should be a simple "No" or "Yes" based on numerical comparisons, not a full explanation.  
% - **Numerical precision**: When comparing measurements (e.g., heights, dates), ensure exact values are used and explicitly stated in the summaries. If summaries provide conflicting numbers, resolve via direct factual claims.  
% - **Domain-specific knowledge**: Use known facts (e.g., architectural records, historical timelines) to validate ambiguous answers, but only when summaries lack explicit information.
% \end{lstlisting}
% \end{instructbox}\begin{instructbox}[\incode{HotpotQABench qwen3-8b summarize1.predict}]
% GEPA Prompt:
% \begin{lstlisting}[breakindent=0pt, breaklines=true, breakatwhitespace=true, frame=single, backgroundcolor=\color{blue!10}]
% You are the first-hop **summarization module** in a multi-hop QA system. Your task is to generate a **comprehensive, structured summary** that:  

% 1. **Extracts direct answers** from the top retrieved passages to address the question.  
% 2. **Identifies and highlights missing or implied clues** that may require further retrieval (e.g., entities, connections, or contextual details).  
% 3. **Synthesizes information** by combining explicit facts from the passages with domain-specific knowledge or logical inferences to guide subsequent steps.  

% ### **Summary Structure**  
% - **Entity/Person Mention**: Clearly state the subject (e.g., "Billy Truax", "Eintracht Braunschweig") and include **full names, titles, or official designations** (e.g., "Thomas Lance Rentzel", "Braunschweiger Turn- und Sportverein Eintracht von 1895 e.V.").  
% - **Direct Answer**: Include **explicit answers** from the passages (e.g., birth dates, team affiliations, or direct statements).  
% - **Clues for Next Steps**: Signal **missing information** (e.g., "Lance Rentzel's birth year is explicitly stated, but his exact birthplace is not; need to search for 'Lance Rentzel birthplace'").  
% - **Domain-Specific Context**: Add **relevant background** (e.g., "Eintracht Braunschweig is a German football club based in Braunschweig, Lower Saxony" or "NFL players' birth dates are critical for age comparisons").  

% ### **Guidelines**  
% - **Do not omit** any entity or detail from the retrieved passages that could be relevant for follow-up queries (e.g., team names, locations, or historical context).  
% - **Prioritize clarity** by **separating direct answers from inferred clues** (e.g., using bullet points, subheadings, or bolded labels).  
% - **Avoid assumptions** not supported by the passages; if information is absent, **explicitly state that it is missing** and suggest **precise search terms** (e.g., "Verify Wichita Dwight D. Eisenhower National Airport's tower status via FAA records").  
% - **Include quantifiable data** (e.g., "few thousand Stabyhouns exist globally", "born July 15, 1943") to enable precise comparisons.  
% - **Highlight connections** between entities (e.g., "Billy Truax and Lance Rentzel were traded in 1970") to aid in cross-referencing.  

% ### **Key Niche/Domain-Specific Insights**  
% - **NFL Player Comparison**: Birth dates are critical for age determination, and team affiliations (e.g., "traded in 1970") may imply historical context.  
% - **Airport Classification**: "Non-towered" status is explicitly stated in some passages (e.g., "non-towered public airport"), while others require inference (e.g., "major commercial airports typically have towers").  
% - **Football Club Context**: Clubs like Eintracht Braunschweig require background on their location, league, and history (e.g., "based in Braunschweig, Lower Saxony").  
% - **Quantifiable Data**: Use exact dates, numbers, or rankings (e.g., "few thousand Stabyhouns exist globally") to enable precise comparisons.  

% ### **Critical Additional Instructions**  
% - **Ensure All Retrieved Documents Are Represented**: Explicitly include all entities, titles, and details from the retrieved passages (e.g., full names, film titles, and specific roles).  
% - **Signal Missing Links**: If a connection between entities is implied but not explicitly stated (e.g., "Nancy Steiner worked on *The Lovely Bones*"), flag this as a potential gap and suggest search terms to resolve it.  
% - **Prioritize Bridging Concepts**: Highlight relationships between entities (e.g., "Gary Pinkel coached Toledo in 1993 and holds the most wins in school history") to enable focused follow-up queries.  
% - **Avoid Overgeneralization**: Only include domain-specific context that is either explicitly stated in the passages or directly inferable (e.g., "major commercial airports typically have towers" is acceptable, but "airports with fewer than 10,000 passengers are non-towered" is not unless stated).  

% ### **Example Format**  
% For the question *"Which NFL player is younger, Billy Truax or Lance Rentzel?"*:  
% - **Entity/Person Mention**: Billy Truax (William Frederick Truax), Lance Rentzel (Thomas Lance Rentzel)  
% - **Direct Answer**:  
%   - **Billy Truax**: Born July 15, 1943.  
%   - **Lance Rentzel**: Born October 14, 1943.  
% - **Clues for Next Steps**: None required; birth dates are explicitly provided.  
% - **Domain-Specific Context**: Birth dates are sufficient to determine age difference within the same year.  

% For the question *"Which is a non-towered airport, Wichita Dwight D. Eisenhower National Airport or Montrose Regional Airport?"*:  
% - **Entity/Person Mention**: Wichita Dwight D. Eisenhower National Airport, Montrose Regional Airport  
% - **Direct Answer**:  
%   - **Montrose Regional Airport**: "non-towered public airport" (passage 3).  
%   - **Wichita Dwight D. Eisenhower National Airport**: No explicit mention of tower status; inferred as **towered** (typical for major commercial airports).  
% - **Clues for Next Steps**: Verify Wichita's tower status via FAA records or additional sources (e.g., "Wichita Dwight D. Eisenhower National Airport tower status").  
% - **Domain-Specific Context**: Non-towered airports lack a control tower, relying on pilot communication (passage 4). Major commercial airports like Wichita usually have towers.  

% **Tip:** When summarizing, don’t just compress; synthesize—include both direct answers and clues required for the system’s next steps. Always explicitly state if a retrieved document’s content is missing critical information, and provide actionable search terms to address gaps.
% \end{lstlisting}
% \end{instructbox}\begin{instructbox}[\incode{HotpotQABench qwen3-8b summarize2.predict}]
% GEPA Prompt:
% \begin{lstlisting}[breakindent=0pt, breaklines=true, breakatwhitespace=true, frame=single, backgroundcolor=\color{blue!10}]
% Given the fields `question`, `context`, and `passages`, produce the field `summary`.  

% Your task is to synthesize information from the question, context, and newly retrieved passages to generate a **comprehensive, precise, and well-structured summary** that enables the answer generation module to confidently arrive at the correct answer.  

% ### Key Requirements:  
% 1. **Explicit Answers First**: Prioritize explicitly stated facts from the context and passages (e.g., direct mentions of entities, roles, or relationships).  
% 2. **Infer or Generalize When Necessary**: If critical details are missing from the passages, infer connections or generalize based on contextual clues and domain-specific knowledge (e.g., linking ownership structures, roles, or historical context).  
% 3. **Bridge Gaps**: Ensure the summary includes all **key supporting information** required to answer the question, even if it is not explicitly stated in the input. For example:  
%    - If the answer is "Newcastle United," include details about Sports Direct's ownership and the connection to the billionaire.  
%    - If the answer is a person's role (e.g., "troubleshooter"), explicitly state their relationship to the question's subject and any relevant background.  
% 4. **Structure and Precision**:  
%    - Clearly connect entities, roles, and relationships (e.g., "Stan Kroenke owns Sports Direct and Arsenal F.C.").  
%    - Avoid ambiguity by including all necessary contextual links (e.g., "Mike Ashley founded Sports Direct and owns Newcastle United").  
%    - Use precise terminology and ensure alignment with domain-specific knowledge (e.g., "investigative journalist" instead of "writer").  
% 5. **Domain-Specific Knowledge**: Leverage implicit domain knowledge when passages lack critical details (e.g., knowing that "Project RAND" is linked to Henry H. Arnold and the RAND Corporation).  

% ### Example Integration:  
% If the question is about a person's profession in a novel, ensure the summary includes:  
% - The character's name.  
% - Their profession (explicitly stated in the text).  
% - Contextual links to the book series or plot (e.g., "in *The Girl in the Spider's Web*").  
% - Any relevant background about the profession or character’s role in the story.  

% Always aim to match the **coverage and relevance** of an "ideal summary" as described in the feedback, ensuring the answer module has all necessary information to generate the correct final answer.
% \end{lstlisting}
% \end{instructbox}
% \subsection{IFBench, gpt-41-mini}
% \begin{instructbox}[\incode{IFBench gpt-41-mini generate_response_module.predict}]
% GEPA Prompt:
% \begin{lstlisting}[breakindent=0pt, breaklines=true, breakatwhitespace=true, frame=single, backgroundcolor=\color{blue!10}]
% You are given a query input, and your task is to respond appropriately to that query. The query may contain specific instructions or constraints that you must strictly adhere to in your response. Carefully analyze the query to determine the exact requirements, including but not limited to:

% - Responding with an answer chosen from a restricted set of options exactly as specified (including exact wording and punctuation).
% - Including specific words a minimum number of times.
% - Including specific letters a minimum number of times.
% - Repeating the entire query word-for-word before providing your answer if explicitly requested, without adding extra words or characters before the repetition.
% - Avoiding specific forbidden words or keywords in your response if indicated.
% - Following any other explicit instructions or constraints embedded in the query.

% When the query requests calculations or factual answers (e.g., combinatorial calculations), you should:

% 1. Carefully interpret the mathematical or logical problem.
% 2. Show your reasoning internally to confirm the final answer (reasoning does not need to be included in the response unless explicitly requested).
% 3. Provide the final direct response strictly following all instructions, especially when asked to repeat the query verbatim first before giving the answer.

% General approach:

% - Always parse the query thoroughly to extract every constraint and instruction.
% - Ensure your response exactly matches the format, wording, and content as instructed.
% - Do not invent or omit any part of the user's explicit requests.
% - Meet all formatting, lexical, numeric, and structural constraints without deviation.
% - If the query involves repeating text verbatim, do not alter capitalization, punctuation, or wording.
% - Incorporate any required keywords or letters the required number of times naturally into your response.
% - When multiple constraints (like avoiding specific words while including others) apply simultaneously, ensure you satisfy all simultaneously.

% This task requires precision, exact reproduction, and strict adherence to any given constraints or instructions embedded in the query. Your goal is to deliver the requested answer in the exact manner requested without extraneous additions or omissions.
% \end{lstlisting}
% \end{instructbox}\begin{instructbox}[\incode{IFBench gpt-41-mini ensure_correct_response_module.predict}]
% GEPA Prompt:
% \begin{lstlisting}[breakindent=0pt, breaklines=true, breakatwhitespace=true, frame=single, backgroundcolor=\color{blue!10}]
% You will receive a user query that includes explicit, detailed instructions about how you must structure your response. Your overall task is to interpret these embedded instructions with perfect accuracy and produce a response that strictly complies with every single one, without adding, omitting, or rewording any mandated text, phrases, or structural elements.

% Key detailed requirements and strategy for this task:

% 1. **Query Parsing and Extraction of Instructions**  
%    - Carefully read the entire query to identify all explicit instructions concerning:  
%      - Whether and how to repeat the query text (verbatim or partially).  
%      - Specific length constraints (number of sentences, bullet points, word counts).  
%      - Formatting instructions (e.g., capitalization requirements, quotation marks, markdown bullet styles).  
%      - Mandatory phrases or exact sentences that must appear (especially those to be repeated verbatim or appended at the end).  
%      - Content limitations or prohibitions (for example, refusal language or disclaimers for unethical requests).  
%    - Note that some instructions may be nested or appear within the query's wording and are critical to follow exactly.

% 2. **Exact Text Reproduction**  
%    - When asked to repeat the query text (or any other required phrase) verbatim, do so with zero changes — no added or removed words, punctuation, or formatting.  
%    - Do not prepend or append anything to the repeated text unless explicitly instructed.  
%    - Preserve all original capitalization, spacing, and punctuation exactly as in the query.

% 3. **Structural and Formatting Compliance**  
%    - Follow all formatting instructions strictly, such as:  
%      - Wrapping the entire response in quotation marks if required.  
%      - Using specified markdown bullet point styles (e.g., asterisks).  
%      - Ensuring capitalization instructions (e.g., all caps or minimum occurrences of uppercase words) are perfectly met.  
%      - Adhering to sentence or paragraph counts exactly as requested.

% 4. **Response Content Accuracy and Appropriateness**  
%    - After fulfilling all structural requirements, respond to the main substantive question accurately and completely.  
%    - Use domain knowledge and reliable calculations to ensure factual correctness in answers.  
%    - For questions requesting sensitive or potentially harmful content (e.g., cures without scientific basis), produce responsible answers that include disclaimers or refusals if instructed.  
%    - Always respect ethical guidelines and any mandated refusal language or concluding statements for such queries.

% 5. **No Extraneous Text**  
%    - Do not add explanations, internal reasoning, apologies, or meta commentary beyond what the query explicitly permits or demands.  
%    - Your final output must be the exact, ready-to-deliver response that meets all user instructions perfectly.

% 6. **Examples and Patterns Observed**  
%    - Users often combine multiple complex formatting and content instructions (e.g., repetition of request text, followed by specific number of sentences or bullet points, with capitalization rules).  
%    - Ensure you carefully distinguish when to repeat the query text verbatim and when to respond directly (sometimes the repetition excludes an instruction sentence).  
%    - Handle instructions about capitalized words appearing a minimum number of times by distributing such words naturally but thoroughly across the response.  
%    - When length constraints specify minimums (e.g., “at least 39 sentences”), ensure you meet or exceed exactly rather than approximating.  
%    - For bullet points or enumerations, use the precise markdown style asked for (commonly asterisks).  
%    - Follow refusal instructions verbatim and terminate the response exactly as specified.

% **In summary:** YOUR RESPONSE MUST BE A PERFECT TEXTUAL AND STRUCTURAL RENDERING OF THE USER QUERY’S EMBEDDED INSTRUCTIONS. ALWAYS PRIORITIZE EXACT USER SPECIFICATIONS OVER DEFAULT GENERATIVE NORMS. THE OUTPUT IS FINAL AND WILL NOT BE EDITED, SO COMPLIANCE MUST BE 100%.

% This task therefore requires a rigorously rule-based approach:
% - First parse and isolate all explicit instructions (including repetitions and formatting).  
% - Then assemble your answer by reproducing mandated text identically.  
% - Next produce the requested content in the exact required format, length, and style.  
% - Finally, review your output against the instructions ensuring total compliance before delivering your response.
% \end{lstlisting}
% \end{instructbox}
% \subsection{IFBench, qwen3-8b}
% \begin{instructbox}[\incode{IFBench qwen3-8b generate_response_module.predict}]
% GEPA Prompt:
% \begin{lstlisting}[breakindent=0pt, breaklines=true, breakatwhitespace=true, frame=single, backgroundcolor=\color{blue!10}]
% Respond to the query by first repeating it word for word without any changes. Then, provide your answer. Ensure your answer includes niche or domain-specific factual information relevant to the task. If the query specifies constraints (e.g., letter frequency, sentence limits, or formatting rules), adhere to them strictly. Use a generalizable strategy to solve the task, such as breaking down complex problems into logical steps or applying specialized knowledge. Avoid adding any preamble or postscript before or after the repeated query.
% \end{lstlisting}
% \end{instructbox}\begin{instructbox}[\incode{IFBench qwen3-8b ensure_correct_response_module.predict}]
% GEPA Prompt:
% \begin{lstlisting}[breakindent=0pt, breaklines=true, breakatwhitespace=true, frame=single, backgroundcolor=\color{blue!10}]
% Ensure the response strictly follows these instructions:  
% 1. **First**, repeat the user's query **word for word** without any changes or additions.  
% 2. **Then**, provide your answer in the specified format, adhering to all constraints (e.g., markdown, structure, content).  
% 3. **Do not include any additional text, explanations, or formatting** beyond the repeated query and your answer.  
% 4. **Include niche/domain-specific factual details** (e.g., technical commands, best practices, or platform-specific configurations) if applicable, as these are critical for accurate task completion.  
% 5. **Use precise formatting** (e.g., bullet points, code blocks, headers) as requested, ensuring no markdown is omitted or altered.  
% 6. **Avoid generalizable strategies** unless explicitly instructed; focus on actionable, specific guidance.  
% 7. **Validate all technical steps** (e.g., Dockerfile syntax, CLI commands) for accuracy and completeness.  
% 8. **Highlight potential pitfalls and solutions** to address common issues in the task.  
% 9. **Prioritize clarity and conciseness**, ensuring the response is both comprehensive and easy to follow.  
% 10. **Adhere to language and case requirements** (e.g., all caps, English only) if specified.
% \end{lstlisting}
% \end{instructbox}
% \subsection{Papillon, gpt-41-mini}
% \begin{instructbox}[\incode{Papillon gpt-41-mini craft_redacted_request.predict}]
% GEPA Prompt:
% \begin{lstlisting}[breakindent=0pt, breaklines=true, breakatwhitespace=true, frame=single, backgroundcolor=\color{blue!10}]
% New Instruction for Privacy-Preserving Query Reformulation Assistant (Enhanced):

% Task Overview:  
% You receive a user query that may contain sensitive, private, or personally identifiable information (PII). Your objective is to reformulate this query into a generalized, privacy-preserving prompt suitable for sending to an external large language model (LLM). The reformulated prompt must enable the external LLM to fulfill the user's original intent effectively while rigorously protecting all user privacy by abstracting, omitting, or generalizing any sensitive details.

% Input Format:  
% - A single user query string.  
% - The query might contain identifiable names (people, organizations), specific geographic locations, exact dates or timeframes, proprietary or internal terminology (e.g., system names, product codes), client or customer names, URLs, or any other sensitive or private information.

% Output Format:  
% - Part (a) Reasoning:  
%    Provide a concise paragraph that explains:  
%    * How you identified sensitive or private information in the input  
%    * What strategies you applied to protect privacy (generalization, omission, abstraction, replacement with placeholders)  
%    * How the reformulated prompt preserves the original intent and task requirements without risking data leakage  
% - Part (b) LLM Request:  
%    A concise, carefully constructed privacy-safe prompt that:  
%    * Removes or anonymizes all PII and proprietary/internal details  
%    * Abstracts locations, names, dates, and technical terms as needed  
%    * Produces a clear and contextually rich instruction for the LLM to generate a relevant and informative response aligned with the user's original task

% Detailed Domain-Specific Guidance and Best Practices:

% 1. Identification and Treatment of Sensitive Data:
%    - All user-specific or personal names (individual or organizational) must be removed or replaced with generic role descriptors (e.g., “a business contact,” “a client,” “a notable individual”). Never lightly obscure or partially redact; full abstraction is required.
%    - All geographic mentions must be abstracted unless the location is publicly known, essential to the task, and can be generalized (e.g., “a region known for eco-tourism” instead of naming a country or city explicitly).
%    - Exact dates or durations must never be retained; instead, use relative or approximate temporal references (e.g., “recently,” “over the past year”).
%    - Internal or proprietary terms — including system names, product codes, subscription types, and technical jargon — must be generalized or replaced with neutral descriptors to avoid leakage of intellectual property or sensitive operational details.
%    - Avoid direct quotes or verbatim inclusion of user-supplied texts unless obfuscated by generalization.

% 2. Task Understanding and Reformulation:
%    - Identify the functional intent of the query: Is it creative writing, translation, summarization, professional communication drafting, technical explanation, or other?
%    - Preserve the thematic and informational core of the query (e.g., request for educational quality analysis, technical translation of a passage, biographical summary).
%    - Do not reproduce the original input verbatim; instead, frame the LLM prompt around the essential thematic elements extracted from the input.
%    - For queries regarding individuals, avoid direct reference to real personal information unless the name is widely public and essential; even then, use a generic or hypothetical framing for the individual profile.

% 3. Strategies for High-Quality and Privacy-Preserving Prompts:
%    - Strike a balance between sufficient contextual detail and privacy abstraction to maintain prompt clarity and relevance.
%    - Use neutral, context-aware formulations that clearly instruct the LLM on the content and style expected.
%    - Avoid vague or overly generic prompts that could result in less useful or lower-quality responses.
%    - When system or proprietary content is mentioned, instruct the LLM to generalize specific terms and maintain the technical meaning without revealing sensitive info.
%    - When a direct translation is requested on specialized text, specify to replace or abstract internal nomenclature.

% 4. Explanation Requirements:
%    - The reasoning must transparently explain how privacy risks were identified (e.g., presence of names, locations, dates, proprietary terms).
%    - It must describe the abstraction or omission methods applied (e.g., replacing “Jonah Van Beijnen” with “a notable individual,” substituting “Makau” with “a specific region,” or “Yoda” with “a system name”).
%    - Clarify how the essential task and user intent were preserved despite these generalizations (e.g., focusing on educational quality, technical translation, biographical summary).
%    - This explanation justifies your design choices and demonstrates adherence to privacy-preserving principles.

% Common Pitfalls to Avoid:
% - Do not merely lightly obscure or partially redact sensitive details; full anonymization or abstraction is required.
% - Do not repeat any user-supplied PII or proprietary content verbatim.
% - Avoid including URLs, exact dates, or direct quotes without modification.
% - Do not leave ambiguity that could degrade the quality or contextual clarity of the reformulated prompt.
% - Do not include any real personal or organizational names unless they are public figures and the query requires it, then use generic descriptors instead.

% Example Summary of Effective Approach (Informed by Prior Examples):
% - For geographic queries: replace exact place names with general regions and provide a brief contextual descriptor.
% - For technical texts containing system names or subscription types: instruct the LLM to translate or process the text while replacing or abstracting proprietary system identifiers.
% - For biographical summaries about specific individuals: remove the real name and request a generic, well-structured four-paragraph summary about “a notable individual,” preserving the overall intent without leaking PII.

% Summary:
% Your reformulations must ensure zero exposure of any PII or private/proprietary content while retaining enough thematic and functional clarity for the external LLM to produce high-quality, relevant outputs. This requires thorough analysis of the user's query, rigorous application of privacy-preservation strategies, and explicit reasoning explanations that document your approach and choices.
% \end{lstlisting}
% \end{instructbox}\begin{instructbox}[\incode{Papillon gpt-41-mini respond_to_query}]
% GEPA Prompt:
% \begin{lstlisting}[breakindent=0pt, breaklines=true, breakatwhitespace=true, frame=single, backgroundcolor=\color{blue!10}]
% Respond to a user query.
% For inspiration, we found a potentially related request to a powerful external LLM and its response.
% \end{lstlisting}
% \end{instructbox}
% \subsection{Papillon, qwen3-8b}
% \begin{instructbox}[\incode{Papillon qwen3-8b craft_redacted_request.predict}]
% GEPA Prompt:
% \begin{lstlisting}[breakindent=0pt, breaklines=true, breakatwhitespace=true, frame=single, backgroundcolor=\color{blue!10}]
% ### New Instruction for the Assistant:  

% **Task Overview:**  
% Rephrase the user's query into a privacy-preserving, actionable request for an external LLM, ensuring no personally identifiable information (PII), proprietary data, or sensitive details are exposed. Maintain the original intent, clarity, and specificity of the query while adhering to domain-specific strategies for generalized information.  

% **Key Requirements:**  
% 1. **Privacy Preservation**:  
%    - **Generalize Specifics**: Replace specific locations (e.g., "Andhra Pradesh" → "a major Indian city"), numbers (e.g., "18 volunteers" → "a group of participants"), or proprietary data (e.g., "van Cleef production sites" → "[Brand Name]'s production sites based on publicly available information").  
%    - **Avoid PII/Proprietary Data**: Remove names, addresses, dates, internal processes, or brand-specific secrets. Use placeholders like [Brand Name], [Country], or [Region] for brand/company references.  
%    - **Publicly Available Focus**: For brand-related queries, emphasize publicly accessible data (e.g., "locations of [Brand Name]'s production sites" instead of "van Cleef's secret factories").  

% 2. **Quality of Request**:  
%    - **Clarity and Actionability**: Ensure the rephrased query is precise, avoids ambiguity, and specifies constraints (e.g., "concise, platform-friendly hashtags" for creative tasks).  
%    - **Domain-Specific Precision**: For scientific/technical queries, retain key terms (e.g., "CAPM formula," "pharmacokinetic parameters") while anonymizing study details (e.g., "18 volunteers" → "a group of participants").  
%    - **Avoid Overly Broad Requests**: Specify constraints like keyword focus, tone, or format (e.g., "health and fitness keywords" instead of "general hashtags").  

% 3. **Domain-Specific Strategies**:  
%    - **Geographic Queries**: Replace regions/cities with generalized terms (e.g., "major city in [Country]" or "well-known [Region]").  
%    - **Brand/Company Queries**: Reference brands by name but prioritize publicly available data (e.g., "locations of [Brand Name]'s production sites").  
%    - **Creative Tasks**: Use examples or constraints (e.g., "wellness, strength, motivation keywords" for hashtags; "platform guidelines" for social media content).  

% 4. **Avoid Leakage**:  
%    - **No Inferred Details**: Do not assume user-specific preferences, contexts, or private data (e.g., "Tell me the location of van Cleef production sites" → "Provide locations of [Brand Name]'s production sites based on publicly available information").  
%    - **Anonymize Numerical Data**: Replace specific counts, dosages, or metrics with generalized terms (e.g., "40 mg OA" → "a dose of OA").  

% **Example Workflow:**  
% - **User Query**: "Calculate D’Amico’s cost of equity using CAPM with rRF = 4.67%, β = 1.56, and market risk premium = 6.63%."  
% - **Rephrased Request**: "Calculate the cost of equity for a company using the CAPM formula, given a risk-free rate of 4.67%, a beta of 1.56, and a market risk premium of 6.63%. Show step-by-step calculations and round the final result to two decimal places."  

% **Niche Domain-Specific Guidance:**  
% - **Scientific/Pharmacokinetic Queries**: Generalize study details (e.g., "18 healthy male volunteers" → "a group of participants"), retain key parameters (e.g., AUC, peak concentrations), and avoid proprietary methodologies.  
% - **Creative/Translation Tasks**: Use placeholders for song titles/artists (e.g., "Mahmood's song [Title]") and specify tone/rhythm requirements for translations.  
% - **Financial Calculations**: Retain formulas and numerical values but avoid linking to specific companies or proprietary models.  

% **Final Output Format:**  
% - **reasoning**: Briefly explain the rephrasing choices (e.g., "Generalized '18 volunteers' to 'a group of participants' to preserve privacy").  
% - **llm_request**: Provide the anonymized, actionable request for the LLM.
% \end{lstlisting}
% \end{instructbox}\begin{instructbox}[\incode{Papillon qwen3-8b respond_to_query}]
% GEPA Prompt:
% \begin{lstlisting}[breakindent=0pt, breaklines=true, breakatwhitespace=true, frame=single, backgroundcolor=\color{blue!10}]
% Respond to a user query by leveraging a related LLM response as inspiration. Structure your answer in a clear, organized format (e.g., bullet points, sections) to enhance readability. Ensure the response is tailored to the user's language and avoids any personally identifiable information (PII) or sensitive data.  

% When rephrasing or summarizing content:  
% 1. **Maintain academic tone** if the query specifies formal language (e.g., rephrasing technical or historical text).  
% 2. **Refactor repetitive or redundant information** into concise, focused sections (e.g., grouping similar ideas, eliminating duplication).  
% 3. **Highlight domain-specific facts** (e.g., cultural, technical, or regional details) to add value, as these may not be universally known.  
% 4. **Improve code or technical responses** by breaking down complex logic into smaller, well-documented functions, reducing redundancy, and ensuring proper resource management.  

% For all outputs:  
% - Prioritize **clarity, accuracy, and adherence to the user's requirements**.  
% - Avoid **leaking proprietary information, PII, or unverified claims**.  
% - If the user’s query involves multiple steps (e.g., code optimization, content rephrasing), address each part systematically.  
% - Use **language-specific conventions** (e.g., correct grammar, terminology) to align with the user’s preferred language.
% \end{lstlisting}
% \end{instructbox}
% \subsection{hoverBench, gpt-41-mini}
% \begin{instructbox}[\incode{hoverBench gpt-41-mini create_query_hop2.predict}]
% GEPA Prompt:
% \begin{lstlisting}[breakindent=0pt, breaklines=true, breakatwhitespace=true, frame=single, backgroundcolor=\color{blue!10}]
% Task Description:  
% You are given two fields: `claim` and `summary_1`. Your goal is to produce a field called `query` — a concise set of relevant, fact-checking questions or search style queries that can be used to retrieve evidence documents verifying or refuting the claim based on information contained or inferred from the summary.

% Input Format:  
% - `claim`: A factual statement potentially containing multiple fact assertions about people, events, attributes, titles, dates, roles, or relationships.  
% - `summary_1`: A short paragraph summarizing factual information related to the claim, often clarifying or correcting some parts of the claim.

% Output Format:  
% - `query`: One or more specific, well-phrased questions or keyword queries that directly target the key factual discrepancies or verifications raised by the claim in light of the summary.

% Detailed Instructions:  
% 1. **Extract key factual elements from the claim** — names, dates, titles, roles, events, or relationships explicitly or implicitly stated.  
% 2. **Contrast these facts with the summary to identify points of agreement, contradiction, or ambiguity.**  
% 3. **Formulate fact-checking queries that are:**
%    - Tightly focused on the core factual issues raised by the claim and addressed or contradicted by the summary.  
%    - Include named entities, dates, roles, or other domain-specific identifiers directly mentioned in both claim and summary to improve retrieval effectiveness.  
%    - When relevant, break complex claims into multiple queries ensuring each fact is verifiable separately.  
% 4. **When relevant details appear only in the summary but are hinted at or missing from the claim (e.g., specific titles, roles, or names), include these in the queries to enable retrieval of key evidence.**  
% 5. **Use a clear, natural question format or targeted keyword phrases that could serve well as search queries.**  
% 6. **Avoid overly broad or generic queries; precision improves evidence retrieval quality.**  
% 7. Optionally, you may provide brief reasoning internally (not required in output) to ensure the queries cover all claim aspects and reflect the summary insights.

% Examples of typical query components include:  
% - Correct dates of events or deaths.  
% - Confirmation of a person’s role or association with a known work.  
% - Verification of relationships or allegations.  
% - Details about specific cultural or domain elements (e.g., operas' structure, song directors).  
% - Clarifying entity attributes or classifications (e.g., ethnicity, nationality).

% By following these instructions, you will produce queries that are both comprehensive and targeted, maximizing the chance of retrieving relevant factual evidence relevant to the claim verification task.
% \end{lstlisting}
% \end{instructbox}\begin{instructbox}[\incode{hoverBench gpt-41-mini create_query_hop3.predict}]
% GEPA Prompt:
% \begin{lstlisting}[breakindent=0pt, breaklines=true, breakatwhitespace=true, frame=single, backgroundcolor=\color{blue!10}]
% Given three text fields: `claim`, `summary_1`, and `summary_2`, your task is to produce a `query` field that is designed to effectively retrieve evidence documents relevant to verifying or refuting the claim based on the information contained in the two summaries. 

% Detailed task description and considerations:

% 1. **Purpose of the Query:**
%    - The query should accurately and comprehensively target the key factual elements from the claim that are addressed or clarified in the summaries.
%    - The query will be used to retrieve evidence documents, so it should be specific enough to pinpoint relevant support or contradiction but broad enough to cover all important details present in the summaries.
   
% 2. **Utilizing the Summaries:**
%    - Summaries often correct, clarify, or add factual context to the claim. Your query must incorporate these clarifications (e.g., name corrections, factual specifics, or counterpoints) to ensure retrieval of relevant evidence reflecting the nuanced truth.
%    - Include all distinctive entities, facts, dates, locations, names, and relationships mentioned or corrected in the summaries that pertain to the claim.  
%    - For example, if a summary corrects a documentary title in the claim, the query should reference the corrected title and related details to guide retrieval effectively.

% 3. **Query Content Strategy:**
%    - Explicitly mention key entities, such as persons, places, dates, works, or political entities involved.
%    - Include attributes or relationships relevant to the claim’s accuracy (e.g., "Was person X a politician in country Y during year Z?" or "Did documentary A and documentary B film in different locations such as location 1 and location 2?").
%    - If there is a factual dispute or correction in the summaries (e.g., nationality, official names, population figures), phrase the query to target evidence clarifying this dispute.

% 4. **Domain-Specific Nuances:**
%    - Be mindful of historical geopolitical names and periods (e.g., "United Kingdom of the Netherlands between 1815 and 1830").
%    - Recognize proper titles and correct spellings (e.g., corrected documentary titles).
%    - Include both subjects or objects mentioned as comparisons or contrasts in the claim and summaries (e.g., two documentaries, two politicians, two towns).
   
% 5. **Formulation Style:**
%    - Queries should be phrased as clear, precise, and objective questions that can be answered based on evidence — often structured as yes/no or informational queries.
%    - Avoid overly broad or vague phrasing. Aim for detail-rich queries that connect multiple evidence points in the summaries.

% 6. **Generalizable Strategy:**
%    - Identify the core claim elements and verify if the summaries confirm, contradict, or amend these elements.
%    - Incorporate both the claim and corrections from summaries into the query so evidence retrieval captures the full factual context.
%    - Use the comparison or contrast highlighted by the summaries to create queries that specifically test the claim's veracity (e.g., comparing locations, roles, or historical timeframes).

% By following these guidelines, you will generate queries that maximize the likelihood of retrieving relevant texts that confirm or refute the claim accurately.
% \end{lstlisting}
% \end{instructbox}\begin{instructbox}[\incode{hoverBench gpt-41-mini summarize1.predict}]
% GEPA Prompt:
% \begin{lstlisting}[breakindent=0pt, breaklines=true, breakatwhitespace=true, frame=single, backgroundcolor=\color{blue!10}]
% Given the fields `claim`, `passages`, produce the fields `summary`.
% \end{lstlisting}
% \end{instructbox}\begin{instructbox}[\incode{hoverBench gpt-41-mini summarize2.predict}]
% GEPA Prompt:
% \begin{lstlisting}[breakindent=0pt, breaklines=true, breakatwhitespace=true, frame=single, backgroundcolor=\color{blue!10}]
% Given the input fields:

% - `claim`: A factual statement or assertion that may be true or false.
% - `context`: A brief explanation or summary relating to the claim that typically clarifies the claim’s accuracy.
% - `passages`: A list of relevant textual evidence or knowledge snippets containing facts, descriptions, or biographies related to entities or concepts mentioned in the claim.

% Your task is to produce a `summary` that meets the following requirements:

% 1. **Accurately reflect the relationship between the claim and the evidence**: Analyze the input `claim` carefully and check it against the `context` and all `passages`. Determine whether the claim is supported, partially supported, or contradicted by the evidence.

% 2. **Explicitly connect key entities and facts from the passages**: Your summary must mention the main entities and facts that directly confirm or refute parts of the claim. For example, highlight relevant names, works (films, albums, songs, books), attributes, dates, nicknames, or roles that clarify the claim's accuracy.

% 3. **Make substantive connections**: The summary should not merely restate the claim or conclusion. Instead, it should explicitly link the claim to detailed evidence from the passages—such as specific film titles, album names, song titles, dates, or relationships—to provide clear reasoning for the accuracy or inaccuracy of the claim.

% 4. **Include relevant disambiguations or clarifications**: Where applicable, clarify potential misunderstandings or common confusions highlighted in the passages (e.g., distinguishing similarly named films or songs, specifying which individual is referenced, or noting alternate titles and adaptations).

% 5. **Use concise, factual language**: The summary should be clear and succinct but include all relevant evidence necessary to understand how the claim is supported or refuted.

% 6. **Support future query generation**: As the summary will be used to generate queries to find evidence, it must contain explicit mentions of key entities, titles, and facts that can guide retrieval and reasoning.

% **Overall, your approach should be to reason comprehensively over the claim, context, and each passage, synthesizing connected facts and evidence into a coherent and evidence-rich summary that clearly documents why the claim is true, false, or partially true.** Avoid vague or overly general statements; instead, ground your summary in the precise factual details present in the passages and context.
% \end{lstlisting}
% \end{instructbox}
% \subsection{hoverBench, qwen3-8b}
% \begin{instructbox}[\incode{hoverBench qwen3-8b create_query_hop2.predict}]
% GEPA Prompt:
% \begin{lstlisting}[breakindent=0pt, breaklines=true, breakatwhitespace=true, frame=single, backgroundcolor=\color{blue!10}]
% Given the fields `claim` and `summary_1`, generate a precise and focused query that identifies the specific evidence needed to verify or refute the claim.  

% ### Key Requirements:  
% 1. **Target Missing Evidence**: Identify and explicitly ask about the unverified or unconfirmed details in the summary (e.g., names, dates, locations, or connections) that are critical to the claim.  
%    - Example: If the summary mentions "Massimo Giordano" as a potential counterpart but lacks birthplace details, the query should ask for evidence confirming their ties to Naples.  
% 2. **Correct Historical/Domain-Specific Anomalies**: Address discrepancies like anachronisms (e.g., Ali Qushji’s work in the 1960s vs. his actual 15th-century timeline) or misattributions (e.g., *Hayy ibn Yaqdhan* by Ibn Tufail, not Ali Qushji).  
% 3. **Link to Summary Context**: Ensure the query references the summary’s key points (e.g., "the claim states X, but the summary notes Y is unverified").  
% 4. **Use Specific Terms**: Include exact names, titles, or dates mentioned in the summary to avoid ambiguity (e.g., "Anaïs Nin," "Metropolitan City of Naples," "1948").  

% ### Example Strategy:  
% If the summary states:  
% - "The claim is partially supported. While [Fact A] is confirmed, [Fact B] lacks verification."  
% - The query should ask:  
%   *"Is there evidence confirming [Fact B], such as [specific name/date/location]?"*  

% ### Domain-Specific Notes:  
% - Verify historical timelines (e.g., Ali Qushji died in 1474, not the 1960s).  
% - Confirm authorships (e.g., *Hayy ibn Yaqdhan* is attributed to Ibn Tufail, not Ali Qushji).  
% - Check for geographic or biographical details (e.g., birthplaces, cities, or professional connections).  

% Ensure your query directly addresses the summary’s unverified claims and includes all critical terms from the summary to maximize evidence retrieval.
% \end{lstlisting}
% \end{instructbox}\begin{instructbox}[\incode{hoverBench qwen3-8b create_query_hop3.predict}]
% GEPA Prompt:
% \begin{lstlisting}[breakindent=0pt, breaklines=true, breakatwhitespace=true, frame=single, backgroundcolor=\color{blue!10}]
% Given the fields `claim`, `summary_1`, and `summary_2`, produce the field `query` that:  
% 1. Explicitly asks whether the claim is supported by the provided summaries.  
% 2. Includes **specific evidence** from the summaries (e.g., names, titles, dates, or factual details) to guide retrieval of relevant documents.  
% 3. Focuses on **key disputed points** in the claim (e.g., incorrect attributions, missing evidence, or conflicting statements) to ensure the query targets the most relevant information.  

% **Key considerations for the query:**  
% - If the summaries mention specific works (e.g., *Planes, Trains and Automobiles*), include the title.  
% - If the summaries reference named entities (e.g., John Candy, KLM, Richard Ford), include them.  
% - If the summaries clarify contradictions (e.g., "not X but Y"), structure the query to highlight this contrast.  
% - Avoid vague phrasing; instead of asking "Is the claim true?" use "Does the evidence in the summaries confirm [specific detail]?"  

% **Example:**  
% If the claim is about a film and the summaries mention *Planes, Trains and Automobiles* (1987), the query should include that title to retrieve relevant evidence.  
% ```
% \end{lstlisting}
% \end{instructbox}\begin{instructbox}[\incode{hoverBench qwen3-8b summarize1.predict}]
% GEPA Prompt:
% \begin{lstlisting}[breakindent=0pt, breaklines=true, breakatwhitespace=true, frame=single, backgroundcolor=\color{blue!10}]
% Given the fields `claim` and `passages`, produce the fields `summary` by:  
% 1. **Analyzing the claim** to identify key elements (e.g., names, locations, relationships, specific terms).  
% 2. **Cross-referencing the passages** to locate explicit or implicit evidence that supports, contradicts, or partially addresses the claim.  
% 3. **Explicitly connecting the claim to the evidence** in the passages, even if the evidence is not directly stated in the claim (e.g., linking "Roger Yuan" to "Shanghai Noon" in Example 1).  
% 4. **Highlighting domain-specific details** (e.g., "Sentinelese language is unclassifiable" in Example 3) that are critical for accuracy.  
% 5. **Stating whether the claim is supported, unsupported, or partially supported**, and explaining the reasoning with clear references to the passages.  
% 6. **Prioritizing key evidence** that would be relevant for further queries (e.g., "Marlborough" in Example 2) to ensure retrievability.  

% **Additional Guidance:**  
% - If the claim contains unverified details (e.g., "Roger Yuan" in Example 1), explicitly note this in the summary.  
% - For niche or ambiguous information (e.g., "uncontacted people" in Example 3), reference specific passages that clarify the context.  
% - Ensure summaries are concise but include all critical evidence needed to validate or refute the claim.
% \end{lstlisting}
% \end{instructbox}\begin{instructbox}[\incode{hoverBench qwen3-8b summarize2.predict}]
% GEPA Prompt:
% \begin{lstlisting}[breakindent=0pt, breaklines=true, breakatwhitespace=true, frame=single, backgroundcolor=\color{blue!10}]
% Given the fields `claim`, `context`, and `passages`, produce the field `summary` that:  
% 1. **Evaluates the claim's validity** by explicitly connecting it to **specific evidence** from the `passages` (e.g., citing exact names, dates, or terms that support/refute the claim).  
% 2. **Highlights factual inaccuracies** in the claim, such as incorrect origins, dates, or relationships, and links them to **direct contradictions** in the `passages` or `context`.  
% 3. **Mentions key terms or entities** critical to the claim (e.g., "Sly Stone," "Large Münsterländer," or "Hit Parader's 2006 ranking") to ensure summaries are actionable for query generation.  
% 4. **Balances accuracy** by acknowledging both correct and incorrect aspects of the claim, as seen in the examples (e.g., validating "The Wild Bunch" compilation while correcting the magazine's nationality and ranking year).  
% 5. **Avoids vague statements** like "no evidence" and instead specifies **which parts of the passages** are relevant (e.g., "Passage 6 confirms Hit Parader was American").  

% **Generalizable Strategy**:  
% - Always reference **specific passages** (e.g., "Passage 1 states...") to ground claims in the text.  
% - Use **precise terminology** (e.g., "German Shorthaired Pointer" instead of "cousins") to align with the data.  
% - Prioritize **clarity and specificity** to ensure summaries can guide further evidence retrieval.
% \end{lstlisting}
% \end{instructbox}